\chapter{Barisan dan Deret} \label{seq:chapter}

%%%%%%%%%%%%%%%%%%%%%%%%%%%%%%%%%%%%%%%%%%%%%%%%%%%%%%%%%%%%%%%%%%%%%%%%%%%%%%

\section{Barisan dan limit}
\label{sec:seqsandlims}

%mbxINTROSUBSECTION

\sectionnotes{2,5 pertemuan kuliah}

Analisis pada dasarnya berkenaan dengan pengambilan limit.  Jenis limit yang paling dasar
adalah limit barisan bilangan real.
Kita telah melihat penggunaan barisan secara informal.  Mari kita berikan
definisi formalnya.

\begin{defn}
Sebuah \emph{\myindex{barisan}} (bilangan real) adalah fungsi $x \colon \N \to
\R$.  Alih-alih $x(n)$, kita
biasanya menyatakan suku ke-$n$ dalam barisan tersebut dengan $x_n$.
Untuk menyatakan sebuah barisan, kita tulis%
\footnote{Untuk menyingkat, $\{ x_n \}$ atau $\{ x_n \}_n$ lazim digunakan.}%
\glsadd{not:sequence}%
\begin{equation*}
\{ x_n \}_{n=1}^\infty.
\end{equation*}

Sebuah barisan $\{ x_n \}_{n=1}^\infty$ disebut \emph{terbatas}\index{barisan terbatas} jika
fungsi yang mendasarinya terbatas.  Artinya, jika
terdapat suatu $B \in \R$ sedemikian sehingga
\begin{equation*}
\sabs{x_n} \leq B \qquad \text{untuk semua } n \in \N.
\end{equation*}
Dengan kata lain, barisan $\{x_n\}_{n=1}^\infty$ terbatas apabila
himpunan $\{ x_n : n \in \N \}$
terbatas.
Dengan cara serupa, kita mendefinisikan istilah
\emph{terbatas di bawah}\index{terbatas di bawah!barisan} dan
\emph{terbatas di atas}\index{terbatas di atas!barisan}.
\end{defn}

Ketika kita perlu
memberikan sebuah barisan konkret, kita sering menyatakan setiap suku dengan rumus dalam
$n$.
Sebagai contoh, $\{ \nicefrac{1}{n} \}_{n=1}^\infty$ menyatakan
barisan $1, \nicefrac{1}{2}, \nicefrac{1}{3}, \nicefrac{1}{4},
\nicefrac{1}{5}, \ldots$.
Barisan $\{ \nicefrac{1}{n} \}_{n=1}^\infty$
merupakan barisan terbatas ($B=1$ sudah memadai).  Di sisi lain, barisan
$\{ n \}_{n=1}^\infty$ menyatakan
$1,2,3,4,\ldots$, dan barisan ini tidak terbatas (mengapa?)

Meskipun notasi untuk barisan
mirip\footnote{\cite{BS} menggunakan $(x_n)_{n=1}^\infty$ untuk menyatakan
sebuah barisan alih-alih $\{ x_n \}_{n=1}^\infty$, yaitu notasi yang digunakan oleh \cite{Rudin:baby}.
Kedua notasi ini lazim digunakan.}
dengan notasi untuk himpunan, kedua pengertian tersebut
berbeda.  Sebagai contoh, barisan $\bigl\{ {(-1)}^n \bigr\}_{n=1}^\infty$ adalah barisan
$-1,1,-1,1,-1,1,\ldots$, sedangkan himpunan nilainya, yaitu
\emph{daerah hasil barisan}\index{daerah hasil barisan},
hanyalah himpunan $\{ -1, 1 \}$.  Kita menuliskan himpunan ini
sebagai $\bigl\{ {(-1)}^n : n \in \N \bigr\}$.
% When ambiguity can arise, we
%use the words \emph{sequence} or \emph{set} to distinguish the two
%concepts.

Contoh barisan lainnya adalah \emph{\myindex{barisan konstan}}.
Barisan itu adalah $\{ c \}_{n=1}^\infty = c,c,c,c,\ldots$ yang terdiri atas satu
konstanta $c \in \R$ yang berulang tanpa henti.

%We now get to the idea of a
%\emph{limit of a sequence}\index{limit!of a sequence}.
%We will see in \propref{prop:limisunique}
%that a limit, if it exists, is unique.
%It makes sense to talk about \emph{the} limit of a sequence.

\begin{defn}
Sebuah barisan $\{ x_n \}_{n=1}^\infty$ dikatakan \emph{konvergen} ke bilangan
$x \in \R$ jika untuk setiap $\epsilon > 0$, terdapat suatu $M \in \N$ sedemikian
sehingga $\sabs{x_n - x} < \epsilon$ untuk semua $n \geq M$.
Bilangan $x$ disebut \emph{limit} barisan tersebut.
Kita akan membuktikan bahwa limit $x$ unik, sehingga kita tulis%
\footnote{Dalam teks, ini dapat ditampilkan sebagai $\lim_{n\to\infty} x_n$.}
\glsadd{not:limseq}%
\begin{equation*}
\lim_{n\to \infty} x_n \coloneqq x .
\end{equation*}

Barisan
yang konvergen disebut \emph{konvergen}\index{konvergen!barisan}.
Jika tidak, kita katakan barisan tersebut \emph{divergen}
atau bahwa barisan itu bersifat
\emph{divergen}\index{divergen!barisan}.
\end{defn}

Sebentar lagi, dalam \propref{prop:limisunique} kita akan menunjukkan bahwa limit $x$ selalu
unik jika ada.  Jadi, masuk akal untuk berbicara tentang \emph{limit} suatu
barisan dan kita hanya perlu menunjukkan bahwa barisan tersebut konvergen ke satu
bilangan.
Untuk beberapa contoh berikutnya, mari kita berpura-pura bahwa
kita telah membuktikan bahwa limit itu unik.

Secara intuitif, limit yang bernilai $x$ berarti bahwa pada akhirnya
setiap suku dalam barisan dekat dengan bilangan $x$.  Lebih tepatnya,
kita dapat berada sedekat apa pun dengan limit, asalkan kita melangkah cukup jauh dalam
barisan.  Ini tidak berarti bahwa kita pernah mencapai limit tersebut.  Mungkin saja,
dan bahkan cukup umum, tidak ada $x_n$ dalam barisan yang sama dengan
limit $x$.
Kita mengilustrasikan konsep ini dalam \figureref{figsequenceconvergence}.  Dalam
gambar tersebut, pertama-tama kita memandang barisan sebagai grafik karena barisan merupakan fungsi pada
$\N$.   Kedua, kita juga menggambarkannya sebagai barisan titik berlabel pada garis
bilangan real.

\begin{myfigureht}
\myincludepdft{sequence-convergence_full}{%
Dua diagram.
Diagram pertama adalah grafik sebuah barisan, yaitu,
nilai-nilainya berada pada sumbu vertikal dan sumbu horizontal menyatakan indeksnya.
Pada sumbu vertikal,
bilangan x ditandai dan sebuah garis horizontal
digambar.  Sebuah pita dari
x dikurangi epsilon hingga x ditambah epsilon diarsir.  Nilai-nilai
barisan sebagian besar berada di luar pita yang diarsir
hingga indeksnya mencapai 7, yang ditandai sebagai M.  Semua nilai
berikutnya berada di dalam pita tersebut.
Pada diagram kedua, sumbu x mendatar, dengan
bilangan x ditandai dan nilai-nilai x subskrip 1, x subskrip 2, hingga x subskrip 10 ditandai
sebagai tanda skala.}
\caption{Ilustrasi konvergensi.
Di bagian atas, kita tampilkan sepuluh titik pertama barisan sebagai grafik
dengan $M$ dan interval di sekitar limit $x$ ditandai.
Di bagian bawah, titik-titik dari barisan yang sama ditandai pada
garis bilangan.\label{figsequenceconvergence}}
\end{myfigureht}

Ketika kita menulis $\lim_{n\to\infty} x_n = x$ untuk suatu bilangan real $x$, kita menyatakan dua
hal: pertama, bahwa $\{ x_n \}_{n=1}^\infty$ konvergen, dan kedua, bahwa limitnya adalah
$x$.

Definisi di atas merupakan salah satu definisi terpenting dalam analisis,
dan definisi tersebut harus dipahami dengan sempurna.  Pokok utama dalam
definisi tersebut adalah bahwa untuk \emph{sebarang} $\epsilon > 0$, kita dapat menemukan suatu $M$.  Nilai
$M$ dapat bergantung pada $\epsilon$, jadi kita baru memilih suatu $M$ setelah mengetahui
$\epsilon$.  Mari kita ilustrasikan konvergensi melalui beberapa contoh.

\begin{example}
Barisan konstan $1,1,1,1,\ldots$ konvergen ke 1.  Untuk
setiap $\epsilon > 0$, pilih $M = 1$.  Artinya, $\sabs{x_n-x} = \sabs{1-1} < \epsilon$ untuk semua $n$.
\end{example}

\begin{example}
Klaim: \emph{Barisan $\{ \nicefrac{1}{n} \}_{n=1}^\infty$ konvergen dan}
\begin{equation*}
\lim_{n\to \infty} \frac{1}{n} = 0 .
\end{equation*}
Bukti: Diberikan sebarang $\epsilon > 0$, cari suatu $M \in \N$ sedemikian sehingga
$0 < \nicefrac{1}{M} < \epsilon$
(di sini \hyperref[thm:arch:i]{sifat Archimedes} digunakan).
Untuk semua $n \geq M$,
\begin{equation*}
\sabs{x_n - x}
=
\abs{\frac{1}{n} - 0} = \abs{\frac{1}{n}} = \frac{1}{n} \leq \frac{1}{M} < \epsilon .
\end{equation*}
\end{example}

\begin{example}
Barisan $\bigl\{ {(-1)}^n \bigr\}_{n=1}^\infty$ divergen.  Bukti: Jika terdapat
limit $x$, maka untuk $\epsilon = \frac{1}{2}$, definisi mengharuskan adanya suatu $M$ yang
memenuhi syarat tersebut.  Andaikan
$M$ seperti itu ada.
Maka, untuk $n \geq M$ yang genap, kita hitung
\begin{equation*}
\nicefrac{1}{2} > \sabs{x_n - x}  = \sabs{1 - x}
\qquad \text{dan} \qquad
\nicefrac{1}{2} > \sabs{x_{n+1} - x}  = \sabs{-1 - x} .
\end{equation*}
Dengan demikian, kita memperoleh kontradiksi
\begin{equation*}
2 = \babs{1 - x - (-1 -x)} \leq
\sabs{1 - x} + \sabs{-1 -x} < \nicefrac{1}{2} + \nicefrac{1}{2} = 1 .
\end{equation*}
\end{example}

\begin{prop} \label{prop:limisunique}
Barisan konvergen memiliki limit yang unik.
\end{prop}

Pembuktian proposisi ini memperlihatkan suatu teknik yang berguna dalam
analisis.  Banyak pembuktian mengikuti skema umum yang sama.  Kita ingin
menunjukkan bahwa suatu kuantitas bernilai nol.  Kita menulis kuantitas itu sebagai
dua kuantitas menggunakan ketaksamaan segitiga, lalu menaksir masing-masing
dengan bilangan yang sekecil apa pun.

\begin{proof}
%NOTE: should be word for word the same as 7.3.3
Andaikan $\{ x_n \}_{n=1}^\infty$ memiliki limit $x$ dan $y$.
Ambil sebarang $\epsilon > 0$.
Berdasarkan definisi, cari suatu $M_1$ sedemikian sehingga untuk semua $n \geq M_1$,
$\sabs{x_n-x} < \nicefrac{\epsilon}{2}$.  Demikian pula, cari suatu $M_2$
sedemikian sehingga untuk semua $n \geq M_2$, berlaku
$\sabs{x_n-y} < \nicefrac{\epsilon}{2}$.
Sekarang ambil suatu $n$ sedemikian sehingga $n \geq M_1$ dan juga $n \geq M_2$, lalu taksir
\begin{equation*}
\begin{split}
\sabs{y-x}
& =
\babs{x_n-x - (x_n -y)} \\
& \leq
\sabs{x_n-x} + \sabs{x_n -y} \\
& <
\frac{\epsilon}{2} + \frac{\epsilon}{2} = \epsilon .
\end{split}
\end{equation*}
Karena $\sabs{y-x} < \epsilon$ untuk setiap $\epsilon > 0$, diperoleh $\sabs{y-x} = 0$
dan $y=x$.  Jadi, limit tersebut (jika ada) unik.
\end{proof}

\begin{prop}
Setiap barisan konvergen $\{ x_n \}_{n=1}^\infty$ merupakan barisan terbatas.
\end{prop}

\begin{proof}
Andaikan $\{ x_n \}_{n=1}^\infty$ konvergen ke $x$.  Jadi, terdapat suatu $M \in \N$
yang dapat diasumsikan setidaknya 2, dan untuk semua $n \geq M$, berlaku
$\sabs{x_n - x} < 1$.  Untuk $n \geq M$,
\begin{equation*}
\begin{split}
\sabs{x_n} & = \sabs{x_n - x + x}
\\
& \leq \sabs{x_n - x} + \sabs{x}
\\
& < 1 + \sabs{x} .
\end{split}
\end{equation*}
Himpunan $\bigl\{ \sabs{x_1}, \sabs{x_2}, \ldots, \sabs{x_{M-1}} , 1+\sabs{x} \bigr\}$
merupakan himpunan berhingga; oleh karena itu, tetapkan
\begin{equation*}
B \coloneqq \max \bigl\{ \sabs{x_1}, \sabs{x_2}, \ldots, \sabs{x_{M-1}}, 1+\sabs{x} \bigr\} .
\end{equation*}
Maka, untuk semua $n \in \N$,
\begin{equation*}
\sabs{x_n} \leq B. \qedhere
\end{equation*}
\end{proof}

Barisan $\bigl\{ {(-1)}^n \bigr\}_{n=1}^\infty$ menunjukkan bahwa pernyataan kebalikannya
tidak berlaku.  Barisan terbatas belum tentu konvergen.

\begin{example}
Mari kita tunjukkan bahwa $\left\{ \frac{n^2+1}{n^2+n} \right\}_{n=1}^\infty$ konvergen dan
\begin{equation*}
\lim_{n\to\infty} \frac{n^2+1}{n^2+n} = 1 .
\end{equation*}

Diberikan sebarang $\epsilon > 0$,
cari $M \in \N$ sedemikian sehingga $\frac{1}{M} < \epsilon$.  Maka, untuk semua $n \geq
M$,
\begin{equation*}
\begin{split}
%\abs{\frac{n^2+1}{n^2+n} - 1} & =
%\abs{\frac{n^2+1 - (n^2+n)}{n^2+n}} \\
\abs{\frac{n^2+1}{n^2+n} - 1}  =
\abs{\frac{n^2+1 - (n^2+n)}{n^2+n}}
& =
\abs{\frac{1 - n}{n^2+n}} \\
& =
\frac{n-1}{n^2+n} \\
& \leq 
\frac{n}{n^2+n} 
 =
\frac{1}{n+1}  \\
& \leq \frac{1}{n}
\leq \frac{1}{M} < \epsilon .
\end{split}
\end{equation*}
Oleh karena itu,
$\lim_{n\to\infty} \frac{n^2+1}{n^2+n} = 1$.
Contoh ini menunjukkan bahwa kadang-kadang, untuk memperoleh yang diinginkan, kita harus membuang
sebagian informasi guna mendapatkan taksiran yang lebih sederhana.
\end{example}

\subsection{Barisan monoton}

Jenis barisan yang paling sederhana adalah barisan monoton.  Memeriksa bahwa
suatu barisan monoton konvergen semudah memeriksa bahwa barisan itu terbatas.
Menentukan
limit suatu barisan monoton
yang konvergen juga mudah, asalkan kita dapat menentukan supremum atau infimum
dari suatu himpunan bilangan yang terhitung.

\begin{defn}
Suatu barisan $\{ x_n \}_{n=1}^\infty$ disebut \emph{monoton naik}\index{barisan
monoton naik} jika $x_n \leq x_{n+1}$ untuk setiap $n \in \N$.  
%
Suatu barisan $\{ x_n \}_{n=1}^\infty$ disebut \emph{monoton turun}\index{barisan
monoton turun} jika $x_n \geq x_{n+1}$ untuk setiap $n \in \N$.  
%
Jika suatu barisan monoton naik atau monoton turun, kita
cukup mengatakan bahwa barisan tersebut \emph{monoton}\index{barisan
monoton}.\footnote{Sebagian
penulis memakai istilah \emph{monotonik}\index{barisan monotonik}.}
\end{defn}

Sebagai contoh,
$\{ n \}_{n=1}^\infty$ monoton naik,
$\{ \nicefrac{1}{n} \}_{n=1}^\infty$ monoton turun,
barisan konstan $\{ 1 \}_{n=1}^\infty$ sekaligus monoton naik dan monoton
turun, sedangkan $\bigl\{ {(-1)}^n \bigr\}_{n=1}^\infty$ tidak monoton.
Beberapa suku pertama dari suatu contoh barisan monoton naik
ditampilkan dalam 
\figureref{figsequenceincreasing}.

\begin{myfigureht}
\myincludegraphics{sequence-increasing}{%
Grafik beberapa suku pertama suatu barisan dengan sumbu horizontal sebagai
indeks.  Setiap nilai berikutnya dalam barisan sama dengan atau lebih besar
daripada nilai sebelumnya, sehingga titik-titik yang menyatakan nilainya bergerak naik saat kita bergerak
ke kanan sepanjang grafik.}
\caption{Beberapa suku pertama barisan monoton naik sebagai sebuah
grafik.\label{figsequenceincreasing}}
\end{myfigureht}

\begin{thm}[Teorema konvergensi monoton] \label{thm:monotoneconv}
\index{teorema konvergensi monoton}%
Suatu barisan monoton $\{ x_n \}_{n=1}^\infty$ terbatas jika dan hanya jika konvergen.

Lebih lanjut, jika $\{ x_n \}_{n=1}^\infty$ monoton naik dan terbatas, maka
\begin{equation*}
\lim_{n\to \infty} x_n = \sup \{ x_n : n \in \N \} .
\end{equation*}
Jika $\{ x_n \}_{n=1}^\infty$ monoton turun dan terbatas, maka
\begin{equation*}
\lim_{n\to \infty} x_n = \inf \{ x_n : n \in \N \} .
\end{equation*}
\end{thm}

\begin{proof}
Perhatikan suatu barisan monoton naik $\{ x_n \}_{n=1}^\infty$.  Andaikan 
terlebih dahulu barisan tersebut terbatas,
yakni,
himpunan $\{ x_n : n \in  \N \}$ terbatas.  Misalkan
\begin{equation*}
x \coloneqq \sup \{ x_n : n \in \N \} .
\end{equation*}
Ambil sembarang $\epsilon > 0$.  Karena $x$ merupakan supremum,
harus terdapat setidaknya satu $M \in \N$ sedemikian sehingga $x_{M} > x-\epsilon$.
Karena $\{ x_n \}_{n=1}^\infty$ monoton naik,
mudah dilihat (dengan \hyperref[induction:thm]{induksi}) bahwa
$x_n \geq x_{M}$ untuk setiap $n \geq M$.  Jadi, untuk setiap $n \geq M$,
\begin{equation*}
\sabs{x_n-x} = x-x_n \leq x-x_{M} < \epsilon  .
\end{equation*}
Dengan demikian, $\{ x_n \}_{n=1}^\infty$ konvergen ke $x$.
Oleh karena itu, suatu barisan monoton naik yang terbatas bersifat konvergen.
Untuk arah sebaliknya, kita telah membuktikan bahwa suatu barisan konvergen
bersifat terbatas.

Bukti untuk barisan monoton turun diserahkan sebagai latihan.
\end{proof}

Suatu barisan monoton naik $\{ x_n \}_{n=1}^\infty$
selalu terbatas di bawah karena $x_1 \leq x_2 \leq \cdots \leq x_n$
untuk setiap $n$, sehingga $x_1$ merupakan batas bawah.  Jadi, untuk menentukan apakah suatu 
barisan monoton naik terbatas, cukup diperiksa apakah barisan tersebut
terbatas di atas.  Demikian pula, suatu barisan monoton turun
selalu terbatas di atas, sehingga cukup diperiksa apakah barisan tersebut
terbatas di bawah.

\begin{example}
Ambil barisan $\bigl\{ \frac{1}{\sqrt{n}} \bigr\}_{n=1}^\infty$.

Barisan tersebut
terbatas di bawah karena
$\frac{1}{\sqrt{n}} > 0$ untuk setiap $n \in \N$.
Mari kita tunjukkan bahwa barisan ini monoton turun.  Kita mulai dari
$\sqrt{n+1} \geq \sqrt{n}$ (mengapa ini benar?).  Dari pertidaksamaan ini
kita peroleh
\begin{equation*}
\frac{1}{\sqrt{n+1}} \leq \frac{1}{\sqrt{n}} .
\end{equation*}
Jadi, barisan tersebut monoton turun dan terbatas di bawah (karena itu
terbatas).  Berdasarkan \thmref{thm:monotoneconv}, kita mendapati bahwa barisan tersebut
konvergen dan
\begin{equation*}
\lim_{n\to \infty} \frac{1}{\sqrt{n}}
=
\inf \left\{ \frac{1}{\sqrt{n}} : n \in \N \right\} .
\end{equation*}
Kita telah mengetahui bahwa infimum tersebut lebih besar daripada atau sama dengan 0, karena
0 merupakan batas bawah.  Ambil bilangan $b \geq 0$ sedemikian
sehingga $b \leq \frac{1}{\sqrt{n}}$ untuk setiap $n$.  Kita kuadratkan kedua ruas untuk
memperoleh
\begin{equation*}
b^2 \leq \frac{1}{n} \qquad \text{untuk setiap } n \in \N.
\end{equation*}
Kita telah melihat sebelumnya bahwa ini menyiratkan $b^2 \leq 0$ (suatu konsekuensi
dari \hyperref[thm:arch:i]{sifat Archimedes}).  Karena $b^2 \geq 0$
juga, kita peroleh $b^2 = 0$
dan dengan demikian $b = 0$.
Oleh karena itu, $b=0$ merupakan batas bawah terbesar, dan $\lim\limits_{n\to\infty} \frac{1}{\sqrt{n}} = 0$.
\end{example}

\begin{example}
Perlu berhati-hati:  Menunjukkan bahwa suatu barisan monoton terbatas
agar dapat menggunakan \thmref{thm:monotoneconv} mungkin sulit.
Barisan
$\{ 1 + \nicefrac{1}{2} + \cdots + \nicefrac{1}{n} \}_{n=1}^\infty$
merupakan barisan monoton naik
yang bertambah dengan lambat dan bahkan bertambah semakin lambat ketika $n$ semakin
besar.  Nanti, ketika sampai pada pembahasan deret, kita akan melihat
bahwa barisan ini tidak memiliki batas atas sehingga tidak konvergen.  Sama sekali tidak
jelas bahwa barisan ini tidak memiliki batas atas.
\end{example}

Contoh umum munculnya barisan monoton diberikan oleh
proposisi berikut.  Buktinya diserahkan sebagai latihan.

\begin{prop} \label{prop:supinfseq}
Misalkan $S \subset \R$ merupakan himpunan terbatas yang tak kosong.
Maka terdapat barisan-barisan monoton
$\{ x_n \}_{n=1}^\infty$ dan $\{ y_n \}_{n=1}^\infty$ sedemikian sehingga $x_n, y_n \in S$ dan
\begin{equation*}
\sup\,S = \lim_{n\to \infty} x_n \qquad \text{dan} \qquad \inf\,S =
\lim_{n\to\infty} y_n .
\end{equation*}
\end{prop}

\subsection{Ekor suatu barisan}

\begin{defn}
Untuk suatu barisan $\{ x_n \}_{n=1}^\infty$,
\emph{ekor-$K$} (dengan $K \in \N$),
atau cukup
\emph{ekor},\index{ekor suatu barisan} dari
$\{ x_n \}_{n=1}^\infty$ adalah barisan yang dimulai dari suku berindeks $K+1$, biasanya ditulis sebagai
\begin{equation*}
\{ x_{n+K} \}_{n=1}^\infty
\qquad \text{atau} \qquad \{ x_n \}_{n=K+1}^\infty .
\end{equation*}
\end{defn}

Sebagai contoh, ekor-$4$ dari $\{ \nicefrac{1}{n} \}_{n=1}^\infty$ adalah
$\nicefrac{1}{5}, \nicefrac{1}{6}, \nicefrac{1}{7}, \nicefrac{1}{8},
\ldots$.  Ekor-$0$ suatu barisan adalah barisan itu sendiri.
Kekonvergenan dan limit suatu barisan hanya bergantung pada ekornya.

\begin{prop}
Misalkan $\{ x_n \}_{n=1}^\infty$ suatu barisan.  Maka pernyataan-pernyataan
berikut ekuivalen:
\begin{enumerate}[(i)]
\item \label{prop:ktail:i}
Barisan $\{ x_n \}_{n=1}^\infty$ konvergen.
\item \label{prop:ktail:ii}
Ekor-$K$ $\{ x_{n+K} \}_{n=1}^\infty$ konvergen untuk semua $K \in \N$.
\item \label{prop:ktail:iii}
Ekor-$K$ $\{ x_{n+K} \}_{n=1}^\infty$ konvergen untuk suatu $K \in \N$.
\end{enumerate}
Lebih lanjut, jika salah satu (dan karenanya semua) limit tersebut ada, maka untuk semua $K \in \N$
\begin{equation*}
\lim_{n\to \infty} x_n = \lim_{n \to \infty} x_{n+K} .
\end{equation*}
\end{prop}

\begin{proof}
Jelas bahwa
\ref{prop:ktail:ii} mengakibatkan \ref{prop:ktail:iii}.
Oleh karena itu, pertama-tama kita akan menunjukkan bahwa
\ref{prop:ktail:i}
mengakibatkan
\ref{prop:ktail:ii},
kemudian menunjukkan bahwa
\ref{prop:ktail:iii}
mengakibatkan
\ref{prop:ktail:i}.  Dengan kata lain, 
%mbxSTARTIGNORE
\begin{equation*}
\begin{tikzcd}[row sep=0pt]
& \text{\ref{prop:ktail:ii}} \ar[dr, Rightarrow] & \\
\text{\ref{prop:ktail:i}} \ar[ur, Rightarrow, "\text{untuk dibuktikan}"] & &
\text{\ref{prop:ktail:iii}} \ar[ll, Rightarrow, "\text{untuk dibuktikan}"] 
\end{tikzcd}
\end{equation*}
%mbxENDIGNORE
%mbxlatex \begin{center}
%mbxlatex \myincludepdft{tail_of_sequence_cd}{Diagram dengan
%mbxlatex tiga panah implikasi.  Panah pertama dari (i) ke (ii)
%mbxlatex berlabel "untuk dibuktikan".  Panah kedua dari (ii) ke (iii)
%mbxlatex tidak berlabel.  Panah ketiga dari (iii) ke (i)
%mbxlatex berlabel "untuk dibuktikan".}
%mbxlatex \end{center}
Dalam prosesnya, kita juga akan menunjukkan bahwa limit-limit tersebut sama.

Kita mulai dengan membuktikan bahwa \ref{prop:ktail:i} mengakibatkan \ref{prop:ktail:ii}.
Andaikan $\{x_n \}_{n=1}^\infty$ konvergen ke suatu $x \in \R$.
Ambil sebarang $K \in \N$, dan
definisikan $y_n \coloneqq x_{n+K}$.  Kita hendak menunjukkan bahwa $\{ y_n \}_{n=1}^\infty$ konvergen
ke $x$.
Untuk suatu $\epsilon > 0$, terdapat $M \in \N$ sedemikian sehingga
$\sabs{x-x_n} < \epsilon$ untuk semua $n \geq M$.
Perhatikan bahwa $n \geq M$ mengakibatkan $n+K \geq M$.  Oleh karena itu, untuk
semua $n \geq M$, kita memiliki
\begin{equation*}
\sabs{x-y_n} = \sabs{x-x_{n+K}} < \epsilon .
\end{equation*}
Akibatnya, $\{ y_n \}_{n=1}^\infty$ konvergen ke $x$.

Selanjutnya, kita buktikan bahwa \ref{prop:ktail:iii} mengakibatkan \ref{prop:ktail:i}.
Ambil suatu $K \in \N$, definisikan
$y_n \coloneqq x_{n+K}$, dan andaikan $\{ y_n \}_{n=1}^\infty$ konvergen ke $x \in \R$.
Artinya, untuk suatu $\epsilon > 0$, terdapat $M' \in \N$ sedemikian sehingga
$\sabs{x-y_n} < \epsilon$ untuk semua $n \geq M'$.
Tetapkan $M \coloneqq M'+K$.  Maka $n \geq M$ mengakibatkan $n-K \geq M'$.
Jadi, setiap kali $n \geq M$, kita memiliki
\begin{equation*}
\sabs{x-x_n} = \sabs{x-y_{n-K}} < \epsilon.
\end{equation*}
Oleh karena itu, $\{ x_n \}_{n=1}^\infty$ konvergen ke $x$.
\end{proof}

Dengan demikian, limit tidak bergantung pada bagaimana suatu barisan dimulai, melainkan hanya
pada ekor barisan tersebut.  Bagian awal barisan
boleh sebarang.

Sebagai contoh, barisan yang didefinisikan oleh $x_n \coloneqq \frac{n}{n^2+16}$ menurun
jika kita mulai dari $n=4$ (sebelumnya barisan tersebut meningkat).  Yaitu:
\begin{equation*}
\begin{aligned}
\{ x_n \}_{n=1}^\infty
={}& \nicefrac{1}{17}, \nicefrac{1}{10}, \nicefrac{3}{25}, \nicefrac{1}{8}, \nicefrac{5}{41},
\\
&
\nicefrac{3}{26}, \nicefrac{7}{65}, \nicefrac{1}{10},
\nicefrac{9}{97}, \nicefrac{5}{58},\ldots .
\end{aligned}
\end{equation*}
Selain itu,
\begin{equation*}
\nicefrac{1}{17} <
\nicefrac{1}{10} <
\nicefrac{3}{25} <
\nicefrac{1}{8} >
\nicefrac{5}{41} >
\nicefrac{3}{26} >
\nicefrac{7}{65} >
\nicefrac{1}{10} >
\nicefrac{9}{97} >
\nicefrac{5}{58} > \ldots .
\end{equation*}
Jika kita membuang 3 suku pertama
dan meninjau ekor-3, barisan tersebut menurun.  Buktinya diserahkan sebagai latihan.  Karena ekor-3
monoton dan terbatas di bawah oleh nol, ekor tersebut konvergen; oleh karena itu, barisan semula konvergen.

\subsection{Subbarisan}

Kadang-kadang berguna untuk meninjau hanya sebagian suku suatu barisan.
Subbarisan dari $\{ x_n \}_{n=1}^\infty$ adalah barisan yang memuat
hanya sebagian bilangan dari $\{ x_n \}_{n=1}^\infty$ dengan urutan yang sama.

\begin{defn}
Misalkan $\{ x_n \}_{n=1}^\infty$ suatu barisan.
Misalkan $\{ n_i \}_{i=1}^\infty$ suatu barisan bilangan asli yang naik
secara ketat, yaitu, $n_i < n_{i+1}$ untuk setiap $i \in \N$ (dengan kata lain $n_1 < n_2 < n_3 < \cdots$).  
Barisan
\begin{equation*}
\{ x_{n_i} \}_{i=1}^\infty
\end{equation*}
disebut \glsadd{not:subsequence}
sebuah \emph{\myindex{subbarisan}} dari $\{ x_n \}_{n=1}^\infty$.
\end{defn}

Jadi, subbarisan tersebut adalah barisan $x_{n_1},x_{n_2},x_{n_3},\ldots$.
Sebagai contoh,
barisan
$\{ \nicefrac{1}{3i} \}_{i=1}^\infty
= \nicefrac{1}{3},\nicefrac{1}{6},\nicefrac{1}{9},\nicefrac{1}{12},\ldots$
merupakan subbarisan dari
$\{ \nicefrac{1}{n} \}_{n=1}^\infty
= 1, \nicefrac{1}{2},\nicefrac{1}{3},\nicefrac{1}{4},\ldots$.
Untuk melihat kesesuaian kedua
barisan ini dengan definisi, ambil $n_i \coloneqq 3i$,
yakni, $\{ n_i \}_{i=1}^\infty$ adalah
barisan $3,6,9,12,\ldots$.  
Bilangan-bilangan $x_{n_i}$ dalam
subbarisan harus berasal dari barisan semula.  Jadi, $1,0,\nicefrac{1}{3},0,
\nicefrac{1}{5},\ldots$
bukan subbarisan dari $\{ \nicefrac{1}{n} \}_{n=1}^\infty$.  Demikian pula, urutan
harus dipertahankan;
barisan $1,\nicefrac{1}{3},\nicefrac{1}{2},\nicefrac{1}{5},\ldots$
bukan subbarisan dari $\{ \nicefrac{1}{n} \}_{n=1}^\infty$.

Ekor suatu barisan merupakan salah satu jenis khusus subbarisan.  Untuk sembarang
subbarisan, kita mempunyai proposisi berikut mengenai kekonvergenan,
yang lebih lemah daripada hasil untuk ekor.

\begin{prop} \label{prop:seqtosubseq}
Jika $\{ x_n \}_{n=1}^\infty$ merupakan barisan konvergen,
maka setiap subbarisan $\{ x_{n_i} \}_{i=1}^\infty$ juga konvergen, dan
\begin{equation*}
\lim_{n\to \infty} x_n = 
\lim_{i\to \infty} x_{n_i} .
\end{equation*}
\end{prop}

\begin{proof}
Andaikan $\lim_{n\to \infty} x_n = x$.  Jadi, untuk setiap
$\epsilon > 0$, terdapat $M \in \N$ sedemikian sehingga untuk setiap $n \geq M$,
\begin{equation*}
\sabs{x_n - x} < \epsilon .
\end{equation*}
Tidak sulit membuktikan (buktikanlah!) dengan \hyperref[induction:thm]{induksi} bahwa
$n_i \geq i$ untuk setiap $i \in \N$.  Jadi, $i \geq M$ mengakibatkan $n_i \geq M$.  Dengan demikian,
untuk setiap $i \geq M$,
\begin{equation*}
\sabs{x_{n_i} - x} < \epsilon ,
\end{equation*}
dan bukti selesai.
\end{proof}

\begin{example}
Adanya suatu subbarisan konvergen tidak mengakibatkan
barisan itu sendiri konvergen.
Ambil barisan $0,1,0,1,0,1,\ldots$.  Artinya,
$x_n = 0$ jika $n$ ganjil, dan $x_n = 1$ jika $n$ genap.  Barisan
$\{ x_n \}_{n=1}^\infty$ divergen; namun, subbarisan
$\{ x_{2i} \}_{i=1}^\infty$ konvergen ke 1 dan subbarisan
$\{ x_{2i+1} \}_{i=1}^\infty$ konvergen ke 0.  Bandingkan dengan \propref{seqconvsubseqconv:prop}.
\end{example}

\subsection{Latihan}

\begin{exnote}
Dalam latihan-latihan berikut, silakan gunakan pengetahuan Anda dari kalkulus untuk
menentukan limitnya, jika limit tersebut ada.  Namun, Anda harus \emph{membuktikan}
bahwa limit yang
Anda peroleh benar, atau membuktikan bahwa barisan tersebut divergen.
\end{exnote}

\begin{exercise}
Apakah barisan
$\{ 3n \}_{n=1}^\infty$
terbatas?  Buktikan atau sangkal.
\end{exercise}

\begin{exercise}
Apakah barisan
$\{ n \}_{n=1}^\infty$
konvergen?  Jika ya, berapakah limitnya?
\end{exercise}

\begin{exercise}
Apakah barisan
$\left\{ \dfrac{{(-1)}^n}{2n} \right\}_{n=1}^\infty$
konvergen?  Jika ya, berapakah limitnya?
\end{exercise}

\begin{exercise}
Apakah barisan
$\{ 2^{-n} \}_{n=1}^\infty$
konvergen?  Jika ya, berapakah limitnya?
\end{exercise}

\begin{exercise}
Apakah barisan
$\left\{ \dfrac{n}{n+1} \right\}_{n=1}^\infty$
konvergen?  Jika ya, berapakah limitnya?
\end{exercise}

\begin{exercise}
Apakah barisan
$\left\{ \dfrac{n}{n^2+1} \right\}_{n=1}^\infty$
konvergen?  Jika ya, berapakah limitnya?
\end{exercise}

\begin{exercise} \label{exercise:absconv}
Misalkan $\{ x_n \}_{n=1}^\infty$ suatu barisan.
\begin{enumerate}[a)]
\item Tunjukkan bahwa $\lim\limits_{n\to\infty} x_n = 0$ (artinya, limit tersebut ada dan bernilai nol)
jika dan hanya jika $\lim\limits_{n\to\infty} \sabs{x_n} = 0$.
\item Temukan sebuah contoh sedemikian sehingga $\{ \sabs{x_n} \}_{n=1}^\infty$ konvergen dan
$\{ x_n \}_{n=1}^\infty$
divergen.
\end{enumerate}
\end{exercise}

\begin{exercise}
Apakah barisan
$\left\{ \dfrac{2^n}{n!} \right\}_{n=1}^\infty$
konvergen?  Jika ya, berapakah limitnya?
\end{exercise}

\begin{exercise}
Tunjukkan bahwa barisan
$\left\{ \dfrac{1}{\sqrt[3]{n}} \right\}_{n=1}^\infty$ monoton dan terbatas.  Kemudian gunakan
\thmref{thm:monotoneconv} untuk menentukan limitnya.
\end{exercise}

\begin{exercise}
Tunjukkan bahwa barisan
$\left\{ \dfrac{n+1}{n} \right\}_{n=1}^\infty$
monoton dan terbatas.  Kemudian gunakan
\thmref{thm:monotoneconv} untuk menentukan limitnya.
\end{exercise}

\begin{exercise}
Lengkapi bukti \thmref{thm:monotoneconv} untuk barisan
monoton turun.
\end{exercise}

\begin{exercise}
Buktikan \propref{prop:supinfseq}.
\end{exercise}

\begin{exercise}
Misalkan $\{ x_n \}_{n=1}^\infty$ suatu barisan monoton yang konvergen.  Andaikan
terdapat $k \in \N$ sedemikian sehingga
\begin{equation*}
\lim_{n\to \infty} x_n = x_k .
\end{equation*}
Tunjukkan bahwa $x_n = x_k$ untuk semua $n \geq k$.
\end{exercise}

\begin{exercise}
\pagebreak[2]
Temukan sebuah subbarisan konvergen dari barisan
$\bigl\{ {(-1)}^n \bigr\}_{n=1}^\infty$.
\end{exercise}

\begin{exercise}
Misalkan $\{x_n\}_{n=1}^\infty$ suatu barisan yang didefinisikan oleh
\begin{equation*}
x_n \coloneqq 
\begin{cases}
n               & \text{jika } n \text{ ganjil} , \\
\nicefrac{1}{n} & \text{jika } n \text{ genap} .
\end{cases}
\avoidbreak
\end{equation*}
\begin{enumerate}[a)]
\item Apakah barisan tersebut terbatas? (buktikan atau sangkal)
\item Apakah terdapat subbarisan konvergen?  Jika ya, temukan subbarisan tersebut.
\end{enumerate}
\end{exercise}

\begin{exercise}
\pagebreak[2]
Misalkan $\{ x_n \}_{n=1}^\infty$ suatu barisan.
Andaikan terdapat dua subbarisan konvergen $\{ x_{n_i} \}_{i=1}^\infty$ dan
$\{ x_{m_i} \}_{i=1}^\infty$.  Andaikan
\begin{equation*}
\lim_{i\to\infty} x_{n_i} = a
\qquad \text{dan} \qquad
\lim_{i\to\infty} x_{m_i} = b,
\end{equation*}
dengan $a \neq b$.  Buktikan bahwa $\{ x_n \}_{n=1}^\infty$ tidak konvergen tanpa
menggunakan \propref{prop:seqtosubseq}.
\end{exercise}

\begin{exercise}[Perlu kecermatan]
Temukan suatu barisan $\{ x_n \}_{n=1}^\infty$ sedemikian sehingga untuk setiap $y \in \R$, terdapat
subbarisan $\{ x_{n_i} \}_{i=1}^\infty$ yang konvergen ke $y$.
\end{exercise}

\begin{exercise}[Mudah]
Misalkan $\{ x_n \}_{n=1}^\infty$ suatu barisan dan $x \in \R$.
Andaikan untuk setiap $\epsilon > 0$, terdapat $M$ sedemikian sehingga
$\sabs{x_n-x} \leq \epsilon$ untuk semua $n \geq M$.
Tunjukkan bahwa $\lim\limits_{n\to\infty} x_n = x$.
\end{exercise}

\begin{exercise}[Mudah]
Misalkan $\{ x_n \}_{n=1}^\infty$ suatu barisan dan $x \in \R$ sedemikian sehingga
terdapat $k \in \N$ sehingga untuk semua $n \geq k$,
$x_n = x$.  Buktikan bahwa $\{ x_n \}_{n=1}^\infty$ konvergen ke $x$.
\end{exercise}

\begin{exercise}
Misalkan $\{ x_n \}_{n=1}^\infty$ suatu barisan dan
definisikan barisan $\{ y_n \}_{n=1}^\infty$ dengan
$y_{2k} \coloneqq x_{k^2}$ dan $y_{2k-1} \coloneqq x_k$ untuk semua $k \in \N$.
Buktikan bahwa $\{ x_n \}_{n=1}^\infty$ konvergen jika dan hanya jika
$\{ y_n \}_{n=1}^\infty$ konvergen.
Lebih lanjut, buktikan bahwa jika keduanya konvergen, maka
$\lim\limits_{n\to\infty} x_n = \lim\limits_{n\to\infty} y_n$.
\end{exercise}

\begin{exercise}
Tunjukkan bahwa ekor-3 dari barisan yang didefinisikan oleh $x_n \coloneqq \frac{n}{n^2+16}$
monoton turun.  Petunjuk: Andaikan $n \geq m \geq 4$ dan tinjau
pembilang dari bentuk $x_n-x_m$.
\end{exercise}

\begin{exercise}
Andaikan $\{ x_n \}_{n=1}^\infty$ suatu barisan sedemikian sehingga
subbarisan $\{ x_{2n} \}_{n=1}^\infty$, $\{ x_{2n-1} \}_{n=1}^\infty$, dan
$\{ x_{3n} \}_{n=1}^\infty$ semuanya konvergen.  Tunjukkan bahwa $\{ x_n \}_{n=1}^\infty$ konvergen.
\end{exercise}

\begin{exercise}
Andaikan $\{ x_n \}_{n=1}^\infty$ suatu barisan monoton naik yang
memiliki subbarisan konvergen.
Tunjukkan bahwa $\{ x_n \}_{n=1}^\infty$ konvergen.
Catatan: Jadi, untuk barisan monoton, \propref{prop:seqtosubseq} merupakan pernyataan \myquote{jika dan hanya jika}.
\end{exercise}

%%%%%%%%%%%%%%%%%%%%%%%%%%%%%%%%%%%%%%%%%%%%%%%%%%%%%%%%%%%%%%%%%%%%%%%%%%%%%%

\sectionnewpage
\section{Fakta tentang limit barisan}
\label{sec:factslimsseqs}

%mbxINTROSUBSECTION

\sectionnotes{2--2.5 pertemuan; barisan yang didefinisikan secara rekursif dapat dilewati
tanpa masalah}

Pada bagian ini, kita membahas beberapa hasil dasar mengenai limit
barisan.
Kita mulai dengan meninjau bagaimana barisan berinteraksi dengan ketaksamaan.

\subsection{Limit dan ketaksamaan}

Lemma dasar mengenai limit dan ketaksamaan disebut lemma apit.
Lemma ini memungkinkan kita menunjukkan kekonvergenan barisan dalam kasus-kasus sulit
jika kita menemukan dua barisan konvergen lain yang lebih sederhana dan 
\myquote{mengapit} barisan semula.

\begin{lemma}[Lemma apit]\index{lemma apit} \label{squeeze:lemma}
Misalkan $\{ a_n \}_{n=1}^\infty$, 
$\{ b_n \}_{n=1}^\infty$, dan 
$\{ x_n \}_{n=1}^\infty$ merupakan barisan-barisan sedemikian sehingga
\begin{equation*}
a_n \leq x_n \leq b_n \quad \text{untuk setiap } n \in \N.
\end{equation*}
Andaikan $\{ a_n \}_{n=1}^\infty$ dan $\{ b_n \}_{n=1}^\infty$ konvergen dan
\begin{equation*}
\lim_{n\to \infty} a_n
=
\lim_{n\to \infty} b_n .
\end{equation*}
Maka $\{ x_n \}_{n=1}^\infty$ konvergen dan
\begin{equation*}
\lim_{n\to \infty} x_n
=
\lim_{n\to \infty} a_n
=
\lim_{n\to \infty} b_n .
\end{equation*}
\end{lemma}

\begin{proof}
Misalkan $x \coloneqq \lim_{n\to\infty} a_n = \lim_{n\to\infty} b_n$.
Ambil $\epsilon > 0$.
Pilih $M_1$ sedemikian sehingga untuk setiap $n \geq M_1$, berlaku
$\sabs{a_n-x} < \epsilon$, dan pilih $M_2$
sedemikian sehingga untuk setiap $n \geq M_2$,
berlaku $\sabs{b_n-x} < \epsilon$.  Tetapkan $M \coloneqq \max \{M_1, M_2 \}$.
Andaikan $n \geq M$.
Secara khusus,
$x - a_n < \epsilon$, atau 
$x - \epsilon < a_n$.  Demikian pula, $b_n < x + \epsilon$.
Dengan menggabungkan semuanya, kita peroleh
\begin{equation*}
x - \epsilon < a_n \leq x_n \leq b_n < x + \epsilon .
\end{equation*}
Dengan kata lain, $-\epsilon < x_n-x < \epsilon$ atau $\sabs{x_n-x} < \epsilon$.
Jadi, $\{x_n\}_{n=1}^\infty$ konvergen ke $x$.
Lihat \figureref{figsqueeze}.
\end{proof}

\begin{myfigureht}
\myincludepdft{figsqueeze}{%
Diagram garis bilangan real.  Titik-titik ditandai pada garis tersebut
secara berurutan dari kiri ke kanan: x dikurangi epsilon, a subskrip n, x, x subskrip n,
b subskrip n, dan x ditambah epsilon.  Panah menegaskan bahwa jarak
dari x dikurangi epsilon ke x adalah epsilon dan bahwa jarak
dari x ke x ditambah epsilon juga epsilon.}
\caption{Bukti lemma apit dalam gambar.\label{figsqueeze}}
\end{myfigureht}

\begin{example}
Salah satu penerapan
\hyperref[squeeze:lemma]{lemma apit} adalah menghitung limit 
barisan dengan menggunakan limit yang telah kita ketahui.  Sebagai contoh, perhatikan
barisan $\bigl\{ \frac{1}{n\sqrt{n}} \bigr\}_{n=1}^\infty$.
Karena $\sqrt{n} \geq 1$ untuk setiap $n \in \N$, berlaku
\begin{equation*}
0 \leq \frac{1}{n\sqrt{n}} \leq \frac{1}{n}
\end{equation*}
untuk setiap $n \in \N$.  Kita telah mengetahui $\lim_{n\to\infty} \nicefrac{1}{n} = 0$. 
Oleh karena itu, dengan menggunakan
barisan konstan $\{ 0 \}_{n=1}^\infty$ dan barisan
$\{ \nicefrac{1}{n} \}_{n=1}^\infty$ dalam
lemma apit, kita simpulkan
\begin{equation*}
\lim_{n\to\infty} \frac{1}{n\sqrt{n}} = 0 .
\end{equation*}
\end{example}

Limit, jika ada, mempertahankan ketaksamaan tak ketat.

\begin{lemma} \label{limandineq:lemma}
Misalkan $\{ x_n \}_{n=1}^\infty$ dan $\{ y_n \}_{n=1}^\infty$ merupakan
barisan konvergen dan
\begin{equation*}
x_n \leq y_n \quad \text{untuk setiap } n \in \N .
\end{equation*}
Maka
\begin{equation*}
\lim_{n\to\infty} x_n \leq
\lim_{n\to\infty} y_n .
\end{equation*}
\end{lemma}

\begin{proof}
Misalkan $x \coloneqq \lim_{n\to\infty} x_n$ dan $y \coloneqq \lim_{n\to\infty} y_n$. 
Ambil 
$\epsilon > 0$.  Pilih $M_1$ sedemikian sehingga untuk setiap $n \geq M_1$,
berlaku $\sabs{x_n-x} < \nicefrac{\epsilon}{2}$.  Pilih $M_2$ sedemikian sehingga
untuk setiap $n \geq M_2$, berlaku
$\sabs{y_n-y} < \nicefrac{\epsilon}{2}$.  Secara khusus,
untuk setiap $n \geq \max\{ M_1, M_2 \}$, berlaku
$x-x_n < \nicefrac{\epsilon}{2}$ dan
$y_n-y < \nicefrac{\epsilon}{2}$.  Kita menjumlahkan kedua ketaksamaan ini untuk
memperoleh
\begin{equation*}
y_n-x_n+x-y < \epsilon, \qquad \text{atau} \qquad
y_n-x_n < y-x+ \epsilon .
\end{equation*}
Karena $x_n \leq y_n$, berlaku
$0 \leq y_n-x_n$ dan dengan demikian $0 < y-x+ \epsilon$.
Dengan kata lain,
\begin{equation*}
x-y < \epsilon .
\end{equation*}
Karena $\epsilon > 0$ dipilih secara sembarang, kita peroleh
$x-y \leq 0$.
Oleh karena itu, $x \leq y$.
\end{proof}

Korolari berikut diperoleh dengan
menggunakan barisan konstan dalam
\lemmaref{limandineq:lemma}.  Buktinya diserahkan sebagai latihan.

\begin{samepage}
\begin{cor} \label{limandineq:cor}
\leavevmode
\begin{enumerate}[(i)]
\item Jika $\{ x_n \}_{n=1}^\infty$ merupakan barisan konvergen sedemikian sehingga $x_n \geq 0$ untuk setiap
$n \in \N$,
maka
\begin{equation*}
\lim_{n\to\infty} x_n \geq 0.
\end{equation*}
\item
Misalkan $a,b \in \R$ dan
misalkan $\{ x_n \}_{n=1}^\infty$ suatu barisan konvergen sedemikian sehingga
\begin{equation*}
a \leq x_n \leq b \quad \text{untuk setiap } n \in \N .
\end{equation*}
Maka
\begin{equation*}
a \leq \lim_{n\to\infty} x_n \leq b.
\end{equation*}
\end{enumerate}
\end{cor}
\end{samepage}

Dalam \lemmaref{limandineq:lemma} dan \corref{limandineq:cor}, kita tidak dapat begitu saja mengganti
semua ketaksamaan tak ketat dengan
ketaksamaan ketat.  Sebagai contoh,
misalkan $x_n \coloneqq \nicefrac{-1}{n}$ dan $y_n \coloneqq \nicefrac{1}{n}$.
Maka $x_n < y_n$, $x_n < 0$,
dan $y_n > 0$ untuk setiap $n$.  Namun, ketaksamaan-ketaksamaan ini
tidak dipertahankan oleh operasi limit karena
$\lim_{n\to\infty} x_n = \lim_{n\to\infty} y_n = 0$.
Pesan dari contoh ini adalah bahwa ketaksamaan ketat dapat menjadi ketaksamaan tak ketat
ketika limit diambil; jika kita mengetahui
$x_n < y_n$ untuk setiap $n$,
kita hanya dapat menyimpulkan 
\begin{equation*}
\lim_{n \to \infty} x_n \leq
\lim_{n \to \infty} y_n .
\end{equation*}
Hal ini merupakan sumber kesalahan yang umum.

\subsection{Kekontinuan operasi aljabar}

Limit berperilaku baik terhadap operasi aljabar.

\begin{prop} \label{prop:contalg}
Misalkan $\{ x_n \}_{n=1}^\infty$ dan $\{ y_n \}_{n=1}^\infty$ adalah barisan konvergen.
\begin{enumerate}[(i)]
\item \label{prop:contalg:i}
Barisan $\{ z_n \}_{n=1}^\infty$, dengan $z_n \coloneqq x_n + y_n$, konvergen dan
\begin{equation*}
\lim_{n \to \infty} (x_n + y_n) = 
\lim_{n \to \infty} z_n = 
\lim_{n \to \infty} x_n + 
\lim_{n \to \infty} y_n .
\end{equation*}
\item \label{prop:contalg:ii}
Barisan $\{ z_n \}_{n=1}^\infty$, dengan $z_n \coloneqq x_n - y_n$, konvergen dan
\begin{equation*}
\lim_{n \to \infty} (x_n - y_n) = 
\lim_{n \to \infty} z_n = 
\lim_{n \to \infty} x_n - 
\lim_{n \to \infty} y_n .
\end{equation*}
\item \label{prop:contalg:iii}
Barisan $\{ z_n \}_{n=1}^\infty$, dengan $z_n \coloneqq x_n y_n$, konvergen dan
\begin{equation*}
\lim_{n \to \infty} (x_n y_n) = 
\lim_{n \to \infty} z_n = 
\left( \lim_{n \to \infty} x_n \right)
\left( \lim_{n \to \infty} y_n \right) .
\end{equation*}
\item \label{prop:contalg:iv}
Jika $\lim_{n\to\infty} y_n \neq 0$ dan $y_n \neq 0$ untuk semua $n \in \N$, maka
barisan $\{ z_n \}_{n=1}^\infty$, dengan $z_n \coloneqq \dfrac{x_n}{y_n}$, konvergen dan
\begin{equation*}
\lim_{n \to \infty} \frac{x_n}{y_n} = 
\lim_{n \to \infty} z_n = 
\frac{\lim_{n \to \infty} x_n}{\lim_{n \to \infty} y_n} .
\end{equation*}
\end{enumerate}
\end{prop}

\begin{proof}
Kita mulai dengan \ref{prop:contalg:i}.
Misalkan $\{ x_n \}_{n=1}^\infty$ dan $\{ y_n \}_{n=1}^\infty$ adalah barisan konvergen dan
tulis $z_n \coloneqq x_n + y_n$.  Misalkan $x \coloneqq \lim_{n\to\infty} x_n$,
$y \coloneqq \lim_{n\to\infty} y_n$, dan $z \coloneqq x+y$.
Diberikan sebarang $\epsilon > 0$.  
Carilah suatu $M_1$ sedemikian sehingga untuk semua $n \geq M_1$,
berlaku
$\sabs{x_n - x} < \nicefrac{\epsilon}{2}$.  
Carilah suatu $M_2$ sedemikian sehingga untuk semua $n \geq M_2$,
berlaku
$\sabs{y_n - y} < \nicefrac{\epsilon}{2}$.  Ambil $M \coloneqq \max \{ M_1, M_2 \}$.
Untuk semua $n \geq M$, berlaku
\begin{equation*}
\begin{split}
\sabs{z_n - z} &=
\babs{(x_n+y_n) - (x+y)} \\
& =
\sabs{x_n-x + y_n-y} \\
& \leq
\sabs{x_n-x} + \sabs{y_n-y} \\
& <
\frac{\epsilon}{2} +
\frac{\epsilon}{2}
= \epsilon.
\end{split}
\end{equation*}
Dengan demikian, \ref{prop:contalg:i} terbukti.
Bukti \ref{prop:contalg:ii} hampir identik dan diserahkan kepada pembaca sebagai
latihan.

Sekarang kita tangani 
\ref{prop:contalg:iii}.
Misalkan lagi bahwa $\{ x_n \}_{n=1}^\infty$ dan $\{ y_n \}_{n=1}^\infty$ adalah barisan konvergen dan
tulis $z_n \coloneqq x_n y_n$.  Misalkan $x \coloneqq \lim_{n\to\infty} x_n$,
$y \coloneqq \lim_{n\to\infty} y_n$, dan $z \coloneqq xy$.
Diberikan sebarang $\epsilon > 0$.
Misalkan $K \coloneqq \max\{ \sabs{x}, \sabs{y}, \nicefrac{\epsilon}{3} , 1 \}$.
Carilah suatu $M_1$ sedemikian sehingga untuk semua $n \geq M_1$,
berlaku
$\sabs{x_n - x} < \frac{\epsilon}{3K}$.
Carilah suatu $M_2$ sedemikian sehingga untuk semua $n \geq M_2$,
berlaku
$\sabs{y_n - y} < \frac{\epsilon}{3K}$.  Ambil $M \coloneqq \max \{ M_1, M_2 \}$.
Untuk semua $n \geq M$, berlaku
\begin{equation*}
\begin{split}
\sabs{z_n - z} &=
\babs{(x_ny_n) - (xy)} \\
& =
\babs{(x_n-x+x)(y_n-y+y) - xy} \\
& =
\babs{(x_n-x)y + x(y_n-y) +(x_n-x)(y_n-y)} \\
& \leq
\babs{(x_n-x)y} + \babs{x(y_n - y)} +
\babs{(x_n-x)(y_n-y)} \\
& =
\sabs{x_n -x}\sabs{y} + 
\sabs{x}\sabs{y_n -y} + 
\sabs{x_n -x}\sabs{y_n -y}
\\
& <
\frac{\epsilon}{3K} K + 
K \frac{\epsilon}{3K} + 
\frac{\epsilon}{3K}
\frac{\epsilon}{3K}
\qquad \qquad \text{(sekarang perhatikan bahwa } \tfrac{\epsilon}{3K} \leq 1
\text{ dan }
K \geq 1\text{)}
\\
& \leq
\frac{\epsilon}{3} + \frac{\epsilon}{3} + \frac{\epsilon}{3}
 = \epsilon .
\end{split}
\end{equation*}

Akhirnya, kita tinjau
\ref{prop:contalg:iv}.
Kita buktikan klaim yang lebih sederhana berikut:

Klaim: \emph{Jika $\{ y_n \}_{n=1}^\infty$ adalah barisan konvergen sedemikian sehingga
$\lim_{n\to\infty} y_n \neq 0$ dan $y_n \neq 0$ untuk semua $n \in \N$, maka
$\{ \nicefrac{1}{y_n} \}_{n=1}^\infty$ konvergen dan}
\begin{equation*}
\lim_{n\to\infty} \frac{1}{y_n} = \frac{1}{\lim_{n\to\infty} y_n}  .
\end{equation*}

Setelah klaim terbukti, kita ambil barisan
$\{ \nicefrac{1}{y_n} \}_{n=1}^\infty$,
mengalikannya dengan barisan $\{ x_n \}_{n=1}^\infty$, lalu menerapkan butir
\ref{prop:contalg:iii}.

Bukti klaim:  Diberikan sebarang $\epsilon > 0$.
Misalkan $y \coloneqq \lim_{n\to\infty} y_n$.
Karena $\sabs{y} \neq 0$, berlaku
$\min \left\{ \sabs{y}^2\frac{\epsilon}{2}, \, \frac{\sabs{y}}{2} \right\} > 0$.
Carilah suatu $M$ sedemikian sehingga untuk semua $n \geq M$,
berlaku
\begin{equation*}
\sabs{y_n - y} < \min \left\{ \sabs{y}^2\frac{\epsilon}{2}, \, \frac{\sabs{y}}{2}
\right\} .
\end{equation*}
Untuk semua $n \geq M$, berlaku
$\sabs{y - y_n} < \nicefrac{\sabs{y}}{2}$, sehingga
\begin{equation*}
\sabs{y} = 
\sabs{y - y_n + y_n } \leq
\sabs{y - y_n} + \sabs{ y_n } < \frac{\sabs{y}}{2} + \sabs{y_n}.
\end{equation*}
Dengan mengurangkan $\nicefrac{\sabs{y}}{2}$ dari kedua ruas, kita peroleh
$\nicefrac{\sabs{y}}{2} < \sabs{y_n}$, atau dengan kata lain,
\begin{equation*}
\frac{1}{\sabs{y_n}} < \frac{2}{\sabs{y}} .
\end{equation*}
Kita selesaikan bukti klaim tersebut:
\begin{equation*}
\begin{split}
\abs{\frac{1}{y_n} - \frac{1}{y}} &=
\abs{\frac{y - y_n}{y y_n}} \\
& =
\frac{\sabs{y - y_n}}{\sabs{y} \sabs{y_n}} \\
& \leq
\frac{\sabs{y - y_n}}{\sabs{y}} \, \frac{2}{\sabs{y}} \\
& <
\frac{\sabs{y}^2 \frac{\epsilon}{2}}{\sabs{y}} \, \frac{2}{\sabs{y}}
= \epsilon .
\end{split}
\avoidbreak
\end{equation*}
Dengan demikian, bukti selesai.
\end{proof}

Dengan menyubstitusikan barisan konstan, kita memperoleh beberapa akibat yang sederhana.
Jika $c \in \R$ dan $\{ x_n \}_{n=1}^\infty$ adalah barisan konvergen, maka
misalnya
\begin{equation*}
\lim_{n \to \infty} c x_n = 
c \left( \lim_{n \to \infty} x_n \right) \qquad
\text{dan}
\qquad
\lim_{n \to \infty} (c + x_n) = 
c + \lim_{n \to \infty} x_n .
\end{equation*}
Dengan cara serupa, kita memperoleh kesamaan semacam itu untuk pengurangan dan pembagian dengan konstanta.

Karena limit dapat dipertukarkan dengan perkalian, kita dapat menunjukkan (latihan)
bahwa $\lim_{n\to\infty} x_n^k = {(\lim_{n\to\infty} x_n)}^k$ untuk semua $k \in \N$.
Artinya, limit dapat dipertukarkan
dengan perpangkatan.  Mari kita lihat apakah hal yang sama berlaku untuk akar.

\begin{prop}
\pagebreak[2]
Misalkan $\{ x_n \}_{n=1}^\infty$ adalah barisan konvergen sedemikian
sehingga $x_n \geq 0$ untuk semua $n \in \N$.
Maka
\begin{equation*}
\lim_{n\to\infty} \sqrt{x_n} =
\sqrt{ \lim_{n\to\infty} x_n } .
\end{equation*}
\end{prop}

Tentu saja, untuk dapat menyatakan hal ini, kita perlu menerapkan
\corref{limandineq:cor} untuk menunjukkan
bahwa
$\lim_{n\to\infty} x_n \geq 0$, sehingga kita dapat mengambil akar kuadrat tanpa
masalah.

\begin{proof}
Misalkan $\{ x_n \}_{n=1}^\infty$ adalah barisan konvergen dan misalkan $x \coloneqq
\lim_{n\to\infty} x_n$.
Seperti baru saja disebutkan, $x \geq 0$.

Pertama, misalkan $x=0$.  Diberikan sebarang $\epsilon > 0$.
Maka terdapat suatu $M$ sedemikian sehingga untuk semua $n \geq M$, berlaku
$x_n = \sabs{x_n} < \epsilon^2$, atau dengan kata lain, $\sqrt{x_n} < \epsilon$.
Jadi,
\begin{equation*}
\abs{\sqrt{x_n} - \sqrt{x}} =
\sqrt{x_n} < \epsilon.
\end{equation*}

Sekarang misalkan $x > 0$ (dan karenanya $\sqrt{x} > 0$).
\begin{equation*}
\begin{split}
\abs{\sqrt{x_n}-\sqrt{x}} &= 
\abs{\frac{x_n-x}{\sqrt{x_n}+\sqrt{x}}} \\
&=
\frac{1}{\sqrt{x_n}+\sqrt{x}}
\sabs{x_n-x} \\
& \leq
\frac{1}{\sqrt{x}}
\sabs{x_n-x} .
\end{split}
\end{equation*}
Sisa bukti diserahkan kepada pembaca.
\end{proof}

Bukti serupa berlaku untuk akar ke-$k$.  Artinya, kita juga
memperoleh
$\lim_{n\to\infty} x_n^{1/k} = {( \lim_{n\to\infty} x_n )}^{1/k}$.  Hal ini diserahkan kepada pembaca
sebagai latihan yang menantang.

Kita juga mungkin ingin mempertukarkan limit dengan tanda nilai mutlak.
Kebalikan proposisi ini tidak benar; lihat
\exerciseref{exercise:absconv} bagian b).

\begin{prop}
Jika $\{ x_n \}_{n=1}^\infty$ adalah barisan konvergen, maka $\{ \sabs{x_n} \}_{n=1}^\infty$
konvergen dan
\begin{equation*}
\lim_{n\to\infty} \sabs{x_n} = 
\abs{\lim_{n\to\infty} x_n} .
\end{equation*}
\end{prop}

\begin{proof}
Kita cukup memperhatikan ketaksamaan segitiga terbalik
\begin{equation*}
\big\lvert \sabs{x_n} - \sabs{x} \big\rvert \leq \sabs{x_n-x} .
\end{equation*}
Jadi, jika $\sabs{x_n -x}$ dapat dibuat sekecil apa pun, hal yang sama berlaku untuk
$\big\lvert \sabs{x_n} - \sabs{x} \big\rvert$.
Rinciannya diserahkan kepada pembaca.
\end{proof}

Mari kita lihat sebuah contoh yang menggabungkan proposisi-proposisi di atas.  Karena
$\lim_{n\to\infty} \nicefrac{1}{n} = 0$, maka
\begin{equation*}
\lim_{n\to \infty}
\abs{\sqrt{1 + \nicefrac{1}{n}} - \nicefrac{100}{n^2}} =  
\abs{\sqrt{1 + \Bigl(\lim_{n\to\infty} \nicefrac{1}{n}\Bigr)}
- 100 \Bigl(\lim_{n\to\infty} \nicefrac{1}{n}\Bigr)\Bigl(\lim_{n\to\infty} \nicefrac{1}{n}\Bigr)} = 1.
\end{equation*}
Artinya, limit pada ruas kiri ada karena ruas kanan
ada.  Kesamaan di atas sebenarnya harus dibaca dari kanan ke kiri.

Di sisi lain, proposisi-proposisi tersebut harus diterapkan dengan hati-hati.
Sebagai contoh, dengan terlebih dahulu menulis ulang bentuk tersebut menggunakan penyebut bersama,
kita peroleh
\begin{equation*}
\lim_{n\to \infty} \left( \frac{n^2}{n+1} - n \right)
= -1 .
\end{equation*}
Namun, 
$\bigl\{ \frac{n^2}{n+1} \bigr\}_{n=1}^\infty$ dan 
$\{n\}_{n=1}^\infty$ tidak konvergen,
sehingga
$\Bigl(\lim\limits_{n\to\infty} \frac{n^2}{n+1}\Bigr) -
\Bigl(\lim\limits_{n\to\infty} n\Bigr)$ tidak bermakna.

\subsection{Barisan yang didefinisikan secara rekursif}

Setelah mengetahui bahwa kita dapat menukar limit dan operasi aljabar, kita dapat
menghitung limit banyak barisan.
Salah satu kelasnya adalah barisan yang didefinisikan secara rekursif, yaitu barisan yang
bilangan berikutnya dihitung dengan suatu rumus dari sejumlah tetap
unsur sebelumnya dalam barisan tersebut.

\begin{example}
Misalkan $\{ x_n \}_{n=1}^\infty$ didefinisikan oleh $x_1 \coloneqq 2$ dan
\begin{equation*}
x_{n+1} \coloneqq x_n - \frac{x_n^2-2}{2x_n} .
\end{equation*}
Pertama-tama, kita harus memastikan bahwa barisan ini terdefinisi dengan baik; kita harus menunjukkan bahwa kita tidak pernah
membagi dengan nol.
Kemudian, kita harus memastikan apakah barisan tersebut konvergen.  Barulah setelah itu
kita dapat mencoba menentukan limitnya.

Mari kita buktikan bahwa untuk semua $n$,
$x_n$ ada dan $x_n > 0$
(sehingga barisan tersebut terdefinisi dengan baik dan terbatas di bawah).
Kita tunjukkan hal ini dengan \hyperref[induction:thm]{induksi}.  Kita tahu bahwa
$x_1 = 2 > 0$.  Untuk langkah induksi, andaikan $x_n$ ada dan
$x_n > 0$.  Maka
\begin{equation*}
x_{n+1} = x_n - \frac{x_n^2-2}{2x_n} =
\frac{2x_n^2 - x_n^2+2}{2x_n} =
\frac{x_n^2+2}{2x_n} .
\end{equation*}
Selalu berlaku bahwa $x_n^2+2 > 0$,
dan karena
$x_n > 0$, maka $x_{n+1} = \frac{x_n^2+2}{2x_n} > 0$.

Selanjutnya, kita
tunjukkan bahwa barisan tersebut monoton turun.  Jika kita menunjukkan bahwa
$x_n^2-2 \geq 0$ untuk semua $n$, maka $x_{n+1} \leq x_n$ untuk semua $n$.
Jelas bahwa $x_1^2-2 = 4-2 = 2 > 0$.  Untuk sebarang $n$, kita memiliki
\begin{equation*}
x_{n+1}^2-2 =
{\left( \frac{x_n^2+2}{2x_n} \right)}^2 - 2
=
\frac{x_n^4+4x_n^2+4 - 8x_n^2}{4x_n^2}
=
\frac{x_n^4-4x_n^2+4}{4x_n^2}
=
\frac{{\left( x_n^2-2 \right)}^2}{4x_n^2} .
\end{equation*}
Karena kuadrat selalu tak negatif,
$x_{n+1}^2-2 \geq 0$ untuk semua $n$.  Oleh karena itu,
$\{ x_n \}_{n=1}^\infty$ monoton turun dan terbatas ($x_n > 0$ untuk semua $n$), sehingga
limitnya ada.  Tinggal menentukan limit tersebut.

Tuliskan
\begin{equation*}
2x_nx_{n+1} = x_n^2+2 .
\end{equation*}
Karena $\{ x_{n+1} \}_{n=1}^\infty$ merupakan ekor-1 dari $\{ x_n \}_{n=1}^\infty$, barisan tersebut konvergen ke
limit yang sama.  Definisikan $x \coloneqq \lim_{n\to\infty} x_n$.  Dengan mengambil limit pada
kedua ruas, kita memperoleh
\begin{equation*}
2x^2 = x^2+2 ,
\end{equation*}
atau $x^2 = 2$.  Karena $x_n > 0$ untuk semua $n$, kita memperoleh $x \geq 0$, sehingga $x = \sqrt{2}$.
\end{example}

Anda mungkin pernah melihat barisan di atas.  Barisan tersebut merupakan
\emph{metode Newton}%
\footnote{%
Nama ini diambil dari fisikawan dan matematikawan Inggris
\href{https://en.wikipedia.org/wiki/Isaac_Newton}{Isaac Newton}
(1642--1726/7).}
untuk menentukan akar kuadrat dari 2.  Metode ini sering muncul dalam
praktik dan konvergen dengan sangat cepat.  Kita menggunakan fakta bahwa
$x_1^2 -2 >0$, meskipun fakta ini sebenarnya tidak diperlukan untuk menunjukkan kekonvergenan jika kita
meninjau suatu ekor barisan.
Barisan tersebut konvergen selama $x_1 \neq 0$, meskipun dengan $x_1$ negatif
kita akan memperoleh $x=-\sqrt{2}$.  Dengan mengganti angka 2 pada pembilang, kita
memperoleh akar kuadrat dari sebarang bilangan positif.  Pernyataan-pernyataan ini diserahkan sebagai
latihan.

Namun, Anda harus berhati-hati.  Sebelum mengambil limit apa pun, Anda harus
memastikan bahwa barisan tersebut konvergen.  Mari kita lihat sebuah contoh.

\begin{example}
Andaikan $x_1 \coloneqq 1$ dan $x_{n+1} \coloneqq x_n^2+x_n$.
Jika kita begitu saja menganggap bahwa limitnya ada (sebut limit itu $x$), maka kita
akan memperoleh persamaan $x = x^2+x$, yang mungkin membuat kita
menyimpulkan bahwa $x=0$.  Namun, tidak sulit
untuk menunjukkan bahwa $\{ x_n \}_{n=1}^\infty$ tak terbatas dan karenanya tidak konvergen.

Hal yang perlu diperhatikan dalam contoh ini adalah bahwa metode tersebut tetap berlaku, tetapi
hasilnya bergantung pada nilai awal $x_1$.  Jika kita menetapkan $x_1 \coloneqq 0$,
maka barisan tersebut konvergen dan limitnya memang 0.
Satu cabang matematika yang disebut dinamika membahas tepat persoalan-persoalan
semacam ini.
Lihat
\exerciseref{exercise:convergentinitialvalues}.
\end{example}

\subsection{Beberapa uji konvergensi}

Kita tidak selalu perlu kembali ke definisi kekonvergenan
untuk membuktikan bahwa suatu barisan konvergen.  Pertama-tama, kita berikan sebuah uji konvergensi sederhana.
Gagasan utamanya adalah bahwa
$\{ x_n \}_{n=1}^\infty$ konvergen ke $x$ jika dan hanya jika
$\bigl\{ \sabs{ x_n - x } \bigr\}_{n=1}^\infty$ konvergen ke nol.

\begin{prop} \label{convzero:prop}
Misalkan $\{ x_n \}_{n=1}^\infty$ suatu barisan.
Andaikan terdapat $x \in \R$
dan suatu barisan konvergen $\{ a_n \}_{n=1}^\infty$
sedemikian sehingga
\begin{equation*}
\lim_{n\to\infty} a_n = 0
\end{equation*}
dan
\begin{equation*}
\sabs{x_n - x} \leq a_n
\qquad
\text{untuk semua $n \in \N$.}
\end{equation*}
Maka $\{ x_n \}_{n=1}^\infty$ konvergen dan $\lim\limits_{n\to\infty} x_n = x$.
\end{prop}

\begin{proof}
Ambil $\epsilon > 0$.  Perhatikan bahwa $a_n \geq 0$
untuk semua $n$.  Pilih $M \in \N$ sedemikian sehingga untuk
semua $n \geq M$, berlaku
$a_n = \sabs{a_n - 0} < \epsilon$.  Maka, untuk semua $n \geq M$,
kita memiliki
\begin{equation*}
\sabs{x_n - x} \leq a_n < \epsilon . \qedhere
\end{equation*}
\end{proof}

Dengan kata lain, mempelajari kapan suatu barisan memiliki limit sama
dengan mempelajari kapan barisan lain menuju nol.
Menentukan apakah suatu barisan tertentu konvergen mungkin sulit, tetapi
untuk beberapa barisan terdapat uji kekonvergenan yang mudah diterapkan.
Mari kita lihat salah satu uji tersebut.
Pertama-tama, kita hitung limit suatu barisan khusus.

\begin{prop}
Misalkan $c > 0$.
\begin{enumerate}[(i)]
\item
Jika $c < 1$, maka
\begin{equation*}
\lim_{n\to\infty} c^n = 0.
\end{equation*}
\item
Jika $c > 1$, maka $\{ c^n \}_{n=1}^\infty$ tak terbatas.
\end{enumerate}
\end{prop}

\begin{proof}
Pertama, tinjau $c < 1$.  Karena $c > 0$, kita memperoleh $c^n > 0$ untuk semua $n \in \N$ melalui
\hyperref[induction:thm]{induksi}.  Kemudian, $c < 1$ mengakibatkan bahwa
$c^{n+1} < c^n$ untuk semua $n$.
Jadi, barisan $\{ c^n \}_{n=1}^\infty$
terbatas di bawah dan
menurun.  Oleh karena itu, barisan tersebut konvergen.
Tetapkan $x \coloneqq \lim_{n\to\infty} c^n$.  Ekor-1 $\{ c^{n+1} \}_{n=1}^\infty$
juga konvergen ke $x$.  Dengan mengambil limit kedua ruas dari $c^{n+1} = c \cdot
c^n$, kita memperoleh $x = cx$, atau $(1-c)x=0$.  Jadi, $x=0$ karena $c \neq 1$.

Sekarang, tinjau $c > 1$.
Ambil sebarang $B > 0$.
Karena $\nicefrac{1}{c} < 1$, barisan $\bigl\{ {(\nicefrac{1}{c})}^n \bigr\}_{n=1}^\infty$
konvergen ke $0$.
Oleh karena itu, untuk suatu $n$ yang cukup besar, kita memperoleh
\begin{equation*}
\frac{1}{c^n} =
{\left(\frac{1}{c}\right)}^n < \frac{1}{B} .
\end{equation*}
Dengan kata lain, $c^n > B$, sehingga $B$ bukan batas atas bagi
$\{ c^n \}_{n=1}^\infty$.  Karena $B$ sebarang, $\{ c^n \}_{n=1}^\infty$ tak terbatas.
\end{proof}

Dalam proposisi di atas,
rasio suku ke-$(n+1)$ terhadap suku ke-$n$ adalah $c$.  Kita
memperumum hasil sederhana ini ke kelas barisan yang lebih luas.
Lema berikut akan muncul kembali ketika kita membahas deret.

\begin{lemma}[Uji rasio untuk barisan]\index{uji rasio untuk barisan}
\label{seq:ratiotest}
Misalkan $\{ x_n \}_{n=1}^\infty$ suatu barisan sedemikian sehingga $x_n \neq 0$ untuk semua $n$ dan
limit berikut
\begin{equation*}
L \coloneqq \lim_{n\to\infty} \frac{\sabs{x_{n+1}}}{\sabs{x_n}}
\qquad
\text{ada.}
\end{equation*}
\begin{enumerate}[(i)]
\item
Jika $L < 1$, maka $\{ x_n \}_{n=1}^\infty$ konvergen dan $\lim\limits_{n\to\infty} x_n = 0$.
\item
Jika $L > 1$, maka $\{ x_n \}_{n=1}^\infty$ tak terbatas (dan karenanya divergen).
\end{enumerate}
\end{lemma}

Jika $L$ ada, tetapi $L=1$, lema tersebut tidak memberikan informasi apa pun.  Kita tidak dapat menarik
kesimpulan hanya berdasarkan informasi itu.  Sebagai contoh,
barisan $\{ \nicefrac{1}{n} \}_{n=1}^\infty$ konvergen ke nol, tetapi $L=1$.
Barisan konstan $\{ 1 \}_{n=1}^\infty$ konvergen ke 1, bukan nol, dan
$L=1$.  Barisan $\bigl\{ {(-1)}^n \bigr\}_{n=1}^\infty$ sama sekali tidak konvergen, dan
$L=1$ pula.
Terakhir, barisan $\{  n \}_{n=1}^\infty$ tak terbatas, dan sekali lagi $L=1$.
Pernyataan lema tersebut dapat sedikit diperkuat; lihat
latihan \ref{exercise:strongerratiotest1} dan
\ref{exercise:strongerratiotest2}.

\begin{proof}
Andaikan $L < 1$.
Karena
$\frac{\sabs{x_{n+1}}}{\sabs{x_n}} \geq 0$ untuk semua $n$, kita memiliki $L \geq 0$.  Pilih
$r$ sedemikian sehingga $L < r < 1$.
Kita hendak membandingkan barisan $\{ x_n \}_{n=1}^\infty$ dengan barisan $\{ r^n \}_{n=1}^\infty$.  Gagasannya adalah bahwa
meskipun rasio $\frac{\sabs{x_{n+1}}}{\sabs{x_n}}$ belum tentu pada akhirnya lebih kecil daripada $L$,
rasio tersebut pada akhirnya akan lebih kecil daripada $r$, yang masih lebih kecil daripada 1.
Gagasan intuitif bukti ini diilustrasikan dalam \figureref{figratseq}.
\begin{myfigureht}
\myincludepdft{figratseq}{%
Diagram garis bilangan real dengan titik L di sebelah kiri, kemudian
sedikit ke kanan terdapat titik r, dan lebih jauh ke kanan terdapat titik 1.
Garis-garis pendek menandai titik-titik yang mengumpul di sekitar L dan semuanya lebih kecil
daripada r.}
\caption{Garis-garis pendek menyatakan rasio
$\frac{\sabs{x_{n+1}}}{\sabs{x_n}}$
yang mendekati $L < 1$.\label{figratseq}}
\end{myfigureht}

Karena $r-L > 0$, terdapat $M \in \N$ sedemikian sehingga untuk
semua $n \geq M$, kita memiliki
\begin{equation*}
\abs{\frac{\sabs{x_{n+1}}}{\sabs{x_n}} - L} < r-L .
\end{equation*}
Oleh karena itu, untuk $n \geq M$,
\begin{equation*}
\frac{\sabs{x_{n+1}}}{\sabs{x_n}} - L < r-L 
\qquad \text{atau} \qquad
\frac{\sabs{x_{n+1}}}{\sabs{x_n}} < r .
\end{equation*}
Untuk $n > M$ (artinya untuk $n \geq M+1$),
tuliskan
\begin{equation*}
\sabs{x_n} =
\sabs{x_M}
\frac{\sabs{x_{M+1}}}{\sabs{x_{M}}}
\frac{\sabs{x_{M+2}}}{\sabs{x_{M+1}}}
\cdots
\frac{\sabs{x_{n}}}{\sabs{x_{n-1}}}
<
\sabs{x_M}
r r \cdots r = \sabs{x_M} r^{n-M} = (\sabs{x_M} r^{-M}) r^n .
\end{equation*}
Barisan $\{ r^n \}_{n=1}^\infty$ konvergen ke nol sehingga
$\sabs{x_M} r^{-M} r^n$ konvergen ke nol.  Berdasarkan \propref{convzero:prop},
ekor-$M$
$\{x_n\}_{n=M+1}^\infty$ konvergen ke nol, sehingga $\{x_n\}_{n=1}^\infty$ konvergen ke nol.

Sekarang, andaikan $L > 1$.  Pilih
$r$ sedemikian sehingga $1 < r < L$.  Karena $L-r > 0$,
terdapat $M \in \N$ sedemikian sehingga untuk
semua $n \geq M$
\begin{equation*}
\abs{\frac{\sabs{x_{n+1}}}{\sabs{x_n}} - L} < L-r .
\end{equation*}
Oleh karena itu,
\begin{equation*}
\frac{\sabs{x_{n+1}}}{\sabs{x_n}} > r .
\end{equation*}
Sekali lagi, untuk $n > M$,
tuliskan
\begin{equation*}
\sabs{x_n} =
\sabs{x_M}
\frac{\sabs{x_{M+1}}}{\sabs{x_{M}}}
\frac{\sabs{x_{M+2}}}{\sabs{x_{M+1}}}
\cdots
\frac{\sabs{x_{n}}}{\sabs{x_{n-1}}}
>
\sabs{x_M}
r r \cdots r = \sabs{x_M} r^{n-M} = (\sabs{x_M} r^{-M}) r^n .
\end{equation*}
Barisan $\{ r^n \}_{n=1}^\infty$ tak terbatas (karena $r > 1$), sehingga
$\{x_n\}_{n=1}^\infty$ tidak mungkin terbatas (jika $\sabs{x_n} \leq B$ untuk semua $n$, maka
$r^n < \frac{B}{\sabs{x_M}} r^{M}$ untuk semua $n > M$, yang mustahil).
Akibatnya, $\{ x_n \}_{n=1}^\infty$ tidak mungkin konvergen.
\end{proof}

\begin{example}
Salah satu penerapan sederhana dari lema di atas adalah membuktikan
\begin{equation*}
\lim_{n\to\infty} \frac{2^n}{n!} = 0 .
\end{equation*}

Bukti:
Hitung
\begin{equation*}
\frac{2^{n+1} / (n+1)!}{2^n/n!}
=
\frac{2^{n+1}}{2^n}\frac{n!}{(n+1)!}
=
\frac{2}{n+1} .
\end{equation*}
Tidak sulit melihat bahwa $\bigl\{ \frac{2}{n+1} \bigr\}_{n=1}^\infty$ konvergen ke nol.
Kesimpulan tersebut mengikuti lema di atas.
\end{example}

\begin{example} \label{example:nto1overn}
Penerapan uji rasio yang lebih rumit (dan berguna) adalah membuktikan
\begin{equation*}
\lim_{n\to\infty} n^{1/n} = 1 .
\end{equation*}

Bukti:
Ambil $\epsilon > 0$.  Tinjau barisan
$\bigl\{ \frac{n}{{(1+\epsilon)}^n} \bigr\}_{n=1}^\infty$.  Hitung
\begin{equation*}
\frac{(n+1)/{(1+\epsilon)}^{n+1}}{n/{(1+\epsilon)}^{n}}
=
\frac{n+1}{n} \frac{1}{1+\epsilon} .
\end{equation*}
Limit dari $\frac{n+1}{n} = 1+\frac{1}{n}$ ketika $n \to \infty$ adalah 1, sehingga
\begin{equation*}
\lim_{n\to \infty} \frac{(n+1)/{(1+\epsilon)}^{n+1}}{n/{(1+\epsilon)}^{n}}
=
\frac{1}{1+\epsilon}  < 1 .
\end{equation*}
Oleh karena itu, $\bigl\{ \frac{n}{{(1+\epsilon)}^n} \bigr\}_{n=1}^\infty$ konvergen ke 0.
Khususnya,
terdapat $M$ sedemikian sehingga untuk $n \geq M$, kita memiliki
$\frac{n}{{(1+\epsilon)}^n} < 1$, atau
$n < {(1+\epsilon)}^n$, atau
$n^{1/n} < 1+\epsilon$.  Karena $n \geq 1$, kita memperoleh $n^{1/n} \geq 1$, sehingga
$0 \leq n^{1/n}-1 < \epsilon$. Akibatnya,
$\lim\limits_{n\to\infty} n^{1/n} = 1$.
\end{example}

\subsection{Latihan}

\begin{exercise}
Buktikan \corref{limandineq:cor}.  Petunjuk: Gunakan barisan konstan
dan \lemmaref{limandineq:lemma}.
\end{exercise}

\begin{exercise}
Buktikan bagian \ref{prop:contalg:ii} dari \propref{prop:contalg}.
\end{exercise}

\begin{exercise}
Buktikan bahwa jika $\{ x_n \}_{n=1}^\infty$ adalah barisan konvergen dan $k \in \N$, maka
\begin{equation*}
\lim_{n\to\infty} x_n^k = 
{\left( \lim_{n\to\infty} x_n \right)}^k .
\end{equation*}
Petunjuk: Gunakan \hyperref[induction:thm]{induksi}.
\end{exercise}

\begin{exercise}
Misalkan $x_1 \coloneqq \frac{1}{2}$ dan $x_{n+1} \coloneqq x_n^2$.  Tunjukkan bahwa
$\{ x_n \}_{n=1}^\infty$ konvergen dan tentukan
$\lim_{n\to\infty} x_n$.  Petunjuk: Anda tidak boleh membagi dengan nol!
\end{exercise}

\begin{exercise}
Misalkan $x_n \coloneqq \frac{n-\cos(n)}{n}$.  Gunakan
\hyperref[squeeze:lemma]{lemma apit} untuk menunjukkan bahwa
$\{ x_n \}_{n=1}^\infty$ konvergen dan tentukan limitnya.
\end{exercise}

\begin{exercise}
Misalkan $x_n \coloneqq \frac{1}{n^2}$ dan $y_n \coloneqq \frac{1}{n}$.  Definisikan
$z_n \coloneqq \frac{x_n}{y_n}$ dan 
$w_n \coloneqq \frac{y_n}{x_n}$.  Apakah $\{ z_n \}_{n=1}^\infty$ dan $\{ w_n \}_{n=1}^\infty$
konvergen?  Berapakah limitnya?  Dapatkah Anda menerapkan \propref{prop:contalg}?
Mengapa atau mengapa tidak?
\end{exercise}

\begin{exercise}
Benar atau salah, buktikan atau berikan contoh tandingan:
Jika $\{ x_n \}_{n=1}^\infty$ adalah barisan
sedemikian sehingga $\{ x_n^2 \}_{n=1}^\infty$ konvergen, maka $\{ x_n \}_{n=1}^\infty$ konvergen.
\end{exercise}

\begin{exercise}
Tunjukkan bahwa
\begin{equation*}
\lim_{n\to\infty} \frac{n^2}{2^n} = 0 .
\end{equation*}
\end{exercise}

\begin{exercise}
Misalkan $\{ x_n \}_{n=1}^\infty$ adalah barisan, $x \in \R$, dan $x_n \neq x$ untuk
semua $n \in \N$.  Misalkan limit
\begin{equation*}
L \coloneqq \lim_{n \to \infty} \frac{\sabs{x_{n+1}-x}}{\sabs{x_n-x}}
\end{equation*}
ada dan $L < 1$.  Tunjukkan bahwa $\{ x_n \}_{n=1}^\infty$ konvergen ke $x$.
\end{exercise}

\begin{exercise}[Menantang]
Misalkan $\{ x_n \}_{n=1}^\infty$ adalah barisan konvergen sedemikian
sehingga $x_n \geq 0$ dan $k \in \N$.
Tunjukkan
\begin{equation*}
\lim_{n\to\infty} x_n^{1/k} =
{\left( \lim_{n\to\infty} x_n \right)}^{1/k} .
\end{equation*}
Petunjuk: Pertama, tangani kasus ketika limitnya nol.
Selanjutnya, carilah suatu bentuk $q$ sedemikian sehingga $\frac{x_n^{1/k}-x^{1/k}}{x_n-x} =
\frac{1}{q}$.
\end{exercise}

\begin{exercise}
Misalkan $r > 0$.  Tunjukkan bahwa dengan memulai dari sebarang $x_1 \neq 0$, barisan
yang didefinisikan oleh
\begin{equation*}
x_{n+1} \coloneqq x_n - \frac{x_n^2-r}{2x_n}
\end{equation*}
konvergen ke $\sqrt{r}$ jika $x_1 > 0$ dan ke $-\sqrt{r}$ jika $x_1 < 0$.
\end{exercise}

\begin{samepage}
\begin{exercise}
Misalkan $\{ a_n \}_{n=1}^\infty$ dan $\{ b_n \}_{n=1}^\infty$ adalah barisan.
\begin{enumerate}[a)]
\item
Misalkan $\{ a_n \}_{n=1}^\infty$ terbatas dan $\{ b_n \}_{n=1}^\infty$ konvergen ke 0.
Tunjukkan bahwa $\{ a_n b_n \}_{n=1}^\infty$ konvergen ke 0.
\item
Berikan contoh ketika $\{ a_n \}_{n=1}^\infty$ tak terbatas, $\{ b_n \}_{n=1}^\infty$ konvergen ke
0, dan $\{ a_n b_n \}_{n=1}^\infty$ tidak konvergen.
\item
Berikan contoh ketika $\{ a_n \}_{n=1}^\infty$ terbatas, $\{ b_n \}_{n=1}^\infty$ konvergen ke
suatu $x \neq 0$, dan $\{ a_n b_n \}_{n=1}^\infty$ tidak konvergen.
\end{enumerate}
\end{exercise}
\end{samepage}

\begin{samepage}
\begin{exercise}[Mudah] \label{exercise:strongerratiotest1}
Buktikan versi yang lebih kuat berikut dari \lemmaref{seq:ratiotest}, yaitu uji
rasio.  
Misalkan $\{ x_n \}_{n=1}^\infty$ adalah barisan sedemikian sehingga $x_n \neq 0$ untuk semua
$n$.
\begin{enumerate}[a)]
\item
Buktikan bahwa jika terdapat $r < 1$ dan $M \in \N$ sedemikian sehingga
\begin{equation*}
\frac{\sabs{x_{n+1}}}{\sabs{x_n}} \leq r \quad \text{untuk semua } n \geq M,
\end{equation*}
maka $\{ x_n \}_{n=1}^\infty$ konvergen ke $0$.
\item
Buktikan bahwa jika terdapat $r > 1$ dan $M \in \N$ sedemikian sehingga
\begin{equation*}
\frac{\sabs{x_{n+1}}}{\sabs{x_n}} \geq r  \quad \text{untuk semua } n \geq M,
\end{equation*}
maka $\{ x_n \}_{n=1}^\infty$ tak terbatas.
\end{enumerate}
\end{exercise}
\end{samepage}

\begin{exercise} \label{exercise:convergentinitialvalues}
Misalkan $x_1 \coloneqq c$ dan $x_{n+1} \coloneqq x_n^2+x_n$.
Tunjukkan bahwa $\{ x_n \}_{n=1}^\infty$ konvergen jika dan hanya jika $-1 \leq c \leq 0$; dalam
hal ini barisan tersebut konvergen ke 0.
\end{exercise}

\begin{exercise}
Buktikan $\lim\limits_{n \to \infty} {(n^2+1)}^{1/n} = 1$.
\end{exercise}

\begin{exercise}
Buktikan bahwa $\bigl\{ {(n!)}^{1/n} \bigr\}_{n=1}^\infty$
tak terbatas.
Petunjuk: Tunjukkan bahwa untuk setiap $C > 0$, $\left\{ \frac{C^n}{n!}
\right\}_{n=1}^\infty$ konvergen ke nol.
\end{exercise}

%%%%%%%%%%%%%%%%%%%%%%%%%%%%%%%%%%%%%%%%%%%%%%%%%%%%%%%%%%%%%%%%%%%%%%%%%%%%%%

\sectionnewpage
\section{Limit superior, limit inferior, dan Bolzano--Weierstrass}
\label{sec:bw}

%mbxINTROSUBSECTION

\sectionnotes{1--2 kuliah, pembuktian alternatif BW opsional}

Pada bagian ini, kita mempelajari barisan terbatas dan subbarisannya.
Kita mendefinisikan apa yang disebut limit superior dan limit inferior
dari suatu barisan terbatas, yaitu limit terbesar dan terkecil
yang mungkin dicapai oleh subbarisan-subbarisannya.
Selain itu, kita membuktikan
teorema Bolzano--Weierstrass%
\footnote{%
Nama teorema ini diambil dari matematikawan Ceko
\href{https://en.wikipedia.org/wiki/Bernard_Bolzano}{Bernhard Placidus Johann Nepomuk Bolzano}
(1781--1848), dan matematikawan Jerman
\href{https://en.wikipedia.org/wiki/Karl_Weierstrass}{Karl Theodor Wilhelm Weierstrass}
(1815--1897).}, sebuah
alat yang sangat penting dalam analisis, yang menunjukkan keberadaan
subbarisan konvergen.

Kita telah membuktikan bahwa setiap barisan konvergen bersifat terbatas; meskipun demikian,
ada banyak barisan terbatas yang divergen.  Sebagai contoh,
barisan $\bigl\{ {(-1)}^n \bigr\}_{n=1}^\infty$ bersifat terbatas
tetapi divergen.  Namun, tidak semuanya hilang, dan kita
masih dapat menghitung limit tertentu dari suatu barisan terbatas yang divergen.

\subsection{Limit superior dan limit inferior}

Ada cara untuk membentuk barisan monoton dari sembarang barisan, dan
dengan cara ini kita
memperoleh \emph{\myindex{limit superior}}\index{limsup} dan
\emph{\myindex{limit inferior}}\index{liminf}.  Limit-limit ini selalu ada untuk barisan
terbatas.

Jika suatu barisan $\{ x_n \}_{n=1}^\infty$ terbatas, maka 
himpunan $\{ x_k : k \in \N \}$ terbatas.  Untuk setiap $n$,
himpunan $\{ x_k : k \geq n \}$ juga terbatas (karena merupakan subhimpunan), sehingga kita
mengambil supremum dan infimumnya.

\begin{defn} \label{liminflimsup:def}
Misalkan $\{ x_n \}_{n=1}^\infty$ suatu barisan terbatas.  Definisikan barisan
$\{ a_n \}_{n=1}^\infty$
dan $\{ b_n \}_{n=1}^\infty$ dengan
$a_n \coloneqq \sup \{ x_k : k \geq n \}$ dan
$b_n \coloneqq \inf \{ x_k : k \geq n \}$.  
Jika limitnya ada, definisikan\glsadd{not:limsupseq}\glsadd{not:liminfseq}
\begin{equation*}
\limsup_{n \to \infty} x_n \coloneqq \lim_{n \to \infty} a_n ,
\qquad
\liminf_{n \to \infty} x_n \coloneqq \lim_{n \to \infty} b_n .
\end{equation*}
\end{defn}

Untuk suatu barisan terbatas, liminf dan limsup selalu ada (lihat di bawah).  Kita dapat
mendefinisikan liminf dan limsup untuk barisan tak terbatas jika kita memperbolehkan $\infty$
dan $-\infty$, dan kita akan melakukannya nanti pada bagian ini.
Tidak sulit untuk memperumum hasil-hasil berikut agar
mencakup barisan tak terbatas; namun, terlebih dahulu kita membatasi perhatian pada
barisan terbatas.

\begin{prop}
Misalkan $\{ x_n \}_{n=1}^\infty$ suatu barisan terbatas.  Misalkan $a_n$ dan $b_n$ seperti dalam
definisi di atas.
\begin{enumerate}[(i)]
\item
Barisan
$\{ a_n \}_{n=1}^\infty$ terbatas dan monoton turun
serta $\{ b_n \}_{n=1}^\infty$ terbatas dan monoton naik.  Secara khusus,
$\liminf\limits_{n\to\infty} x_n$ dan $\limsup\limits_{n\to\infty} x_n$ ada.
\item
$\displaystyle \limsup_{n \to \infty} x_n = \inf \{ a_n : n \in \N \}$
dan
$\displaystyle \liminf_{n \to \infty} x_n = \sup \{ b_n : n \in \N \}$.
\item
$\displaystyle \liminf_{n \to \infty} x_n \leq \limsup_{n \to \infty} x_n$.
\end{enumerate}
\end{prop}

\begin{proof}
Mari kita lihat mengapa $\{ a_n \}_{n=1}^\infty$ merupakan barisan turun.  Karena $a_n$ adalah batas
atas terkecil untuk $\{ x_k : k \geq n \}$, bilangan ini juga merupakan
batas atas untuk subhimpunan $\{ x_k : k \geq n+1 \}$.  Oleh karena itu,
$a_{n+1}$, yakni batas atas terkecil untuk
$\{ x_k : k \geq n+1 \}$, harus lebih kecil daripada atau sama dengan $a_n$,
yakni batas atas terkecil untuk
$\{ x_k : k \geq n \}$.
Dengan kata lain,
$a_n \geq a_{n+1}$ untuk semua $n$.  Dengan cara serupa (sebagai latihan),
$\{ b_n \}_{n=1}^\infty$ merupakan barisan naik.
Sebagai latihan, tunjukkan bahwa
jika $\{ x_n \}_{n=1}^\infty$ terbatas, maka $\{ a_n \}_{n=1}^\infty$ dan
$\{ b_n \}_{n=1}^\infty$ juga harus terbatas.

Butir kedua berlaku karena barisan
$\{ a_n \}_{n=1}^\infty$ dan $\{ b_n \}_{n=1}^\infty$ monoton dan terbatas.

Untuk butir ketiga, perhatikan bahwa $b_n \leq a_n$, sebab $\inf$ suatu himpunan tak kosong
lebih kecil daripada atau sama dengan $\sup$-nya.  Barisan $\{ a_n \}_{n=1}^\infty$
dan $\{ b_n \}_{n=1}^\infty$
masing-masing konvergen ke limsup dan liminf.
Terapkan \lemmaref{limandineq:lemma} untuk memperoleh
\begin{equation*}
\lim_{n\to \infty} b_n \leq \lim_{n\to \infty} a_n.  \qedhere
\end{equation*}
\end{proof}
\begin{myfigureht}
\myincludegraphics{sequence-limsupliminf_an_bn}{%
Grafik contoh barisan x indeks n berupa titik-titik
untuk 50 titik, yaitu n dari 1 hingga 50.
Di tengah terdapat dua garis putus-putus horizontal,
garis atas diberi label lim sup x indeks n saat n menuju tak hingga
dan
garis bawah diberi label lim inf x indeks n saat n menuju tak hingga.
Titik-titik berada di atas maupun di bawah garis-garis tersebut, tetapi tampak
semakin berkumpul di dekat daerah antara kedua garis pada 
sisi kanan.  Dua barisan tambahan ditandai,
di bagian bawah dengan wajik, yaitu barisan yang lebih kecil daripada atau sama dengan
semua titik dan bernilai konstan sampai bertemu sebuah titik.  Lalu barisan itu naik
dan kembali konstan sampai bertemu sebuah titik, dan pola ini terus berulang.
Jadi, barisan itu tampak seperti deretan ruas horizontal yang berakhir ketika
wajik berada di tempat yang sama dengan sebuah titik, sedangkan semua titik lainnya
berada di atas wajik-wajik tersebut.
Semua wajik berada di bawah garis putus-putus yang lebih rendah.
Pola serupa tetapi terbalik dibentuk oleh lingkaran dari atas,
yaitu deretan lingkaran sampai bertemu sebuah titik, kemudian turun
dan berulang, dan semua lingkaran berada di atas atau tepat pada titik-titik.
Semua lingkaran berada di atas garis putus-putus yang lebih tinggi.}
\caption{50 suku pertama dari suatu contoh barisan.  Suku-suku $x_n$ dari barisan
ditandai dengan titik
%mbxSTARTIGNORE
(\raisebox{0.25ex}{\tiny$\bullet$}),
%mbxENDIGNORE
%mbxlatex ($\boldsymbol{\cdot}$),
$a_n$ ditandai dengan
lingkaran ($\circ$), dan
$b_n$ ditandai dengan wajik ($\diamond$).\label{sequence-limsupliminf_an_bn}}
\end{myfigureht}

\begin{example} \label{example:liminfsupex}
Misalkan $\{ x_n \}_{n=1}^\infty$ didefinisikan oleh
\begin{equation*}
x_n \coloneqq
\begin{cases}
\frac{n+1}{n} & \text{jika } n \text{ ganjil,} \\
0             & \text{jika } n \text{ genap.}
\end{cases}
\end{equation*}
Mari kita hitung $\liminf$ dan $\limsup$ barisan ini.  Lihat juga
\figureref{sequence-limsupliminf_an_bn-example}.  Pertama, kita hitung
limit inferior:
\begin{equation*}
\liminf_{n\to\infty} x_n = 
\lim_{n\to\infty}
\bigl(
\inf \{ x_k : k \geq n \}
\bigr)
=
\lim_{n\to\infty} 0 = 0 .
\end{equation*}
Untuk limit superior, kita tulis
\begin{equation*}
\limsup_{n\to\infty} x_n = 
\lim_{n\to\infty}
\bigl(
\sup \{ x_k : k \geq n \}
\bigr) .
\end{equation*}
Tidak sulit untuk melihat bahwa
\begin{equation*}
\sup \{ x_k : k \geq n \} =
\begin{cases}
\frac{n+1}{n}   & \text{jika } n \text{ ganjil,} \\
\frac{n+2}{n+1} & \text{jika } n \text{ genap.}
\end{cases}
\end{equation*}
Kita serahkan kepada pembaca untuk menunjukkan bahwa limitnya adalah 1.  Artinya,
\begin{equation*}
\limsup_{n\to\infty} x_n = 1 .
\end{equation*}
Perhatikan bahwa barisan $\{ x_n \}_{n=1}^\infty$ bukan barisan konvergen.
\begin{myfigureht}
\myincludegraphics{sequence-limsupliminf_an_bn-example}{%
Penandaan sama seperti pada gambar sebelumnya, tetapi kini untuk suatu
barisan tertentu.  Garis putus-putus bawah (lim inf) adalah sumbu horizontal
dan garis putus-putus atas kira-kira berada di tengah gambar.
Titik pertama berada jauh di atas garis putus-putus atas
dan sebuah lingkaran berada di
tempat yang sama, sedangkan wajik berada pada sumbu horizontal.
Sesungguhnya, semua wajik
berada pada sumbu horizontal.
Titik kedua berada pada sumbu horizontal, tetapi lingkarannya berada di atas
garis putus-putus atas namun lebih rendah daripada titik pertama.
Titik ketiga berada pada ketinggian yang sama dengan lingkaran kedua
dan lingkaran ketiga berada di tempat yang sama.
Pola ini kemudian terus berulang: titik-titik untuk n ganjil
dan lingkaran-lingkaran semakin mendekati garis putus-putus atas
dari atas.}
\caption{20 suku pertama barisan pada \exampleref{example:liminfsupex}.
Penandaannya seperti pada \figureref{sequence-limsupliminf_an_bn}.
\label{sequence-limsupliminf_an_bn-example}}
\end{myfigureht}
\end{example}

Kita mengaitkan subbarisan tertentu dengan $\limsup$ dan $\liminf$.
Penting untuk diperhatikan bahwa $\{ a_n \}_{n=1}^\infty$ dan
$\{ b_n \}_{n=1}^\infty$ bukan
subbarisan dari $\{ x_n \}_{n=1}^\infty$, dan keduanya juga tidak harus
terdiri atas bilangan-bilangan yang sama.
Sebagai contoh, jika barisannya adalah $\{ \nicefrac{1}{n} \}_{n=1}^\infty$, maka
$b_n = 0$ untuk semua $n \in \N$.

\begin{thm} \label{subseqlimsupinf:thm}
Jika $\{ x_n \}_{n=1}^\infty$ suatu barisan terbatas, maka ada suatu subbarisan
$\{ x_{n_k} \}_{k=1}^\infty$ sedemikian sehingga
\begin{equation*}
\lim_{k\to \infty} x_{n_k} = \limsup_{n \to \infty} x_n .
\end{equation*}
Dengan cara serupa, ada suatu subbarisan (yang mungkin berbeda)
$\{ x_{m_k} \}_{k=1}^\infty$ sedemikian sehingga
\begin{equation*}
\lim_{k\to \infty} x_{m_k} = \liminf_{n \to \infty} x_n .
\end{equation*}
\end{thm}

\begin{proof}
Definisikan $a_n \coloneqq \sup \{ x_k : k \geq n \}$.
Tuliskan
$x \coloneqq \limsup_{n\to\infty} x_n = \lim_{n\to\infty} a_n$.
Kita mendefinisikan subbarisan tersebut secara induktif.
Misalkan $n_1 \coloneqq 1$, dan andaikan $n_1,n_2,\ldots,n_{k-1}$ sudah didefinisikan untuk
suatu $k \geq 2$.  Pilih $m \geq n_{k-1} + 1$
sedemikian sehingga
\begin{equation*}
a_{(n_{k-1}+1)} - x_m < \frac{1}{k} .
\end{equation*}
$m$ seperti itu ada
karena $a_{(n_{k-1}+1)}$ adalah supremum dari
himpunan $\{ x_\ell : \ell \geq n_{k-1} + 1 \}$ sehingga terdapat suku-suku
barisan yang sedekat apa pun dengan supremum (atau bahkan mungkin sama dengannya).
Tetapkan $n_{k} \coloneqq  m$.  Subbarisan $\{ x_{n_k} \}_{k=1}^\infty$ kini terdefinisi.  Selanjutnya, kita
harus membuktikan bahwa subbarisan itu konvergen ke $x$.

Untuk semua $k \geq 2$, kita mempunyai
$a_{(n_{k-1}+1)} \geq a_{n_k}$ (mengapa?) dan $a_{n_{k}} \geq x_{n_k}$.
Oleh karena itu, untuk setiap $k \geq 2$,
\begin{equation*}
\begin{split}
\sabs{a_{n_k} - x_{n_k}} & = 
a_{n_k} - x_{n_k}
\\
& \leq
a_{(n_{k-1}+1)} - x_{n_k}
\\
& < \frac{1}{k} .
\end{split}
\end{equation*}

Mari kita tunjukkan bahwa $\{ x_{n_k} \}_{k=1}^\infty$ konvergen ke $x$.
Perhatikan bahwa subbarisan tersebut tidak harus monoton.  Ambil $\epsilon > 0$.
Karena $\{ a_n \}_{n=1}^\infty$ konvergen ke $x$, subbarisan
$\{ a_{n_k} \}_{k=1}^\infty$ konvergen ke $x$.
Jadi, terdapat $M_1 \in \N$
sedemikian sehingga untuk semua $k \geq M_1$, berlaku
\begin{equation*}
\sabs{a_{n_k} - x} < \frac{\epsilon}{2} .
\end{equation*}
Pilih $M_2 \in \N$ sedemikian sehingga
\begin{equation*}
\frac{1}{M_2} \leq \frac{\epsilon}{2}.
\end{equation*}
Ambil $M \coloneqq \max \{M_1 , M_2 , 2 \}$.  Untuk semua $k \geq M$,
\begin{equation*}
\begin{split}
\sabs{x- x_{n_k}} & =
\sabs{a_{n_k} - x_{n_k} + x - a_{n_k}}
\\
& \leq \sabs{a_{n_k} - x_{n_k}} + \sabs{x - a_{n_k}}
\\
& < \frac{1}{k} + \frac{\epsilon}{2}
\\
& \leq \frac{1}{M_2} + \frac{\epsilon}{2} \leq \frac{\epsilon}{2} +
\frac{\epsilon}{2} = \epsilon .
\end{split}
\end{equation*}

Pernyataan untuk $\liminf$ diserahkan sebagai latihan.
\end{proof}

\subsection{Menggunakan limit inferior dan limit superior}

Keunggulan $\liminf$ dan $\limsup$ adalah bahwa keduanya selalu dapat dituliskan
untuk sebarang barisan (terbatas).
Jika kita dapat menghitungnya dengan suatu cara, kita juga dapat menghitung limit
barisan itu jika ada, atau menunjukkan bahwa barisan itu divergen.
Perhitungan dengan $\liminf$ dan $\limsup$ agak
menyerupai perhitungan dengan limit, meskipun terdapat perbedaan-perbedaan halus.

\begin{prop} \label{liminfsupconv:prop}
Misalkan $\{ x_n \}_{n=1}^\infty$ adalah barisan terbatas.  Maka
$\{ x_n \}_{n=1}^\infty$ konvergen
jika dan hanya jika
\begin{equation*}
\liminf_{n\to \infty} x_n = 
\limsup_{n\to \infty} x_n.
\end{equation*}
Lebih lanjut, jika $\{ x_n \}_{n=1}^\infty$ konvergen, maka
\begin{equation*}
\lim_{n\to \infty} x_n = 
\liminf_{n\to \infty} x_n = 
\limsup_{n\to \infty} x_n.
\end{equation*}
\end{prop}

\begin{proof}
Misalkan $a_n$ dan $b_n$ seperti dalam \defnref{liminflimsup:def}.
Secara khusus, untuk semua $n \in \N$,
\begin{equation*}
b_n \leq x_n \leq a_n .
\end{equation*}
Pertama, misalkan
$\liminf_{n\to\infty} x_n = \limsup_{n\to\infty} x_n$.
Maka $\{ a_n \}_{n=1}^\infty$ dan $\{ b_n \}_{n=1}^\infty$
keduanya konvergen ke limit yang sama.
Berdasarkan lemma apit
(\lemmaref{squeeze:lemma}), $\{ x_n \}_{n=1}^\infty$ konvergen dan
\begin{equation*}
\lim_{n\to \infty} b_n
=
\lim_{n\to \infty} x_n
=
\lim_{n\to \infty} a_n .
\end{equation*}

Sekarang, misalkan $\{ x_n \}_{n=1}^\infty$ konvergen ke $x$.
Berdasarkan \thmref{subseqlimsupinf:thm},
terdapat subbarisan $\{ x_{n_k} \}_{k=1}^\infty$
yang konvergen ke $\limsup_{n\to\infty} x_n$.
Karena $\{ x_n \}_{n=1}^\infty$ konvergen ke $x$,
setiap subbarisan konvergen ke $x$ dan
karena itu $\limsup_{n\to\infty} x_n = \lim_{k\to\infty} x_{n_k} = x$.
Demikian pula, $\liminf_{n\to\infty} x_n = x$.
\end{proof}

Limit superior dan limit inferior berperilaku baik
terhadap subbarisan.

\begin{prop} \label{prop:subseqslimsupinf}
Misalkan $\{ x_n \}_{n=1}^\infty$ adalah barisan terbatas dan
$\{ x_{n_k} \}_{k=1}^\infty$ adalah subbarisan.  Maka
\begin{equation*}
\liminf_{n\to\infty} x_n \leq
\liminf_{k\to\infty} x_{n_k} \leq
\limsup_{k\to\infty} x_{n_k} \leq
\limsup_{n\to\infty} x_n .
\end{equation*}
\end{prop}

\begin{proof}
Ketaksamaan di tengah telah dibuktikan.  Kita akan membuktikan ketaksamaan
ketiga dan menyerahkan ketaksamaan pertama sebagai latihan.

Kita ingin membuktikan bahwa
$\limsup_{k\to\infty} x_{n_k} \leq \limsup_{n\to\infty} x_n$.  Definisikan
$a_n \coloneqq \sup \{ x_k : k \geq n \}$ 
seperti biasa.
Definisikan juga
$c_n \coloneqq \sup \{ x_{n_k} : k \geq n \}$.
Tidak selalu benar bahwa $\{ c_n \}_{n=1}^\infty$ merupakan subbarisan
dari $\{ a_n \}_{n=1}^\infty$.
Namun, karena $n_k \geq k$ untuk semua $k$, kita memperoleh
$\{ x_{n_k} : k \geq n \} \subset \{ x_k : k \geq n \}$.
Supremum suatu himpunan bagian kurang dari atau sama dengan supremum
himpunan tersebut; oleh karena itu,
\begin{equation*}
c_n \leq a_n \qquad \text{untuk semua $n$}.
\end{equation*}
\lemmaref{limandineq:lemma} memberikan
\begin{equation*}
\lim_{n\to\infty} c_n \leq \lim_{n\to\infty} a_n ,
\end{equation*}
yang merupakan kesimpulan yang diinginkan.
\end{proof}

Limit superior dan limit inferior
masing-masing merupakan limit subbarisan terbesar dan terkecil.
Jika subbarisan $\{ x_{n_k} \}_{k=1}^\infty$ dalam proposisi sebelumnya
konvergen, maka
$\liminf_{k\to\infty} x_{n_k} = \lim_{k\to\infty} x_{n_k} =
\limsup_{k\to\infty} x_{n_k}$.  Oleh karena itu,
\begin{equation*}
\liminf_{n\to\infty} x_n \leq
\lim_{k\to\infty} x_{n_k} \leq
\limsup_{n\to\infty} x_n .
\end{equation*}

Demikian pula, kita memperoleh uji berguna berikut untuk kekonvergenan
barisan terbatas.  Buktinya kita serahkan sebagai latihan.

\begin{prop} \label{seqconvsubseqconv:prop}
Barisan terbatas $\{ x_n \}_{n=1}^\infty$ konvergen ke $x$
jika dan hanya jika
setiap subbarisan konvergen
$\{ x_{n_k} \}_{k=1}^\infty$ juga konvergen ke $x$.
\end{prop}

\subsection{Teorema Bolzano--Weierstrass}

Meskipun tidak setiap barisan terbatas konvergen, teorema
Bolzano--Weierstrass memberi tahu kita bahwa setidaknya kita dapat menemukan
sebuah subbarisan yang konvergen.
Versi Bolzano--Weierstrass yang kami sajikan dalam bagian ini adalah
teorema Bolzano--Weierstrass untuk barisan bilangan real.

\begin{thm}[Bolzano--Weierstrass]\index{teorema Bolzano--Weierstrass}\label{thm:bwseq}
Andaikan suatu barisan bilangan real $\{ x_n \}_{n=1}^\infty$ terbatas.
Maka terdapat subbarisan konvergen $\{ x_{n_i} \}_{i=1}^\infty$.
\end{thm}

\begin{proof}
\thmref{subseqlimsupinf:thm} menyatakan bahwa terdapat
sebuah subbarisan yang limitnya adalah $\limsup_{n\to\infty} x_n$.
\end{proof}

Pembaca mungkin menyanggah bahwa
\thmref{subseqlimsupinf:thm} secara ketat lebih kuat daripada teorema
Bolzano--Weierstrass yang disajikan di atas.  Hal itu benar.
Namun,
\thmref{subseqlimsupinf:thm} hanya berlaku pada garis bilangan real, sedangkan
Bolzano--Weierstrass berlaku dalam konteks yang lebih umum (yaitu, dalam $\R^n$)
dengan pernyataan yang nyaris persis sama.

Karena teorema ini begitu penting dalam analisis, kami menyajikan sebuah
bukti eksplisit.
Gagasan dalam bukti berikut juga dapat diperumum ke konteks-konteks lain.

\begin{proof}[Bukti alternatif teorema Bolzano--Weierstrass]
Karena barisan tersebut terbatas, terdapat dua bilangan $a_1 < b_1$
sedemikian sehingga $a_1 \leq x_n \leq b_1$ untuk semua $n \in \N$.
Kita akan mendefinisikan sebuah subbarisan $\{ x_{n_i} \}_{i=1}^\infty$ dan dua
barisan $\{ a_i \}_{i=1}^\infty$ dan $\{ b_i \}_{i=1}^\infty$, sedemikian sehingga
$\{ a_i \}_{i=1}^\infty$ monoton naik, $\{ b_i \}_{i=1}^\infty$ monoton turun,
$a_i \leq x_{n_i} \leq b_i$, dan $\lim_{i\to\infty} a_i =
\lim_{i\to\infty} b_i$.   Dengan demikian,
konvergensi $\{ x_{n_i} \}_{i=1}^\infty$ diperoleh dari \hyperref[squeeze:lemma]{lemma apit}.

Kita mendefinisikan barisan-barisan tersebut secara induktif.  Kita akan mendefinisikannya
sedemikian sehingga, untuk semua $i$, berlaku $a_i < b_i$,
dan $x_n \in [a_i,b_i]$ untuk tak berhingga banyak $n \in \N$.
Kita telah mendefinisikan $a_1$ dan $b_1$.  Ambil $n_1 \coloneqq 1$, yaitu
$x_{n_1} = x_1$.
Andaikan bahwa sampai suatu $k \in \N$,
kita telah mendefinisikan subbarisan $x_{n_1}, x_{n_2}, \ldots,
x_{n_k}$, serta barisan $a_1,a_2,\ldots,a_k$
dan $b_1,b_2,\ldots,b_k$.
Misalkan $y \coloneqq \frac{a_k+b_k}{2}$.
Jelas bahwa
$a_k < y < b_k$.  Jika terdapat tak berhingga banyak $j \in \N$
sedemikian sehingga $x_j \in [a_k,y]$, tetapkan $a_{k+1} \coloneqq a_k$, $b_{k+1}
\coloneqq y$,
dan pilih $n_{k+1} > n_{k}$
sedemikian sehingga $x_{n_{k+1}} \in [a_k,y]$.  Jika hanya terdapat berhingga banyak
$j$ yang memenuhi
$x_j \in [a_k,y]$, maka harus terdapat tak berhingga banyak $j \in
\N$ sedemikian sehingga
$x_j \in [y,b_k]$.  Dalam kasus ini, tetapkan $a_{k+1} \coloneqq y$, $b_{k+1}
\coloneqq b_k$,
dan pilih $n_{k+1} > n_{k}$
sedemikian sehingga $x_{n_{k+1}} \in [y,b_k]$.

Sekarang barisan-barisan tersebut telah terdefinisi.  Tinggal dibuktikan bahwa
$\lim_{i\to\infty} a_i = \lim_{i\to\infty} b_i$.  Kedua limit tersebut ada karena barisan-barisan
tersebut monoton.  Dalam konstruksi ini,
$b_i - a_i$ dibagi dua pada setiap langkah.  Oleh karena itu,
$b_{i+1} - a_{i+1} = \frac{b_i-a_i}{2}$.  Berdasarkan
\hyperref[induction:thm]{induksi},
\begin{equation*}
b_i - a_i = \frac{b_1-a_1}{2^{i-1}} .
\end{equation*}

Misalkan $x \coloneqq \lim_{i\to\infty} a_i$.  Karena $\{ a_i \}_{i=1}^\infty$ monoton,
\begin{equation*}
x = \sup \{ a_i : i \in \N \} .
\end{equation*}
Misalkan $y \coloneqq \lim_{i\to\infty} b_i = \inf \{ b_i : i \in \N \}$.  Karena $a_i < b_i$ untuk
semua $i$, kita memperoleh $x \leq y$.
Karena barisan-barisan tersebut monoton, untuk semua $i$ berlaku (mengapa?)
\begin{equation*}
y-x \leq b_i-a_i = \frac{b_1-a_1}{2^{i-1}} .
\end{equation*}
Karena $\frac{b_1-a_1}{2^{i-1}}$ dapat dibuat sekecil apa pun dan $y-x \geq 0$,
kita memperoleh $y-x = 0$.  Selesaikan dengan \hyperref[squeeze:lemma]{lemma apit}.
\end{proof}

Bukti lain lagi untuk teorema Bolzano--Weierstrass diperoleh dengan menunjukkan
klaim berikut,
yang diserahkan sebagai latihan yang menantang.
\emph{Klaim: Setiap barisan memiliki subbarisan monoton}.

\subsection{Limit tak hingga}

Sama seperti pada infimum dan supremum, kita dapat memperbolehkan limit tertentu
bernilai tak hingga.  Dengan kata lain, kita menulis $\lim_{n\to\infty} x_n = \infty$ atau
$\lim_{n\to\infty} x_n = -\infty$ untuk barisan divergen tertentu.

\begin{defn}
\index{limit tak hingga!dari barisan}\index{limit!tak hingga}
Kita mengatakan
$\{ x_n \}_{n=1}^\infty$ \emph{\myindex{divergen menuju tak hingga}}%
\footnote{Kadang-kadang dikatakan bahwa $\{ x_n \}_{n=1}^\infty$ \emph{konvergen menuju tak hingga}.}
jika untuk setiap $K \in
\R$, terdapat $M \in \N$ sedemikian sehingga untuk semua $n \geq M$, berlaku $x_n >
K$.  Dalam hal ini kita menulis
\begin{equation*}
\lim_{n \to \infty} x_n \coloneqq \infty .  
\end{equation*}
Demikian pula,
jika untuk setiap $K \in \R$ terdapat $M \in \N$ sedemikian sehingga
untuk semua $n \geq M$, berlaku $x_n < K$, kita mengatakan $\{ x_n \}_{n=1}^\infty$
\emph{\myindex{divergen menuju minus tak hingga}} dan kita menulis
\begin{equation*}
\lim_{n \to \infty} x_n \coloneqq -\infty .  
\end{equation*}
\end{defn}

Dengan definisi ini dan dengan memperbolehkan $\infty$ dan $-\infty$,
kita dapat menulis $\lim_{n\to\infty} x_n$ untuk setiap barisan monoton.

\begin{prop} \label{prop:unboundedmonotone}
Misalkan $\{ x_n \}_{n=1}^\infty$ suatu barisan monoton tak terbatas.  Maka
\begin{equation*}
\lim_{n \to \infty} x_n =
\begin{cases}
\infty  & \text{jika } \{ x_n \}_{n=1}^\infty \text{ monoton naik,} \\
-\infty & \text{jika } \{ x_n \}_{n=1}^\infty \text{ monoton turun.}
\end{cases}
\end{equation*}
\end{prop}

\begin{proof}
Kasus monoton naik merupakan akibat dari
\exerciseref{exercise:infseqlimlims} bagian c) di bawah.
Misalkan $\{x_n\}_{n=1}^\infty$ monoton turun dan tak terbatas.
Bahwa barisan ini tak terbatas berarti
untuk setiap $K \in \R$, terdapat $M \in \N$ sedemikian sehingga $x_M < K$.
Karena barisan tersebut monoton, $x_n \leq x_M < K$ untuk semua $n \geq M$.
Oleh karena itu, $\lim_{n\to\infty} x_n = -\infty$.
\end{proof}

\begin{example}
\begin{equation*}
\lim_{n\to \infty} n = \infty,
\qquad
\lim_{n\to \infty} n^2 = \infty,
\qquad
\lim_{n\to \infty} -n = -\infty.
\end{equation*}
Verifikasinya diserahkan kepada pembaca.
\end{example}

Kita juga dapat memperbolehkan $\liminf$ dan $\limsup$ mengambil
nilai $\infty$ dan $-\infty$, sehingga
kita dapat menerapkan $\liminf$ dan $\limsup$
pada sembarang barisan, bukan hanya
barisan terbatas.  Sayangnya, barisan $\{ a_n \}_{n=1}^\infty$ dan
$\{ b_n \}_{n=1}^\infty$
bukanlah barisan bilangan real, melainkan barisan bilangan real diperluas.  Secara
khusus, $a_n$ dapat sama dengan $\infty$ untuk suatu $n$, dan $b_n$ dapat sama dengan
$-\infty$.  Jadi kita tidak memiliki definisi untuk limit-limitnya.
Namun, karena bilangan real diperluas masih merupakan himpunan terurut, kita
dapat mengambil supremum dan infimumnya.

\begin{defn}
Misalkan $\{ x_n \}_{n=1}^\infty$ suatu barisan bilangan real yang tak terbatas.  Definisikan
barisan-barisan bilangan real diperluas dengan
$a_n \coloneqq \sup \{ x_k : k \geq n \}$ dan
$b_n \coloneqq \inf \{ x_k : k \geq n \}$.
Definisikan\glsadd{not:limsupseq}\glsadd{not:liminfseq}
\index{limit inferior}\index{limit superior}\index{limsup}\index{liminf}
\begin{equation*}
\limsup_{n \to \infty} x_n \coloneqq \inf \{ a_n : n \in \N \}, \qquad \text{dan} \qquad
\liminf_{n \to \infty} x_n \coloneqq \sup \{ b_n : n \in \N \}.
\end{equation*}
\end{defn}

Definisi ini sesuai dengan definisi untuk barisan
terbatas.

\begin{prop}
Misalkan $\{ x_n \}_{n=1}^\infty$ suatu barisan tak terbatas.  Definisikan
$\{ a_n \}_{n=1}^\infty$ dan $\{ b_n \}_{n=1}^\infty$ seperti di atas.
Maka $\{ a_n \}_{n=1}^\infty$ monoton turun, dan $\{ b_n \}_{n=1}^\infty$ monoton naik.
Jika $a_n$ merupakan bilangan real untuk setiap $n$, maka
$\limsup_{n\to\infty} x_n = \lim_{n\to\infty} a_n$.
Jika $b_n$ merupakan bilangan real untuk setiap $n$, maka
$\liminf_{n\to\infty} x_n = \lim_{n\to\infty} b_n$.
\end{prop}

\begin{proof}
Seperti sebelumnya,
$a_n = \sup \{ x_k : k \geq n \} \geq \sup \{ x_k : k \geq n+1 \} =
a_{n+1}$.  Jadi $\{ a_n \}_{n=1}^\infty$ monoton turun. Demikian pula,
$\{ b_n \}_{n=1}^\infty$ monoton naik.

Jika barisan $\{ a_n \}_{n=1}^\infty$ merupakan barisan bilangan real, maka
$\lim_{n\to\infty} a_n = \inf \{ a_n : n \in \N \}$.  Hal ini mengikuti
\thmref{thm:monotoneconv} jika $\{ a_n \}_{n=1}^\infty$ terbatas dan
\propref{prop:unboundedmonotone} jika $\{a_n \}_{n=1}^\infty$
tak terbatas.  Kita melanjutkan dengan cara serupa untuk $\{ b_n \}_{n=1}^\infty$.
\end{proof}

Definisi ini berperilaku sebagaimana yang diharapkan terhadap
$\limsup$ dan $\liminf$; lihat latihan \ref{exercise:infseqlimex}
dan \ref{exercise:infseqlimlims}.

\begin{example}
Misalkan
$x_n \coloneqq 0$ untuk $n$ ganjil dan $x_n \coloneqq n$ untuk $n$ genap.
Maka $a_n = \infty$ untuk semua $n$, karena untuk setiap $M$,
terdapat $k$ genap sedemikian sehingga $x_k = k \geq M$.
Di sisi lain, $b_n = 0$ untuk semua $n$, karena
untuk setiap $n$, himpunan
$\{ x_k : k \geq n \}$ terdiri atas $0$ dan bilangan-bilangan positif.
Jadi,
\begin{equation*}
\lim_{n\to \infty} x_n \quad \text{tidak ada},
\qquad
\limsup_{n\to \infty} x_n = \infty ,
\qquad
\liminf_{n\to \infty} x_n = 0.
\end{equation*}
\end{example}

\subsection{Latihan}

\begin{exercise}
Misalkan $\{ x_n \}_{n=1}^\infty$ adalah barisan terbatas.  Definisikan $a_n$ dan
$b_n$ seperti dalam \defnref{liminflimsup:def}.  Tunjukkan bahwa $\{ a_n \}_{n=1}^\infty$
dan $\{ b_n \}_{n=1}^\infty$ terbatas.
\end{exercise}

\begin{exercise}
Misalkan $\{ x_n \}_{n=1}^\infty$ adalah barisan terbatas.
Definisikan $b_n$ seperti dalam \defnref{liminflimsup:def}.  Tunjukkan bahwa
$\{ b_n \}_{n=1}^\infty$ adalah barisan monoton naik.
\end{exercise}

\begin{exercise}
Selesaikan bukti \propref{prop:subseqslimsupinf}.  Dengan kata lain,
misalkan $\{ x_n \}_{n=1}^\infty$ adalah barisan terbatas dan
$\{ x_{n_k} \}_{k=1}^\infty$ adalah subbarisan.  Buktikan
$\displaystyle \liminf_{n\to\infty} x_n \leq
\liminf_{k\to\infty} x_{n_k}$.
\end{exercise}

\begin{exercise}
Buktikan \propref{seqconvsubseqconv:prop}.
\end{exercise}

\begin{exercise}
\leavevmode
\begin{enumerate}[a)]
\item
Misalkan $x_n \coloneqq \dfrac{{(-1)}^n}{n}$.
Tentukan $\limsup\limits_{n\to\infty} x_n$ dan $\liminf\limits_{n\to\infty} x_n$.
\item
Misalkan $x_n \coloneqq \dfrac{(n-1){(-1)}^n}{n}$.
Tentukan $\limsup\limits_{n\to\infty} x_n$ dan $\liminf\limits_{n\to\infty} x_n$.
\end{enumerate}
\end{exercise}

\begin{exercise}
Misalkan $\{ x_n \}_{n=1}^\infty$ dan $\{ y_n \}_{n=1}^\infty$ adalah barisan terbatas sedemikian sehingga
$x_n \leq y_n$ untuk semua $n$.  Tunjukkan
\begin{equation*}
\limsup_{n\to\infty} x_n \leq
\limsup_{n\to\infty} y_n
\qquad \text{dan} \qquad
\liminf_{n\to\infty} x_n \leq
\liminf_{n\to\infty} y_n .
\end{equation*}
\end{exercise}

\begin{exercise}
Misalkan $\{ x_n \}_{n=1}^\infty$ dan $\{ y_n \}_{n=1}^\infty$ adalah barisan terbatas.
\begin{enumerate}[a)]
\item
Tunjukkan bahwa $\{ x_n + y_n \}_{n=1}^\infty$ terbatas.
\item
Tunjukkan bahwa
\begin{equation*}
\Bigl(\liminf_{n\to \infty} x_n\Bigr)
+
\Bigl(\liminf_{n\to \infty} y_n\Bigr)
\leq
\liminf_{n\to \infty} \, (x_n+y_n) .
\end{equation*}
Petunjuk: Salah satu cara pembuktian adalah mencari subbarisan $\{ x_{n_m}+y_{n_m} \}_{m=1}^\infty$ dari
$\{ x_n + y_n \}_{n=1}^\infty$ yang konvergen.
Kemudian cari subbarisan $\{ x_{n_{m_i}} \}_{i=1}^\infty$ dari
$\{ x_{n_m} \}_{m=1}^\infty$ yang konvergen.
\item
Temukan $\{ x_n \}_{n=1}^\infty$ dan $\{ y_n \}_{n=1}^\infty$ secara eksplisit sedemikian sehingga
\begin{equation*}
\Bigl(\liminf_{n\to \infty} x_n\Bigr)
+
\Bigl(\liminf_{n\to \infty} y_n\Bigr)
<
\liminf_{n\to \infty} \, (x_n+y_n) .
\end{equation*}
Petunjuk: Carilah contoh yang tidak memiliki limit.
\end{enumerate}
\end{exercise}

\begin{samepage}
\begin{exercise}
Misalkan $\{ x_n \}_{n=1}^\infty$ dan $\{ y_n \}_{n=1}^\infty$ adalah barisan terbatas (menurut
latihan sebelumnya, $\{ x_n + y_n \}_{n=1}^\infty$ terbatas).
\begin{enumerate}[a)]
\item
Tunjukkan bahwa
\begin{equation*}
\Bigl(\limsup_{n\to \infty} x_n\Bigr)
+
\Bigl(\limsup_{n\to \infty} y_n\Bigr)
\geq
\limsup_{n\to \infty} \, (x_n+y_n) .
\end{equation*}
Petunjuk: Lihat latihan sebelumnya.
\item
Temukan $\{ x_n \}_{n=1}^\infty$ dan $\{ y_n \}_{n=1}^\infty$ secara eksplisit sedemikian sehingga
\begin{equation*}
\Bigl(\limsup_{n\to \infty} x_n\Bigr)
+
\Bigl(\limsup_{n\to \infty} y_n\Bigr)
>
\limsup_{n\to \infty} \, (x_n+y_n) .
\end{equation*}
Petunjuk: Lihat latihan sebelumnya.
\end{enumerate}
\end{exercise}
\end{samepage}

\begin{exercise}
Jika $S \subset \R$ adalah suatu himpunan, maka $x \in \R$ disebut \emph{\myindex{titik
akumulasi}}
jika untuk setiap $\epsilon > 0$, himpunan $(x-\epsilon,x+\epsilon) \cap S
\setminus \{ x \}$ tidak kosong.  Dengan kata lain, terdapat titik-titik dalam $S$
yang berjarak sedekat apa pun dengan $x$.
Sebagai contoh, $S \coloneqq \{ \nicefrac{1}{n} : n \in \N \}$ memiliki titik akumulasi tunggal $0$, tetapi $0 \notin S$.
Buktikan versi teorema Bolzano--Weierstrass berikut:

\medskip

\noindent
\emph{\textbf{Teorema.} Misalkan $S \subset \R$ adalah himpunan tak berhingga yang terbatas,
maka terdapat setidaknya satu titik akumulasi dari $S$}.

\medskip

Petunjuk: Jika $S$ tak berhingga, maka $S$ memuat suatu himpunan bagian tak berhingga yang terhitung.
Dengan kata lain, terdapat barisan $\{ x_n \}_{n=1}^\infty$ yang terdiri atas bilangan-bilangan berbeda dalam $S$.
\end{exercise}

\begin{samepage}
\begin{exercise}[Menantang]
\leavevmode
\begin{enumerate}[a)]
\item
Buktikan bahwa setiap barisan memuat suatu subbarisan monoton.
Petunjuk: Sebut $n \in \N$ sebagai suatu \emph{puncak}
dari barisan
$\{ x_n \}_{n=1}^\infty$
jika $x_m \leq x_n$ untuk semua $m \geq n$.
Ada dua kemungkinan: Barisan tersebut hanya memiliki berhingga banyak
puncak,
atau memiliki tak berhingga banyak puncak.
\item
Simpulkan teorema Bolzano--Weierstrass.
\end{enumerate}
\end{exercise}
\end{samepage}

\begin{exercise}
\pagebreak[2]
Buktikan versi yang lebih kuat dari \propref{seqconvsubseqconv:prop}.
Misalkan $\{ x_n \}_{n=1}^\infty$ adalah barisan sedemikian sehingga setiap subbarisan
$\{ x_{n_m} \}_{m=1}^\infty$ memiliki subbarisan
$\{ x_{n_{m_i}} \}_{i=1}^\infty$ yang konvergen ke $x$.
\begin{enumerate}[a)]
\item
Pertama, tunjukkan bahwa $\{ x_n \}_{n=1}^\infty$
terbatas.
\item
Sekarang tunjukkan bahwa $\{ x_n \}_{n=1}^\infty$ konvergen ke $x$.
\end{enumerate}
\end{exercise}

\begin{exercise}
Misalkan $\{x_n\}_{n=1}^\infty$ adalah barisan terbatas.
\begin{enumerate}[a)]
\item
Buktikan bahwa terdapat $s$ sedemikian sehingga untuk setiap $r > s$, terdapat 
$M \in \N$ sedemikian sehingga untuk semua $n \geq M$, berlaku
$x_n < r$.
\item
Jika $s$ adalah bilangan seperti pada a), buktikan bahwa $\limsup\limits_{n\to\infty} x_n \leq s$.
\item
Tunjukkan bahwa jika $S$ adalah himpunan semua $s$ seperti pada a), maka
$\limsup\limits_{n\to\infty} x_n = \inf \, S$.
\end{enumerate}
\end{exercise}

\begin{exercise}[Mudah] \label{exercise:infseqlimex}
\pagebreak[2]
Misalkan $\{ x_n \}_{n=1}^\infty$ sedemikian sehingga $\liminf\limits_{n\to\infty} x_n = -\infty$,
$\limsup\limits_{n\to\infty} x_n = \infty$.
\begin{enumerate}[a)]
\item
Tunjukkan bahwa $\{ x_n \}_{n=1}^\infty$ tidak konvergen, dan juga
bahwa baik $\lim\limits_{n\to\infty} x_n = \infty$
maupun $\lim\limits_{n\to\infty} x_n = -\infty$
tidak berlaku.
\item
Temukan contoh barisan semacam itu.
\end{enumerate}
\end{exercise}

\begin{exercise} \label{exercise:infseqlimlims}
Misalkan $\{ x_n \}_{n=1}^\infty$ adalah suatu barisan.
\begin{enumerate}[a)]
\item
Tunjukkan bahwa
$\lim\limits_{n\to\infty} x_n = \infty$ jika dan hanya jika
$\liminf\limits_{n\to\infty} x_n = \infty$.
\item
Kemudian tunjukkan bahwa $\lim\limits_{n\to\infty} x_n = - \infty$ jika dan hanya jika
$\limsup\limits_{n\to\infty} x_n = -\infty$.
\item
Jika $\{ x_n \}_{n=1}^\infty$ monoton naik, tunjukkan bahwa
$\lim_{n\to\infty} x_n$ ada dan berhingga atau $\lim_{n\to\infty} x_n = \infty$.  Dalam kedua
kasus, $\lim_{n\to\infty} x_n = \sup \{ x_n : n \in \N \}$.
\end{enumerate}
\end{exercise}

\begin{exercise} \label{exercise:strongerratiotest2}
Buktikan versi yang lebih kuat berikut dari \lemmaref{seq:ratiotest}, yaitu uji
rasio.
Misalkan $\{ x_n \}_{n=1}^\infty$ adalah barisan sedemikian sehingga $x_n \neq 0$ untuk semua
$n$.
\begin{enumerate}[a)]
\item
Buktikan bahwa jika
\begin{equation*}
\limsup_{n \to \infty} \frac{\sabs{x_{n+1}}}{\sabs{x_n}} < 1 ,
\end{equation*}
maka $\{ x_n \}_{n=1}^\infty$ konvergen ke $0$.
\item
Buktikan bahwa jika
\begin{equation*}
\liminf_{n \to \infty} \frac{\sabs{x_{n+1}}}{\sabs{x_n}} > 1 ,
\end{equation*}
maka $\{ x_n \}_{n=1}^\infty$ tak terbatas.
\end{enumerate}
\end{exercise}

\begin{exercise}
Misalkan $\{ x_n \}_{n=1}^\infty$ adalah barisan terbatas dan, seperti sebelumnya, $a_n \coloneqq \sup \{ x_k : k \geq n \}$.
Andaikan bahwa untuk suatu $\ell \in \N$,
$a_\ell \notin \{ x_k : k \geq \ell \}$.  Tunjukkan bahwa $a_j = a_\ell$ untuk semua $j \geq \ell$, sehingga
$\limsup\limits_{n\to\infty} x_n = a_\ell$.
\end{exercise}

\begin{exercise}
Misalkan $\{ x_n \}_{n=1}^\infty$ adalah suatu barisan,
dan, seperti sebelumnya, $a_n \coloneqq \sup \{ x_k : k \geq n \}$ serta
$b_n \coloneqq \inf \{ x_k : k \geq n \}$.
\begin{enumerate}[a)]
\item
Buktikan bahwa jika $a_\ell = \infty$ untuk suatu $\ell \in \N$, maka 
$\limsup\limits_{n\to\infty} x_n = \infty$.
\item
Buktikan bahwa jika $b_\ell = -\infty$ untuk suatu $\ell \in \N$, maka 
$\liminf\limits_{n\to\infty} x_n = -\infty$.
\end{enumerate}
\end{exercise}

\begin{exercise}
Misalkan $\{ x_n \}_{n=1}^\infty$ adalah barisan
sedemikian sehingga $\liminf_{n\to\infty} x_n$ dan
$\limsup_{n\to\infty} x_n$ keduanya berhingga.  Buktikan bahwa $\{ x_n \}_{n=1}^\infty$ terbatas.
\end{exercise}

\begin{exercise}
Misalkan $\{ x_n \}_{n=1}^\infty$ adalah barisan terbatas dan diberikan $\epsilon > 0$.
Buktikan bahwa terdapat $M$ sedemikian sehingga untuk semua $k \geq M$,
\begin{equation*}
x_k - \Bigl(\limsup_{n\to\infty} x_n\Bigr) < \epsilon \qquad \text{dan} \qquad
\Bigl(\liminf_{n\to\infty} x_n\Bigr) - x_k < \epsilon .
\end{equation*}
\end{exercise}

\begin{exercise} \label{exercise:extendsubseqlimsupinf}
Perluas \thmref{subseqlimsupinf:thm} ke barisan tak terbatas:
Misalkan $\{ x_n \}_{n=1}^\infty$ adalah suatu barisan.
Jika $\limsup_{n\to\infty} x_n = \infty$, buktikan bahwa
terdapat subbarisan $\{ x_{n_i} \}_{i=1}^\infty$ yang konvergen ke $\infty$.
Kemudian buktikan hasil yang sama untuk $-\infty$, lalu buktikan kedua pernyataan tersebut untuk
$\liminf$.
\end{exercise}


%%%%%%%%%%%%%%%%%%%%%%%%%%%%%%%%%%%%%%%%%%%%%%%%%%%%%%%%%%%%%%%%%%%%%%%%%%%%%%

\sectionnewpage
\section{Barisan Cauchy}
\label{sec:cauchy}

\sectionnotes{kurang dari satu pertemuan kuliah}

Sering kali kita ingin mendeskripsikan suatu bilangan tertentu
dengan barisan yang konvergen ke bilangan tersebut.  Dalam hal ini, bilangan
itu sendiri tidak mungkin digunakan dalam bukti bahwa barisan tersebut konvergen.
Akan sangat berguna jika kita dapat memeriksa kekonvergenan tanpa mengetahui
limitnya.

\begin{defn}
Sebuah barisan $\{ x_n \}_{n=1}^\infty$ disebut \emph{\myindex{barisan Cauchy}}%
\footnote{%
Dinamai menurut matematikawan Prancis
\href{https://en.wikipedia.org/wiki/Cauchy}{Augustin-Louis Cauchy} (1789--1857).} jika
untuk setiap $\epsilon > 0$ terdapat $M \in \N$ sedemikian sehingga
untuk semua $n \geq M$ dan semua $k \geq M$, berlaku
\begin{equation*}
\sabs{x_n - x_k} < \epsilon .
\end{equation*}
\end{defn}

Secara informal, sifat Cauchy berarti bahwa pada akhirnya semua suku barisan
dapat dibuat sedekat apa pun satu sama lain.  Kita mungkin menduga bahwa barisan
semacam itu konvergen, dan dugaan ini benar karena $\R$ memiliki
\hyperref[defn:lub]{sifat batas atas terkecil}.  Sebelum membuktikan fakta ini,
mari kita lihat beberapa contoh.

\begin{example}
Barisan $\{ \nicefrac{1}{n} \}_{n=1}^\infty$ merupakan barisan Cauchy.

Bukti:  Untuk $\epsilon > 0$ yang diberikan, pilih $M$ sedemikian sehingga
$M > \nicefrac{2}{\epsilon}$.  Maka, untuk $n,k \geq M$,
berlaku $\nicefrac{1}{n} < \nicefrac{\epsilon}{2}$
dan
$\nicefrac{1}{k} < \nicefrac{\epsilon}{2}$.  Oleh karena itu, untuk $n, k \geq M$,
berlaku
\begin{equation*}
\abs{\frac{1}{n} - \frac{1}{k}}
\leq
\abs{\frac{1}{n}} + \abs{\frac{1}{k}}
< \frac{\epsilon}{2} + \frac{\epsilon}{2} = \epsilon.
\end{equation*}
\end{example}

\begin{example}
Barisan $\bigl\{ {(-1)}^n \bigr\}_{n=1}^\infty$ bukan barisan Cauchy.

Bukti:
Untuk sembarang $M \in \N$, pilih sembarang bilangan genap $n \geq M$, lalu
tetapkan $k \coloneqq n+1$.  Maka
\begin{equation*}
\begin{split}
\babs{{(-1)}^n - {(-1)}^k}
& =
\babs{{(-1)}^n - {(-1)}^{n+1}}
\\
& =
\babs{1 - (-1)} = 2 .
\end{split}
\end{equation*}
Oleh karena itu, untuk sembarang $\epsilon \leq 2$, definisi tersebut tidak dapat dipenuhi,
sehingga barisan ini bukan barisan Cauchy.
\end{example}

%\begin{example}
%The sequence $\{ \frac{n+1}{n} \}_{n=1}^\infty$ is a Cauchy sequence.
%
%Proof:  Given $\epsilon > 0$, find $M$ such that
%$M > \nicefrac{2}{\epsilon}$.  Then for $n,k \geq M$,
%we have $\nicefrac{1}{n} < \nicefrac{\epsilon}{2}$
%and
%$\nicefrac{1}{k} < \nicefrac{\epsilon}{2}$.  Therefore, for $n, k \geq M$,
%we have
%\begin{equation*}
%\begin{split}
%\abs{\frac{n+1}{n} - \frac{k+1}{k}}
%& =
%\abs{\frac{k(n+1)-n(k+1)}{nk}}
%\\
%& =
%\abs{\frac{kn+k-nk-n}{nk}}
%\\
%& =
%\abs{\frac{k-n}{nk}}
%\\
%& \leq
%\abs{\frac{k}{nk}}
%+
%\abs{\frac{-n}{nk}}
%\\
%& = \frac{1}{n} + \frac{1}{k}
%< \frac{\epsilon}{2} + \frac{\epsilon}{2} = \epsilon .
%\end{split}
%\end{equation*}
%\end{example}

\begin{prop}
Jika suatu barisan adalah barisan Cauchy, maka barisan tersebut terbatas.
\end{prop}

\begin{proof}
Andaikan $\{ x_n \}_{n=1}^\infty$ merupakan barisan Cauchy.  Pilih $M$ sedemikian sehingga untuk semua
$n,k \geq M$, berlaku $\sabs{x_n-x_k} < 1$.  Khususnya,
untuk semua $n \geq M$,
\begin{equation*}
\sabs{x_n - x_M} < 1 .
\end{equation*}
Berdasarkan ketaksamaan segitiga terbalik,
$\sabs{x_n} - \sabs{x_M} \leq \sabs{x_n - x_M} < 1$.  Jadi, untuk $n \geq M$,
\begin{equation*}
\sabs{x_n} < 1 + \sabs{x_M}.
\end{equation*}
Tetapkan
\begin{equation*}
B \coloneqq \max \bigl\{ \sabs{x_1}, \sabs{x_2}, \ldots,
\sabs{x_{M-1}}, 1+ \sabs{x_M} \bigr\} .
\end{equation*}
Maka $\sabs{x_n} \leq B$ untuk semua $n \in \N$.
\end{proof}

\begin{thm}
Suatu barisan bilangan real merupakan barisan Cauchy jika dan hanya jika barisan tersebut konvergen.
\end{thm}

\begin{proof}
Andaikan $\{ x_n \}_{n=1}^\infty$ konvergen ke $x$,
dan misalkan $\epsilon > 0$ diberikan.
Maka terdapat $M$ sedemikian sehingga untuk $n \geq M$,
\begin{equation*}
\sabs{x_n - x} < \frac{\epsilon}{2} .
\end{equation*}
Jadi, untuk $n \geq M$ dan $k \geq M$,
\begin{equation*}
\sabs{x_n - x_k} = 
\sabs{x_n - x + x - x_k}
\leq \sabs{x_n-x} + \sabs{x-x_k} < \frac{\epsilon}{2} + \frac{\epsilon}{2} =
\epsilon .
\end{equation*}

Arah tersebut mudah.  Sekarang andaikan $\{ x_n \}_{n=1}^\infty$ merupakan barisan Cauchy.
Kita telah menunjukkan bahwa $\{ x_n \}_{n=1}^\infty$ terbatas.
Untuk barisan terbatas, liminf dan limsup ada; di sinilah kita menggunakan
\hyperref[defn:lub]{sifat batas atas terkecil}.
Jika dapat ditunjukkan bahwa
\begin{equation*}
\liminf_{n\to \infty} x_n = \limsup_{n\to\infty} x_n ,
\end{equation*}
maka $\{ x_n \}_{n=1}^\infty$ harus konvergen berdasarkan \propref{liminfsupconv:prop}.


Definisikan $a \coloneqq \limsup_{n\to\infty} x_n$ dan
$b \coloneqq \liminf_{n\to\infty} x_n$.
Berdasarkan \thmref{subseqlimsupinf:thm}, terdapat subbarisan
$\{ x_{n_i} \}_{i=1}^\infty$ dan
$\{ x_{m_i} \}_{i=1}^\infty$ sedemikian sehingga
\begin{equation*}
\lim_{i\to\infty} x_{n_i} = a
\qquad \text{dan} \qquad
\lim_{i\to\infty} x_{m_i} = b.
\end{equation*}
Untuk $\epsilon > 0$ yang diberikan,
terdapat $M_1$ sedemikian sehingga
$\sabs{x_{n_i} - a} < \nicefrac{\epsilon}{3}$ untuk semua $i \geq M_1$
dan terdapat $M_2$ sedemikian sehingga
$\sabs{x_{m_i} - b} < \nicefrac{\epsilon}{3}$ untuk semua $i \geq M_2$.
Terdapat pula $M_3$ sedemikian sehingga
$\sabs{x_n-x_k} < \nicefrac{\epsilon}{3}$
untuk semua $n,k \geq M_3$.
Tetapkan $M \coloneqq \max \{ M_1, M_2, M_3 \}$.
Jika $i \geq M$, maka $n_i \geq M$ dan $m_i \geq M$.  Jadi,
\begin{equation*}
\begin{split}
\sabs{a-b} & =
\sabs{a-x_{n_i}+x_{n_i}
-x_{m_i}+x_{m_i}
-b} \\
& \leq
\sabs{a-x_{n_i}}
+ \sabs{x_{n_i} -x_{m_i}}
+ \sabs{x_{m_i} -b} \\
& <
\frac{\epsilon}{3}
+
\frac{\epsilon}{3}
+
\frac{\epsilon}{3}
= \epsilon .
\end{split}
\end{equation*}
Karena $\sabs{a-b} < \epsilon$ untuk semua $\epsilon > 0$, kita memperoleh $a=b$ dan
barisan tersebut konvergen.
\end{proof}

\begin{remark}
Pernyataan teorema ini kadang digunakan untuk mendefinisikan
sifat kelengkapan bilangan real.  Suatu himpunan disebut
\emph{\myindex{lengkap secara Cauchy}} (atau kadang cukup disebut \emph{\myindex{lengkap}})
jika setiap barisan Cauchy konvergen ke suatu unsur di dalam himpunan tersebut.
Di atas, kita telah membuktikan bahwa karena $\R$ memiliki
\hyperref[defn:lub]{sifat batas atas terkecil}, $\R$ lengkap secara Cauchy.
Kita dapat membentuk $\R$ dengan
\myquote{melengkapi} $\Q$, yaitu dengan \myquote{menambahkan} titik secukupnya agar semua
barisan Cauchy konvergen (rincian ini kita lewatkan).
Medan yang dihasilkan memiliki
sifat batas atas terkecil.
Keuntungan mendefinisikan kelengkapan melalui barisan Cauchy adalah bahwa
definisi tersebut dapat digeneralisasi ke konteks yang lebih abstrak,
seperti ruang metrik; lihat \chapterref{ms:chapter}.
\end{remark}

Kriteria Cauchy lebih kuat daripada sekadar syarat bahwa
$\sabs{x_{n+1}-x_n}$ (atau $\sabs{x_{n+j}-x_n}$ untuk $j$ tetap) menuju nol ketika
$n$ menuju tak hingga.  Ketika kita membahas jumlah parsial deret harmonik
(lihat \exampleref{example:harmonicseries} pada bagian berikutnya), kita akan memperoleh
barisan yang memenuhi $x_{n+1}-x_n = \frac{1}{n+1}$,
tetapi $\{ x_n \}_{n=1}^\infty$ divergen.
Bahkan, untuk barisan tersebut,
$\lim_{n\to\infty} \sabs{x_{n+j}-x_n} = 0$ untuk
setiap $j \in \N$ (bandingkan dengan \exerciseref{exercise:badnocauchy}).
Pokok penting dalam definisi Cauchy adalah bahwa $n$ dan $k$
bervariasi secara independen dan dapat terpisah sejauh apa pun.

\subsection{Latihan}

\begin{exercise}
Buktikan secara langsung dari definisi barisan Cauchy bahwa
$\bigl\{ \frac{n^2-1}{n^2} \bigr\}_{n=1}^\infty$ adalah barisan Cauchy.
\end{exercise}

\begin{exercise}
Misalkan $\{ x_n \}_{n=1}^\infty$ adalah barisan sedemikian sehingga
terdapat suatu konstanta positif $C < 1$ dan, untuk semua $n \geq 2$,
\begin{equation*}
\sabs{x_{n+1} - x_n} \leq C \sabs{x_{n}-x_{n-1}} .
\end{equation*}
Buktikan bahwa $\{ x_n \}_{n=1}^\infty$ adalah barisan Cauchy.
Petunjuk: Anda boleh langsung menggunakan rumus berikut (untuk $C \neq 1$):
\begin{equation*}
1+ C+ C^2 + \cdots + C^n = \frac{1-C^{n+1}}{1-C}.
\end{equation*}
\end{exercise}

\begin{exercise}[Menantang]
Misalkan $F$ adalah medan terurut yang memuat bilangan-bilangan rasional
$\Q$, dan $\Q$ padat,
yaitu, setiap kali $x,y \in F$ memenuhi $x < y$,
terdapat $q \in \Q$ sedemikian sehingga $x < q < y$.
Katakan bahwa barisan bilangan rasional $\{ x_n \}_{n=1}^\infty$ adalah barisan Cauchy
jika untuk setiap $\epsilon \in \Q$ dengan $\epsilon > 0$, terdapat
$M$ sedemikian sehingga untuk semua $n,k \geq M$, berlaku $\sabs{x_n-x_k} < \epsilon$.
Misalkan setiap barisan Cauchy yang terdiri atas bilangan rasional mempunyai limit di $F$.
Buktikan bahwa $F$ mempunyai \hyperref[defn:lub]{sifat batas atas terkecil}.
\end{exercise}

\begin{exercise}
Misalkan $\{ x_n \}_{n=1}^\infty$ dan $\{ y_n \}_{n=1}^\infty$ adalah barisan sedemikian
sehingga $\lim_{n\to\infty} y_n =0$.  Misalkan bahwa untuk semua $k \in \N$
dan
untuk semua $m \geq k$, berlaku
\begin{equation*}
\sabs{x_m-x_k} \leq y_k .
\end{equation*}
Tunjukkan bahwa $\{ x_n \}_{n=1}^\infty$ adalah barisan Cauchy.
\end{exercise}

\begin{exercise}
Misalkan barisan Cauchy $\{ x_n \}_{n=1}^\infty$ sedemikian sehingga untuk setiap $M \in \N$,
terdapat $k \geq M$ dan $n \geq M$ sedemikian sehingga
$x_k < 0$ dan $x_n > 0$.  Dengan hanya menggunakan definisi barisan Cauchy
dan barisan konvergen, tunjukkan bahwa
barisan tersebut konvergen ke $0$.
\end{exercise}

\begin{exercise}
Misalkan $\sabs{x_n-x_k} \leq \nicefrac{n}{k^2}$ untuk semua $n$ dan $k$.
Tunjukkan bahwa $\{ x_n \}_{n=1}^\infty$ adalah barisan Cauchy.\footnote{Setelah membuktikan bahwa
barisan tersebut adalah barisan Cauchy, sebagai latihan tambahan, buktikan bahwa barisan itu konstan.}
\end{exercise}

\begin{exercise}
Misalkan $\{ x_n \}_{n=1}^\infty$ adalah barisan Cauchy sedemikian sehingga untuk tak berhingga banyak
$n$, berlaku $x_n = c$.  Dengan hanya menggunakan definisi barisan Cauchy, buktikan 
bahwa $\lim\limits_{n\to\infty} x_n = c$.
\end{exercise}

\begin{exercise}
Tentukan apakah pernyataan berikut benar atau salah; buktikan atau berikan contoh tandingan: Jika $\{ x_n \}_{n=1}^\infty$ adalah barisan
Cauchy, maka terdapat $M$
sedemikian sehingga untuk semua $n \geq M$, berlaku
$\sabs{x_{n+1}-x_n}
\leq
\sabs{x_{n}-x_{n-1}}$.
\end{exercise}

\begin{exercise}
Misalkan 
$\{ x_n \}_{n=1}^\infty$ adalah barisan sedemikian sehingga untuk setiap $\ell$ dan $k$,
berlaku
$\sabs{x_\ell-x_k} \leq y_\ell + z_k$, dengan
$\{ y_n \}_{n=1}^\infty$ dan 
$\{ z_n \}_{n=1}^\infty$ merupakan barisan yang konvergen ke $0$.
Tunjukkan bahwa 
$\{ x_n \}_{n=1}^\infty$ adalah barisan Cauchy.
\end{exercise}


%%%%%%%%%%%%%%%%%%%%%%%%%%%%%%%%%%%%%%%%%%%%%%%%%%%%%%%%%%%%%%%%%%%%%%%%%%%%%%

\sectionnewpage
\section{Deret}
\label{sec:series}

%mbxINTROSUBSECTION

\sectionnotes{2 kuliah}

Salah satu objek mendasar dalam matematika adalah deret.  Bahkan, ketika
landasan analisis sedang dikembangkan, motivasinya adalah
memahami deret.  Pemahaman tentang deret penting dalam penerapan
analisis.  Sebagai contoh, solusi persamaan diferensial sering kali
diberikan sebagai deret, dan persamaan diferensial menjadi dasar untuk memahami
hampir seluruh ilmu pengetahuan modern.

\subsection{Definisi}

\begin{defn}
Diberikan barisan $\{ x_n \}_{n=1}^\infty$, kita tuliskan objek formal
\begin{equation*}
\sum_{n=1}^\infty x_n
\end{equation*}
\glsadd{not:series}
dan menyebutnya \emph{\myindex{deret}}.  Suatu deret
\emph{konvergen}\index{konvergen!deret} jika barisan
$\{ s_k \}_{k=1}^\infty$
yang didefinisikan oleh
\glsadd{not:finitesum}
\begin{equation*}
s_k \coloneqq \sum_{n=1}^k x_n = x_1 + x_2 + \cdots + x_k
\end{equation*}
konvergen.
Bilangan-bilangan $s_k$ disebut
\emph{\myindex{jumlah parsial}}.
Jika deret tersebut konvergen, kita tuliskan
\begin{equation*}
\sum_{n=1}^\infty x_n =  \lim\limits_{k\to\infty} s_k .
\end{equation*}
Dalam hal ini, kita sedikit menyederhanakan persoalan dan memperlakukan
$\sum_{n=1}^\infty x_n$ sebagai bilangan.

Jika barisan $\{ s_k \}_{k=1}^\infty$ divergen,
kita mengatakan bahwa deret tersebut \emph{divergen}\index{divergen!deret}.
Dalam hal ini, $\sum_{n=1}^\infty x_n$ hanyalah objek formal dan bukan bilangan.
\end{defn}

Dengan kata lain, untuk deret konvergen, kita mempunyai
\begin{equation*}
\sum_{n=1}^\infty x_n
=
\lim_{k\to\infty} 
\sum_{n=1}^k x_n .
\end{equation*}
Kesamaan ini hanya berlaku jika limit di ruas
kanan benar-benar ada.  Jika deret tersebut tidak konvergen, ruas kanan
tidak bermakna (limitnya tidak ada).
Karena itu, berhati-hatilah karena 
$\sum_{n=1}^\infty x_n$ memiliki dua makna berbeda (notasi untuk deret itu sendiri 
atau limit jumlah-jumlah parsial), dan Anda harus menggunakan konteks
untuk membedakannya.

\begin{remark}
Kadang-kadang lebih mudah untuk memulai
deret dari indeks selain 1.  Sebagai contoh, kita dapat menuliskan
\begin{equation*}
\sum_{n=0}^\infty r^n = \sum_{n=1}^\infty r^{n-1} .
\end{equation*}
Ruas kiri lebih mudah dituliskan.
\end{remark}

\begin{remark}
Deret $\sum_{n=1}^\infty x_n$ lazim dituliskan sebagai
\begin{equation*}
x_1 + x_2 + x_3 + \cdots
\end{equation*}
dengan pengertian bahwa elipsis tersebut menyatakan deret dan
bukan sekadar penjumlahan.  Kita tidak menggunakan notasi ini karena notasi informal
semacam itu dapat menimbulkan kesalahan dalam pembuktian.
\end{remark}

\begin{example}
Deret
\begin{equation*}
\sum_{n=1}^\infty \frac{1}{2^n}
\end{equation*}
konvergen dan limitnya adalah 1.  Dengan kata lain,
\begin{equation*}
\sum_{n=1}^\infty \frac{1}{2^n} = 
\lim_{k\to\infty} \sum_{n=1}^k \frac{1}{2^n} = 
1 .
\end{equation*}

Bukti: Kita memerlukan kesamaan
\begin{equation*}
\left( \sum_{n=1}^k \frac{1}{2^n} \right)
+ \frac{1}{2^k}
= 1 .
\end{equation*}
Kesamaan tersebut langsung berlaku ketika $k=1$.  Bukti untuk $k$ umum
mengikuti dari \hyperref[induction:thm]{induksi}, yang kita serahkan kepada
pembaca.  Lihat \figureref{figcutseries} untuk ilustrasinya.
\begin{myfigureht}
\myincludepdft{figcutseries}{%
Diagram garis bilangan real.  Titik 0 berada di kiri dan titik
1 berada di kanan.  Jarak sebesar setengah ditandai dari 0 ke
setengah.  Kemudian jarak sebesar seperempat ditandai dari setengah
ke titik tiga perempat.  Selanjutnya, jarak sebesar seperdelapan ditandai dari
tiga perempat ke tujuh perdelapan.  Jarak setengah ditambah seperempat ditambah
seperdelapan juga ditandai, yang menggabungkan ketiga jarak tersebut.
Terakhir, sisa jarak sebesar seperdelapan dari tujuh perdelapan ke 1 juga
ditandai.}
\caption{Kesamaan 
$\left( \sum_{n=1}^k \frac{1}{2^n} \right)
+ \frac{1}{2^k}
= 1$ yang diilustrasikan untuk $k=3$.\label{figcutseries}}
\end{myfigureht}

Misalkan $s_k$ adalah jumlah parsial.  Kita tuliskan
\begin{equation*}
\sabs{
1 - s_k 
}
=
\abs{
1 - 
\sum_{n=1}^k \frac{1}{2^n}
}
=
\abs{\frac{1}{2^k}} = 
\frac{1}{2^k} .
\end{equation*}
Barisan $\bigl\{ \frac{1}{2^k} \bigr\}_{k=1}^\infty$, dan
karena itu $\bigl\{ \sabs{1-s_k} \bigr\}_{k=1}^\infty$,
konvergen ke nol.  Jadi, $\{ s_k \}_{k=1}^\infty$ konvergen ke $1$.
\end{example}

\begin{prop} \label{geometric:prop}
Misalkan
$-1 < r < 1$.  Maka \emph{\myindex{deret geometri}}
$\sum_{n=0}^\infty r^n$ konvergen, dan
\begin{equation*}
\sum_{n=0}^\infty r^n = \frac{1}{1-r} .
\end{equation*}
\end{prop}

Rincian buktinya diserahkan sebagai latihan.
Pembuktiannya dilakukan dengan menunjukkan bahwa 
\begin{equation*}
\sum_{n=0}^{k-1} r^n = \frac{1-r^k}{1-r} ,
\end{equation*}
kemudian mengambil limit ketika $k$ menuju $\infty$.
Deret geometri merupakan salah satu deret terpenting, dan bahkan merupakan
salah satu dari sedikit deret yang limitnya dapat kita tentukan secara begitu eksplisit.

\medskip

\pagebreak[2]
Seperti pada barisan, kita dapat membicarakan \emph{\myindex{ekor deret}}.

\begin{prop}
Misalkan $\sum_{n=1}^\infty x_n$ adalah deret.  Misalkan $M \in \N$.  Maka
\begin{equation*}
\sum_{n=1}^\infty x_n \quad \text{konvergen jika dan hanya jika} \quad
\sum_{n=M}^\infty x_n \quad \text{konvergen.}
\end{equation*}
\end{prop}

\begin{proof}
Kita meninjau jumlah-jumlah parsial kedua deret tersebut (untuk $k \geq M$):
\begin{equation*}
\sum_{n=1}^{k} x_n
=
\left(
\sum_{n=1}^{M-1} x_n
\right)
+
\sum_{n=M}^{k} x_n .
\end{equation*}
Perhatikan bahwa 
$\sum_{n=1}^{M-1} x_n$ adalah bilangan tetap.  Gunakan
\propref{prop:contalg} untuk menyelesaikan bukti.
\end{proof}

\subsection{Deret Cauchy}

\begin{defn}
Deret $\sum_{n=1}^\infty x_n$ disebut \emph{Cauchy} atau
\emph{deret Cauchy}\index{deret Cauchy}
jika barisan jumlah parsial $\{ s_n \}_{n=1}^\infty$ merupakan barisan Cauchy.
\end{defn}

Barisan bilangan real konvergen jika dan hanya jika barisan tersebut
bersifat Cauchy.  Oleh karena itu, suatu deret konvergen jika dan hanya jika deret tersebut bersifat Cauchy.
Deret $\sum_{n=1}^\infty x_n$ bersifat Cauchy jika dan hanya jika untuk setiap $\epsilon > 0$,
terdapat $M \in \N$ sedemikian sehingga untuk setiap $n \geq M$
dan $k \geq M$, berlaku
\begin{equation*}
\abs{ \left( \sum_{i=1}^k x_i \right) - \left( \sum_{i=1}^n x_i \right) }
< \epsilon .
\end{equation*}
Tanpa mengurangi keumuman, kita andaikan $n < k$.  Maka kita tuliskan
\begin{equation*}
\abs{ \left( \sum_{i=1}^k x_i \right) - \left( \sum_{i=1}^n x_i \right) }
=
\abs{ \sum_{i={n+1}}^k x_i }
< \epsilon .
\end{equation*}
Kita telah membuktikan proposisi sederhana berikut.

\begin{prop} \label{prop:cachyser}
Deret $\sum_{n=1}^\infty x_n$ bersifat Cauchy
jika dan hanya jika untuk setiap $\epsilon > 0$, 
terdapat $M \in \N$ sedemikian sehingga untuk setiap $n \geq M$
dan setiap $k > n$,
\begin{equation*}
\abs{ \sum_{i={n+1}}^k x_i }
< \epsilon .
\end{equation*}
\end{prop}

\subsection{Sifat-sifat dasar}

\begin{prop}
Misalkan $\sum_{n=1}^\infty x_n$ adalah deret konvergen.  Maka
barisan $\{ x_n \}_{n=1}^\infty$ konvergen dan
\begin{equation*}
\lim_{n\to\infty} x_n = 0.
\end{equation*}
\end{prop}

\begin{proof}
Diberikan $\epsilon > 0$.  Karena $\sum_{n=1}^\infty x_n$ konvergen, deret tersebut bersifat Cauchy.
Jadi, kita memperoleh $M$ sedemikian sehingga untuk setiap $n \geq M$,
\begin{equation*}
\epsilon > 
\abs{ \sum_{i={n+1}}^{n+1} x_i }
=
\sabs{ x_{n+1} } .
\end{equation*}
Dengan demikian, untuk setiap $n \geq M+1$, berlaku $\sabs{x_{n}} < \epsilon$.
\end{proof}

\begin{example}
Jika $r \geq 1$ atau $r \leq -1$, deret geometri $\sum_{n=0}^\infty r^n$
divergen.

Bukti: $\sabs{r^n} = \sabs{r}^n \geq 1^n = 1$.  Suku-sukunya tidak menuju nol
sehingga deret tersebut tidak mungkin konvergen.
\end{example}

Jadi, jika suatu deret konvergen, suku-suku deret tersebut menuju nol.
Namun, implikasi ini hanya berlaku satu arah.
Mari kita berikan sebuah contoh.

\begin{example} \label{example:harmonicseries}
Deret $\sum_{n=1}^\infty \nicefrac{1}{n}$ divergen (meskipun
$\lim_{n\to\infty} \nicefrac{1}{n} = 0$).
Ini adalah \emph{\myindex{deret harmonik}} yang terkenal%
\footnote{Divergensi deret harmonik telah diketahui 
jauh sebelum teori deret dirumuskan secara ketat.  Bukti yang
kita berikan merupakan bukti paling awal dan dikemukakan oleh
\href{https://en.wikipedia.org/wiki/Oresme}{Nicole Oresme}
(1323?--1382).}.

Bukti: Kita akan menunjukkan bahwa barisan jumlah parsial tak terbatas, sehingga
tidak mungkin konvergen.
Tuliskan jumlah parsial $s_n$ untuk $n = 2^k$ sebagai berikut:
\begin{align*}
 s_1 & = 1 , \\
 s_2 & = \left( 1 \right) + \left( \frac{1}{2} \right) , \\
 s_4 & = \left( 1 \right) + \left( \frac{1}{2} \right) +
        \left( \frac{1}{3} + \frac{1}{4} \right) , \\
 s_8 & = \left( 1 \right) + \left( \frac{1}{2} \right) +
        \left( \frac{1}{3} + \frac{1}{4} \right) +
        \left( \frac{1}{5} + \frac{1}{6} + \frac{1}{7} + \frac{1}{8} \right) , \\
& ~~ \vdots \\
 s_{2^k} & = 
1 + 
\sum_{i=1}^k
\left(
\sum_{m=2^{i-1}+1}^{2^i} \frac{1}{m}
\right) .
\end{align*}
Perhatikan bahwa $\nicefrac{1}{3} + \nicefrac{1}{4} \geq \nicefrac{1}{4} + \nicefrac{1}{4} =
\nicefrac{1}{2}$ dan
$\nicefrac{1}{5} + \nicefrac{1}{6} + \nicefrac{1}{7} + \nicefrac{1}{8}
\geq \nicefrac{1}{8} + \nicefrac{1}{8} + \nicefrac{1}{8} + \nicefrac{1}{8} =
\nicefrac{1}{2}$.  Secara lebih umum,
\begin{equation*}
\sum_{m=2^{k-1}+1}^{2^k} \frac{1}{m}
\geq
\sum_{m=2^{k-1}+1}^{2^k} \frac{1}{2^k}
=
(2^{k-1}) \frac{1}{2^k} = \frac{1}{2} .
\end{equation*}
Oleh karena itu,
\begin{equation*}
s_{2^k} = 
1 + 
\sum_{i=1}^k
\left(
\sum_{m=2^{i-1}+1}^{2^i} \frac{1}{m}
\right) 
\geq
1 + \sum_{i=1}^k \frac{1}{2} = 1 + \frac{k}{2} .
\end{equation*}
Karena $\{ \nicefrac{k}{2} \}_{k=1}^\infty$ tak terbatas berdasarkan
\hyperref[thm:arch:i]{Sifat Archimedes},
$\{ s_{2^k} \}_{k=1}^\infty$ tak terbatas,
sehingga $\{ s_n \}_{n=1}^\infty$ juga tak terbatas.
Jadi, $\{ s_n \}_{n=1}^\infty$ divergen, dan akibatnya $\sum_{n=1}^\infty \nicefrac{1}{n}$ divergen.
\end{example}

Seperti jumlah berhingga, deret konvergen berperilaku linear.
Artinya, kita dapat mengalikannya dengan konstanta
dan menjumlahkannya, dan operasi-operasi ini dilakukan suku demi suku.

\begin{prop}[Linearitas deret]\index{linearitas deret}
Misalkan $\alpha \in \R$, sedangkan $\sum_{n=1}^\infty x_n$ dan $\sum_{n=1}^\infty y_n$ adalah
deret konvergen.  Maka
\begin{enumerate}[(i)]
\item
$\sum_{n=1}^\infty \alpha x_n$ adalah deret konvergen dan
\begin{equation*}
\sum_{n=1}^\infty \alpha x_n
=
\alpha \sum_{n=1}^\infty x_n .
\end{equation*}
\item
$\sum_{n=1}^\infty ( x_n + y_n )$ adalah deret konvergen dan
\begin{equation*}
\sum_{n=1}^\infty ( x_n + y_n ) 
=
\left( \sum_{n=1}^\infty x_n \right)
+
\left( \sum_{n=1}^\infty y_n \right) .
\end{equation*}
\end{enumerate}
\end{prop}

\begin{proof}
Untuk butir pertama,
kita cukup menuliskan jumlah parsial ke-$k$:
\begin{equation*}
\sum_{n=1}^k \alpha x_n
=
\alpha \left( \sum_{n=1}^k x_n \right) .
\end{equation*}
Kita meninjau ruas kanan dan memperhatikan bahwa kelipatan konstan dari
barisan konvergen
bersifat konvergen.  Dengan demikian, kita mengambil limit kedua ruas untuk memperoleh
hasil tersebut.

Untuk butir kedua, kita juga meninjau
jumlah parsial ke-$k$:
\begin{equation*}
\sum_{n=1}^k ( x_n + y_n ) 
=
\left( \sum_{n=1}^k x_n \right)
+
\left( \sum_{n=1}^k y_n \right) .
\end{equation*}
Kita meninjau ruas kanan dan memperhatikan bahwa jumlah barisan-barisan konvergen
bersifat konvergen.  Dengan demikian, kita mengambil limit kedua ruas untuk memperoleh
proposisi tersebut.
\end{proof}

Contoh penerapan berguna dari butir pertama adalah rumus
berikut.  Jika $\sabs{r} < 1$ dan $i \in \N$, maka
\begin{equation*}
\sum_{n=i}^\infty r^n = \frac{r^i}{1-r} .
\end{equation*}
Rumus tersebut diperoleh dengan menggunakan deret geometri dan mengalikannya dengan
$r^i$:
\begin{equation*}
r^i \sum_{n=0}^\infty r^n =
\sum_{n=0}^\infty r^{n+i}
=
\sum_{n=i}^\infty r^n .
\end{equation*}

Mengalikan deret tidak sesederhana menjumlahkannya; lihat bagian
berikutnya.
Tentu saja, kita tidak mengalikannya
suku demi suku.  Cara tersebut bahkan tidak berlaku untuk jumlah berhingga:
$(a+b)(c+d) \neq ac+bd$.

\subsection{Konvergensi mutlak}

Karena barisan monoton lebih mudah ditangani daripada barisan sebarang,
deret $\sum_{n=1}^\infty x_n$ dengan $x_n \geq 0$ untuk semua $n$
umumnya juga lebih mudah ditangani.  Barisan jumlah parsialnya monoton naik
dan konvergen jika terbatas di atas.
Mari kita nyatakan pernyataan ini secara formal sebagai proposisi.

\begin{prop}
Jika $x_n \geq 0$ untuk semua $n$, maka $\sum_{n=1}^\infty x_n$ konvergen jika dan hanya jika
barisan jumlah parsialnya terbatas di atas.
\end{prop}

Limit barisan monoton naik adalah supremumnya.
Jadi, ketika $x_n \geq 0$ untuk semua $n$, kita memperoleh
ketaksamaan
\begin{equation*}
\sum_{n=1}^k x_n \leq
\sum_{n=1}^\infty x_n .
\end{equation*}
Jika kita memperbolehkan limit tak hingga, ketaksamaan tersebut tetap
berlaku bahkan ketika deret divergen menuju tak hingga, meskipun dalam hal itu
ketaksamaan tersebut tidak terlalu berguna.

Kita akan melihat bahwa kriteria umum berikut untuk kekonvergenan deret 
berdampak besar pada cara deret dapat dimanipulasi.

\begin{defn}
Deret $\sum_{n=1}^\infty x_n$
\emph{\myindex{konvergen mutlak}}\index{konvergensi mutlak} jika
deret $\sum_{n=1}^\infty \sabs{x_n}$ konvergen.
Jika suatu deret konvergen tetapi tidak konvergen mutlak, kita mengatakan bahwa
deret tersebut \emph{\myindex{konvergen bersyarat}}.\index{konvergensi bersyarat}
\end{defn}

\begin{prop}
Jika deret $\sum_{n=1}^\infty x_n$ konvergen mutlak, maka deret tersebut konvergen.
\end{prop}

\begin{proof}
Suatu deret konvergen jika dan hanya jika deret tersebut bersifat Cauchy.
Misalkan $\sum_{n=1}^\infty \sabs{x_n}$ bersifat Cauchy.  Artinya, untuk setiap $\epsilon > 0$,
terdapat $M$ sedemikian sehingga untuk semua $k \geq M$ dan semua $n > k$, berlaku 
\begin{equation*}
\sum_{i=k+1}^n \sabs{x_i} 
=
\abs{ \sum_{i=k+1}^n \sabs{x_i} }
<
\epsilon .
\end{equation*}
Kita menerapkan ketaksamaan segitiga untuk jumlah berhingga dan memperoleh
\begin{equation*}
\abs{ \sum_{i=k+1}^n x_i }
\leq
\sum_{i=k+1}^n \sabs{x_i}
<
\epsilon .
\end{equation*}
Dengan demikian, $\sum_{n=1}^\infty x_n$ bersifat Cauchy, sehingga deret tersebut konvergen.
\end{proof}

Jika $\sum_{n=1}^\infty x_n$ konvergen mutlak, limit
$\sum_{n=1}^\infty x_n$ dan $\sum_{n=1}^\infty \sabs{x_n}$ pada umumnya berbeda.  Menghitung yang satu
tidak membantu kita menghitung yang lain.  Namun, perhitungan di atas menghasilkan
ketaksamaan yang berguna untuk deret konvergen mutlak,
yaitu versi deret dari ketaksamaan segitiga,
yang pembuktiannya kita serahkan sebagai latihan:
\begin{equation*}
\abs{ \sum_{i=1}^\infty x_i }
\leq
\sum_{i=1}^\infty \sabs{x_i} .
\end{equation*}

Deret konvergen mutlak mempunyai banyak sifat yang baik.
Sebagai contoh, deret konvergen mutlak
dapat disusun ulang secara sebarang, atau deret-deret semacam itu dapat
dikalikan dengan mudah.  Sebaliknya, deret konvergen bersyarat
sering kali tidak berperilaku seperti yang diharapkan.  Lihat bagian berikutnya.

Kita serahkan sebagai latihan untuk menunjukkan bahwa
\begin{equation*}
\sum_{n=1}^\infty \frac{{(-1)}^n}{n}
\end{equation*}
konvergen, meskipun pembaca sebaiknya menyelesaikan bagian ini sebelum mencobanya.
Di sisi lain, kita telah membuktikan bahwa
\begin{equation*}
\sum_{n=1}^\infty \frac{1}{n}
\end{equation*}
divergen.  Oleh karena itu,
$\sum_{n=1}^\infty \frac{{(-1)}^n}{n}$ merupakan deret konvergen bersyarat.

\subsection{Uji perbandingan dan deret-\texorpdfstring{$p$}{p}}

Kita telah mencatat bahwa agar deret dengan suku-suku positif konvergen,
suku-sukunya tidak hanya harus menuju nol, tetapi harus menuju nol
\myquote{cukup cepat.}  Jika kita mengetahui kekonvergenan suatu deret,
kita dapat menggunakan uji perbandingan berikut untuk melihat apakah suku-suku deret
lain menuju nol \myquote{cukup cepat.}

\begin{samepage}
\begin{prop}[Uji perbandingan]\index{uji perbandingan untuk deret}
Misalkan $\sum_{n=1}^\infty x_n$ dan $\sum_{n=1}^\infty y_n$ adalah deret sedemikian sehingga $0 \leq x_n \leq y_n$
untuk semua $n \in \N$.
\begin{enumerate}[(i)]
\item Jika $\sum_{n=1}^\infty y_n$ konvergen, maka $\sum_{n=1}^\infty x_n$ juga konvergen.
\item Jika $\sum_{n=1}^\infty x_n$ divergen, maka $\sum_{n=1}^\infty y_n$ juga divergen.
\end{enumerate}
\end{prop}
\end{samepage}

\begin{proof}
Karena semua suku deret tersebut tak negatif, kedua barisan
jumlah parsialnya monoton naik.
Karena $x_n \leq y_n$ untuk semua $n$, jumlah-jumlah parsialnya
memenuhi, untuk semua $k$,
\begin{equation} \label{comptest:eq}
\sum_{n=1}^k x_n \leq \sum_{n=1}^k y_n .
\end{equation}
Jika deret $\sum_{n=1}^\infty y_n$ konvergen, jumlah-jumlah parsial deret tersebut
terbatas.  Oleh karena itu, ruas kanan \eqref{comptest:eq}
terbatas untuk semua $k$; terdapat $B \in \R$ sedemikian sehingga
$\sum_{n=1}^k y_n \leq B$ untuk semua $k$, sehingga
\begin{equation*}
\sum_{n=1}^k x_n \leq \sum_{n=1}^k y_n \leq B.
\end{equation*}
Dengan demikian, jumlah-jumlah parsial untuk $\sum_{n=1}^\infty x_n$
juga terbatas.  Karena jumlah-jumlah parsial tersebut membentuk barisan monoton naik,
barisan itu konvergen.  Jadi, butir pertama telah dibuktikan.

Di sisi lain, jika $\sum_{n=1}^\infty x_n$ divergen, barisan jumlah parsialnya
harus tak terbatas karena monoton naik.  Artinya, jumlah-jumlah parsial
$\sum_{n=1}^\infty x_n$ pada akhirnya lebih besar daripada sebarang bilangan real.  Dengan menggabungkan hal ini
dengan \eqref{comptest:eq}, kita melihat bahwa untuk setiap $B \in
\R$, terdapat $k$ sedemikian sehingga 
\begin{equation*}
B \leq \sum_{n=1}^k x_n \leq \sum_{n=1}^k y_n .
\end{equation*}
Dengan demikian, jumlah-jumlah parsial untuk $\sum_{n=1}^\infty y_n$ juga tak terbatas, dan
$\sum_{n=1}^\infty y_n$ juga divergen.
\end{proof}

Deret yang berguna untuk diterapkan bersama uji perbandingan adalah
deret-$p$%
\footnote{Kita belum mendefinisikan $x^p$ untuk $x > 0$ dan
sebarang $p \in \R$.  Definisinya adalah $x^p \coloneqq \exp ( p \ln x )$.
Kita akan mendefinisikan logaritma dan eksponensial dalam \sectionref{sec:logandexp}.
Untuk sementara, Anda dapat memikirkan $p$ rasional,
dengan $x^{k/m} = {(x^{1/m})}^{k}$.  Lihat juga \exerciseref{exercise:realpower}.}.

\begin{prop}[Deret-$p$ atau uji-$p$]%
\index{deret-p@deret-$p$}\index{uji-p@uji-$p$}
Untuk $p \in \R$, 
deret
\begin{equation*}
\sum_{n=1}^\infty \frac{1}{n^p}
\end{equation*}
konvergen jika dan hanya jika $p > 1$.
\end{prop}

\begin{proof}
Pertama, misalkan $p \leq 1$.
Karena $n \geq 1$, berlaku
$\frac{1}{n^p} \geq \frac{1}{n}$.  Karena
$\sum_{n=1}^\infty \frac{1}{n}$ divergen,
$\sum_{n=1}^\infty \frac{1}{n^p}$ harus divergen untuk semua $p \leq 1$ berdasarkan uji perbandingan.

Sekarang, misalkan $p > 1$.
Kita melanjutkan seperti pada
deret harmonik, tetapi alih-alih menunjukkan bahwa barisan
jumlah parsial tak terbatas, kita menunjukkan bahwa barisan tersebut terbatas.
Suku-suku deretnya positif, sehingga barisan jumlah parsial
monoton naik dan konvergen jika terbatas
di atas.
Misalkan $s_n$ menyatakan jumlah parsial ke-$n$.
\begin{align*}
 s_1 & = 1 , \\
 s_3 & = \left( 1 \right) + \left( \frac{1}{2^p} + \frac{1}{3^p} \right) , \\
 s_7 & = \left( 1 \right) + \left( \frac{1}{2^p} + \frac{1}{3^p} \right) +
        \left( \frac{1}{4^p} + \frac{1}{5^p} + \frac{1}{6^p} + \frac{1}{7^p} \right) , \\
& ~~ \vdots \\
 s_{2^k - 1} &= 
1 + 
\sum_{i=1}^{k-1}
\left(
\sum_{m=2^i}^{2^{i+1}-1} \frac{1}{m^p}
\right) .
\end{align*}
Alih-alih menaksir dari bawah, kita menaksir dari atas.
Karena $p$ positif, $2^p < 3^p$, sehingga
$\frac{1}{2^p} + \frac{1}{3^p} <
\frac{1}{2^p} + \frac{1}{2^p}$.  Demikian pula,
$\frac{1}{4^p} + \frac{1}{5^p} +
\frac{1}{6^p} + \frac{1}{7^p} <
\frac{1}{4^p} + \frac{1}{4^p} +
\frac{1}{4^p} + \frac{1}{4^p}$.  Oleh karena itu,
untuk semua $k \geq 2$,
\begin{equation*}
\begin{split}
s_{2^k-1}
& =
1+
\sum_{i=1}^{k-1}
\left(
\sum_{m=2^{i}}^{2^{i+1}-1} \frac{1}{m^p}
\right) 
\\
& <
1+
\sum_{i=1}^{k-1}
\left(
\sum_{m=2^{i}}^{2^{i+1}-1} \frac{1}{{(2^i)}^p}
\right) 
\\
& =
1+
\sum_{i=1}^{k-1}
\left(
\frac{2^i}{{(2^i)}^p}
\right) 
\\
& =
1+
\sum_{i=1}^{k-1}
{\left(
\frac{1}{2^{p-1}}
\right)}^i .
\end{split}
\end{equation*}
Karena $p > 1$, berlaku $\frac{1}{2^{p-1}} < 1$.
\propref{geometric:prop} menyatakan bahwa
\begin{equation*}
\sum_{i=1}^\infty
{\left(
\frac{1}{2^{p-1}}
\right)}^i
\end{equation*}
konvergen.  Jadi,
\begin{equation*}
s_{2^k-1} < 
1+
\sum_{i=1}^{k-1}
{\left(
\frac{1}{2^{p-1}}
\right)}^i 
\leq 
1+
\sum_{i=1}^\infty
{\left(
\frac{1}{2^{p-1}}
\right)}^i .
\end{equation*}
Untuk setiap $n$, terdapat $k \geq 2$ sedemikian sehingga $n \leq 2^k-1$, dan
karena $\{ s_n \}_{n=1}^\infty$ merupakan barisan monoton, $s_n \leq s_{2^k-1}$.
Jadi, untuk semua $n$,
\begin{equation*}
s_n < 
1+
\sum_{i=1}^\infty
{\left(
\frac{1}{2^{p-1}}
\right)}^i
\end{equation*}
Dengan demikian, barisan jumlah parsial terbatas dan deret tersebut konvergen.
\end{proof}

Baik uji deret-$p$ maupun uji perbandingan tidak
memberi tahu kita nilai limit jumlahnya.  Keduanya hanya memberi tahu bahwa limit
jumlah-jumlah parsial ada.  Sebagai contoh, meskipun kita mengetahui bahwa
$\sum_{n=1}^\infty \nicefrac{1}{n^2}$ konvergen, jauh lebih sulit untuk
menemukan\footnote{Pembuktian fakta inilah yang
membuat matematikawan Swiss
\href{https://en.wikipedia.org/wiki/Leonhard_Euler}{Leonhard Euler}
(1707--1783)
terkenal.}
bahwa limitnya adalah $\nicefrac{\pi^2}{6}$.
Jika kita memperlakukan $\sum_{n=1}^\infty \nicefrac{1}{n^p}$ sebagai fungsi dari $p$,
kita memperoleh fungsi $\zeta$ (zeta) Riemann.  Kajian mengenai
perilaku fungsi ini mencakup
salah satu masalah tak terpecahkan paling terkenal dalam matematika saat ini dan memiliki penerapan
dalam bidang yang tampaknya tidak berkaitan, seperti kriptografi modern.

\begin{example}
Deret $\sum_{n=1}^\infty \frac{1}{n^2+1}$ konvergen.

Bukti: Pertama, $\frac{1}{n^2+1} < \frac{1}{n^2}$ untuk semua $n \in \N$.
Deret $\sum_{n=1}^\infty \frac{1}{n^2}$ konvergen berdasarkan uji deret-$p$.
Oleh karena itu, berdasarkan uji perbandingan, $\sum_{n=1}^\infty \frac{1}{n^2+1}$ konvergen.
\end{example}

\subsection{Uji rasio}

Misalkan $r > 0$.  Rasio dua suku berurutan dalam deret geometri
$\sum_{n=0}^\infty r^n$ adalah $\frac{r^{n+1}}{r^n} = r$, dan deret tersebut konvergen
setiap kali $r < 1$.  Seperti pada barisan, fakta ini
dapat diperumum untuk deret yang lebih umum
asalkan rasio semacam itu ada \myquote{dalam limit.}  Kemudian kita membandingkan
ekor suatu deret dengan deret geometri.


\begin{prop}[Uji rasio]\index{uji rasio untuk deret}
Misalkan $\sum_{n=1}^\infty x_n$ adalah suatu deret dengan $x_n \neq 0$ untuk semua $n$, dan misalkan
\begin{equation*}
L \coloneqq \lim_{n\to\infty} \frac{\sabs{x_{n+1}}}{\sabs{x_n}}
\qquad \text{ada.}
\end{equation*}
\begin{enumerate}[(i)]
\item
Jika $L < 1$, maka $\sum_{n=1}^\infty x_n$ konvergen mutlak.
\item
Jika $L > 1$, maka $\sum_{n=1}^\infty x_n$ divergen.
\end{enumerate}
\end{prop}

Meskipun uji sebagaimana dinyatakan di atas sering kali memadai, uji tersebut dapat sedikit
diperkuat; lihat \exerciseref{exercise:strongerratiotest}.

\begin{proof}
Jika $L > 1$, maka
\lemmaref{seq:ratiotest} menyatakan bahwa barisan $\{ x_n \}_{n=1}^\infty$
divergen.  Karena syarat perlu bagi kekonvergenan deret adalah
suku-sukunya menuju nol, kita mengetahui bahwa $\sum_{n=1}^\infty x_n$ harus divergen.

Jadi, misalkan $L < 1$.
Kita akan menunjukkan bahwa $\sum_{n=1}^\infty \sabs{x_n}$ harus konvergen.
Buktinya serupa dengan bukti \lemmaref{seq:ratiotest}.  Tentu saja $L \geq
0$.  
Pilih
$r$ sedemikian sehingga $L < r < 1$.  Karena $r-L > 0$, terdapat $M \in \N$ sedemikian sehingga untuk
semua $n \geq M$,
\begin{equation*}
\abs{\frac{\sabs{x_{n+1}}}{\sabs{x_n}} - L} < r-L .
\end{equation*}
Oleh karena itu,
\begin{equation*}
\frac{\sabs{x_{n+1}}}{\sabs{x_n}} < r .
\end{equation*}
Untuk $n > M$ (yaitu untuk $n \geq M+1$),
tuliskan
\begin{equation*}
\sabs{x_n} =
\sabs{x_M}
\frac{\sabs{x_{M+1}}}{\sabs{x_{M}}}
\frac{\sabs{x_{M+2}}}{\sabs{x_{M+1}}}
\cdots
\frac{\sabs{x_{n}}}{\sabs{x_{n-1}}}
<
\sabs{x_M}
r r \cdots r = \sabs{x_M} r^{n-M} =
\bigl(\sabs{x_M} r^{-M}\bigr) r^n .
\end{equation*}
Untuk $k > M$, tuliskan jumlah parsial sebagai
\begin{equation*}
\begin{split}
\sum_{n=1}^k \sabs{x_n}
& =
\left(\sum_{n=1}^{M} \sabs{x_n} \right)
+
\left(\sum_{n=M+1}^{k} \sabs{x_n} \right)
\\
& <
\left(\sum_{n=1}^{M} \sabs{x_n} \right)
+
\left(\sum_{n=M+1}^{k} 
\bigl(\sabs{x_M} r^{-M}\bigr) r^n
\right)
\\
& =
\left(\sum_{n=1}^{M} \sabs{x_n} \right)
+
\bigl(\sabs{x_M} r^{-M}\bigr)
\left( \sum_{n=M+1}^{k} r^n \right) .
\end{split}
\end{equation*}
Karena $0 < r < 1$, deret geometri
$\sum_{n=0}^{\infty} r^n$ konvergen, sehingga
$\sum_{n=M+1}^{\infty} r^n$ juga konvergen.  Kita mengambil
limit ketika $k$ menuju tak hingga pada ruas kanan di atas dan memperoleh
\begin{equation*}
\begin{split}
\sum_{n=1}^k \sabs{x_n}
& <
\left(\sum_{n=1}^{M} \sabs{x_n} \right)
+
\bigl(\sabs{x_M} r^{-M}\bigr)
\left( \sum_{n=M+1}^{k} r^n \right) 
\\
& \leq
\left(\sum_{n=1}^{M} \sabs{x_n} \right)
+
\bigl(\sabs{x_M} r^{-M}\bigr)
\left( \sum_{n=M+1}^{\infty} r^n \right) .
\end{split}
\end{equation*}
Ruas kanan merupakan bilangan yang tidak bergantung pada $k$.
Dengan demikian, barisan jumlah parsial $\sum_{n=1}^\infty \sabs{x_n}$ terbatas
dan $\sum_{n=1}^\infty \sabs{x_n}$ konvergen.  Jadi, $\sum_{n=1}^\infty x_n$
konvergen mutlak.
\end{proof}

\begin{example}
Deret
\begin{equation*}
\sum_{n=1}^\infty \frac{2^n}{n!}
\end{equation*}
konvergen mutlak.

Bukti: Kita tuliskan
\begin{equation*}
\lim_{n\to\infty} \frac{2^{(n+1)}/(n+1)!}{2^n / n!} =
\lim_{n\to\infty} \frac{2}{n+1} = 0 .
\end{equation*}
Oleh karena itu, deret tersebut konvergen mutlak berdasarkan uji rasio.
\end{example}

\subsection{Latihan}

\begin{exercise}
Misalkan jumlah parsial ke-$k$ dari $\displaystyle \sum_{n=1}^\infty x_n$ adalah $s_k = \frac{k}{k+1}$.
Tentukan deret tersebut (yaitu tentukan $x_n$), buktikan bahwa deret tersebut konvergen, lalu
tentukan limitnya.
\end{exercise}

\begin{exercise} \label{geometric:exr}
Buktikan \propref{geometric:prop}; yaitu, untuk $-1 < r < 1$, buktikan
\begin{equation*}
\sum_{n=0}^\infty r^n = \frac{1}{1-r} .
\end{equation*}
Petunjuk: Lihat \exampleref{example:geometricsum}.
\end{exercise}

\begin{exercise}
Tentukan kekonvergenan atau kedivergenan deret-deret berikut.

\medskip

\noindent
a)
$\displaystyle \sum_{n=1}^\infty \frac{3}{9n+1}$
\qquad
b)
$\displaystyle \sum_{n=1}^\infty \frac{1}{2n-1}$
\qquad
c)
$\displaystyle \sum_{n=1}^\infty \frac{{(-1)}^n}{n^2}$
\qquad
d)
$\displaystyle \sum_{n=1}^\infty \frac{1}{n(n+1)}$
\qquad
e)
$\displaystyle \sum_{n=1}^\infty n e^{-n^2}$
\end{exercise}

\begin{samepage}
\begin{exercise}
\leavevmode
\begin{enumerate}[a)]
\item Buktikan bahwa jika
$\displaystyle
\sum_{n=1}^\infty x_n
$
konvergen, maka
$\displaystyle
\sum_{n=1}^\infty ( x_{2n} + x_{2n+1} )
$
deret tersebut juga konvergen.
\item
Berikan contoh eksplisit ketika kebalikannya tidak berlaku.
\end{enumerate}
\end{exercise}
\end{samepage}

\begin{exercise}
Untuk $i=1,2,\ldots,n$, misalkan $\{ x_{i,k} \}_{k=1}^\infty$ menyatakan $n$
barisan.  Misalkan bahwa untuk setiap $i$,
\begin{equation*}
\sum_{k=1}^\infty x_{i,k}
\end{equation*}
konvergen.  Buktikan
\begin{equation*}
\sum_{i=1}^n \left( \sum_{k=1}^\infty x_{i,k} \right)
=
\sum_{k=1}^\infty \left( \sum_{i=1}^n x_{i,k} \right) .
\end{equation*}
\end{exercise}

\begin{exercise} \label{exercise:strongerratiotest}
Buktikan versi uji rasio yang lebih kuat berikut:
Misalkan $\sum_{n=1}^\infty x_n$ adalah deret.
\begin{enumerate}[a)]
\item
Jika terdapat $N$ dan $\rho < 1$ sedemikian sehingga
$\frac{\sabs{x_{n+1}}}{\sabs{x_n}} < \rho$ untuk semua $n \geq N$,
maka deret tersebut konvergen mutlak.
(Catatan: Secara ekuivalen, syarat ini dapat dinyatakan sebagai
$\limsup\limits_{n\to\infty} \frac{\sabs{x_{n+1}}}{\sabs{x_n}} < 1$.)
\item
Jika terdapat $N$ sedemikian sehingga
$\frac{\sabs{x_{n+1}}}{\sabs{x_n}} \geq 1$
untuk semua $n \geq N$,
maka deret tersebut divergen. 
\end{enumerate}
\end{exercise}

\begin{exercise}[Menantang]
Misalkan $\{ x_n \}_{n=1}^\infty$ adalah barisan menurun dan
$\sum_{n=1}^\infty x_n$ konvergen.
Buktikan $\displaystyle \lim_{n\to\infty} n x_n = 0$.
\end{exercise}

\begin{exercise}
Tunjukkan bahwa $\displaystyle \sum_{n=1}^\infty \frac{{(-1)}^n}{n}$ konvergen.
Petunjuk: Tinjau jumlah dua suku berurutan.
\end{exercise}

\begin{exercise}
\leavevmode
\begin{enumerate}[a)]
\item Buktikan bahwa jika $\sum_{n=1}^\infty x_n$ dan $\sum_{n=1}^\infty y_n$ konvergen mutlak, maka
$\sum_{n=1}^\infty x_ny_n$ konvergen mutlak.
\item Berikan contoh eksplisit ketika kebalikannya tidak berlaku.
\item Berikan contoh eksplisit ketika ketiga deret tersebut konvergen mutlak,
bukan sekadar jumlah berhingga,
dan $\left(\sum_{n=1}^\infty x_n\right)\left(\sum_{n=1}^\infty y_n\right)
\neq \sum_{n=1}^\infty x_ny_n$.  Artinya, tunjukkan bahwa deret tidak
dikalikan suku demi suku.
\end{enumerate}
\end{exercise}

\begin{exercise}
Buktikan ketaksamaan segitiga untuk deret:
Jika $\sum_{n=1}^\infty x_n$ konvergen mutlak, maka
\begin{equation*}
\abs{\sum_{n=1}^\infty x_n} \leq
\sum_{n=1}^\infty \sabs{x_n} .
\end{equation*}
\end{exercise}

\begin{exercise}
Buktikan \emph{\myindex{uji perbandingan limit}}.  Artinya, buktikan bahwa jika
$a_n > 0$ dan $b_n > 0$ untuk semua $n$, dan
\begin{equation*}
0 < \lim_{n\to\infty} \frac{a_n}{b_n} < \infty ,
\end{equation*}
maka $\sum_{n=1}^\infty a_n$ dan $\sum_{n=1}^\infty b_n$ keduanya konvergen atau keduanya divergen.
\end{exercise}

\begin{exercise} \label{exercise:badnocauchy}
Misalkan $x_n \coloneqq \sum_{i=1}^n \nicefrac{1}{i}$.  Tunjukkan bahwa untuk setiap $k$,
kita memperoleh
$\displaystyle \lim_{n\to\infty} \sabs{x_{n+k}-x_n} = 0$, namun
$\{ x_n \}_{n=1}^\infty$ bukan barisan Cauchy.
\end{exercise}

\begin{samepage}
\begin{exercise}
Misalkan $s_k$ adalah jumlah parsial ke-$k$ dari $\sum_{n=1}^\infty x_n$.
\begin{enumerate}[a)]
\item
Misalkan terdapat $m \in \N$ sedemikian sehingga $\lim\limits_{k\to\infty}
s_{mk}$ ada dan $\lim\limits_{n\to\infty} x_n = 0$.  Tunjukkan bahwa
$\sum_{n=1}^\infty x_n$ konvergen.
\item
Berikan contoh ketika $\lim\limits_{k\to\infty} s_{2k}$ ada dan
$\lim\limits_{n\to\infty} x_n \neq 0$ (sehingga $\sum_{n=1}^\infty x_n$ divergen).
\item
(Menantang) Berikan contoh ketika $\lim_{n\to\infty} x_n = 0$, dan terdapat
subbarisan $\{ s_{k_i} \}_{i=1}^\infty$ sedemikian sehingga $\displaystyle \lim_{i\to\infty} s_{k_i}$ ada,
tetapi $\sum_{n=1}^\infty x_n$ tetap divergen.
\end{enumerate}
\end{exercise}
\end{samepage}

\begin{exercise} \label{exercise:squareseriesconv}
Misalkan $\sum_{n=1}^\infty x_n$ konvergen dan $x_n \geq 0$ untuk semua $n$.
Buktikan bahwa $\sum_{n=1}^\infty x_n^2$ konvergen.
\end{exercise}

\begin{exercise}[Menantang] \label{exercise:cauchycondensation}
Misalkan $\{ x_n\}_{n=1}^\infty$ adalah barisan menurun yang terdiri atas bilangan positif.
Bukti kekonvergenan/kedivergenan deret-$p$ dapat diperumum.
Buktikan
\emph{\myindex{prinsip kondensasi Cauchy}}:
\begin{equation*}
\sum_{n=1}^\infty x_n
\qquad \text{konvergen jika dan hanya jika} \qquad
\sum_{n=1}^\infty 2^n x_{2^n} \qquad \text{konvergen}.
\end{equation*}
\end{exercise}

\begin{exercise}
Gunakan prinsip kondensasi Cauchy
(lihat \exerciseref{exercise:cauchycondensation})
untuk menentukan kekonvergenan deret-deret berikut:

\medskip

\noindent
a) $\displaystyle \sum_{n=1}^\infty \frac{\ln n}{n^2}$
\qquad
b) $\displaystyle \sum_{n=2}^\infty \frac{1}{n \ln n}$
\qquad
c) $\displaystyle \sum_{n=2}^\infty \frac{1}{n {(\ln n)}^2}$
\qquad
d) $\displaystyle \sum_{n=2}^\infty \frac{1}{n (\ln n ){(\ln \ln n)}^2}$

\medskip

\noindent
Perhatikan bahwa hanya ekor beberapa deret ini yang memenuhi
hipotesis prinsip tersebut; jelaskan mengapa hal itu sudah memadai.

\noindent
Petunjuk: Anda boleh menggunakan identitas $\ln (2^n) = n \ln 2$.
\end{exercise}

\begin{exercise}[Menantang]
Buktikan \emph{\myindex{teorema Abel}}:

\medskip

\noindent
\emph{\textbf{Teorema.} Misalkan $\sum_{n=1}^\infty x_n$ adalah deret yang barisan jumlah parsialnya
terbatas, $\{ \lambda_n \}_{n=1}^\infty$ adalah barisan dengan
$\lim_{n\to\infty} \lambda_n = 0$, dan
$\sum_{n=1}^\infty \sabs{ \lambda_{n+1} - \lambda_n }$ konvergen.
Maka $\sum_{n=1}^\infty \lambda_n x_n$ konvergen.}
\end{exercise}



%%%%%%%%%%%%%%%%%%%%%%%%%%%%%%%%%%%%%%%%%%%%%%%%%%%%%%%%%%%%%%%%%%%%%%%%%%%%%%

\sectionnewpage
\section{Lebih lanjut tentang deret}
\label{sec:moreonseries}

%mbxINTROSUBSECTION

\sectionnotes{hingga 2--3 kuliah (opsional, dapat dilewati dengan aman atau dibahas sebagian)}

\subsection{Uji akar}

Uji yang serupa dengan uji rasio adalah
\emph{\myindex{uji akar}}.  Bukti uji ini pun serupa.
Sekali lagi, gagasannya adalah memperumum apa yang terjadi pada deret geometri.

\begin{prop}[Uji akar]
Misalkan $\sum_{n=1}^\infty x_n$ adalah deret dan misalkan
\begin{equation*}
L \coloneqq \limsup_{n\to\infty} {\sabs{x_n}}^{1/n} .
\end{equation*}
\begin{enumerate}[(i)]
\item Jika $L < 1$, maka $\sum_{n=1}^\infty x_n$ konvergen mutlak.
\item Jika $L > 1$, maka $\sum_{n=1}^\infty x_n$ divergen.
\end{enumerate}
\end{prop}

\begin{proof}
Jika $L > 1$, maka terdapat\footnote{%
Untuk kasus $L=\infty$, lihat \exerciseref{exercise:extendsubseqlimsupinf}.
Sebagai alternatif, perhatikan bahwa jika $L=\infty$, maka $\bigl\{ {\sabs{x_{n}}}^{1/n}
\bigr\}_{n=1}^\infty$
tak terbatas, sehingga
$\{ x_n \}_{n=1}^\infty$ juga tak terbatas.}
subbarisan $\{ x_{n_k} \}_{k=1}^\infty$ sedemikian sehingga
$L = \lim_{k\to\infty} {\sabs{x_{n_k}}}^{1/n_k}$.  Misalkan
$r$ sedemikian sehingga $L > r > 1$.  Terdapat $M$ sedemikian
sehingga untuk semua $k \geq M$, berlaku 
${\sabs{x_{n_k}}}^{1/n_k} > r > 1$, atau dengan kata lain,
$\sabs{x_{n_k}} > r^{n_k} > 1$.
Subbarisan 
$\bigl\{ \sabs{x_{n_k}} \bigr\}_{k=1}^\infty$,
dan karena itu juga
$\bigl\{ \sabs{x_{n}} \bigr\}_{n=1}^\infty$,
tidak mungkin konvergen ke nol, sehingga deret tersebut divergen.

Sekarang, misalkan $L < 1$.  Pilih $r$ sedemikian sehingga $L < r < 1$.
Berdasarkan definisi limit superior,
terdapat $M$ sedemikian sehingga untuk semua $n \geq M$,
\begin{equation*}
\sup \bigl\{ {\sabs{x_k}}^{1/k} : k \geq n \bigr\} < r .
\end{equation*}
Oleh karena itu, untuk semua $n \geq M$,
\begin{equation*}
{\sabs{x_n}}^{1/n} < r , \qquad \text{atau dengan kata lain} \qquad \sabs{x_n} < r^n .
\end{equation*}
Misalkan $k > M$, lalu taksir jumlah parsial ke-$k$:
\begin{equation*}
\sum_{n=1}^k \sabs{x_n} = 
\left( \sum_{n=1}^M \sabs{x_n} \right) + 
\left( \sum_{n=M+1}^k \sabs{x_n} \right)
\leq
\left( \sum_{n=1}^M \sabs{x_n} \right) + 
\left( \sum_{n=M+1}^k r^n \right) .
\end{equation*}
Karena $0 < r < 1$,
deret geometri $\sum_{n=M+1}^\infty r^n$ konvergen ke
$\frac{r^{M+1}}{1-r}$.  Karena semuanya positif,
\begin{equation*}
\sum_{n=1}^k \sabs{x_n} 
\leq
\left( \sum_{n=1}^M \sabs{x_n} \right) + 
\frac{r^{M+1}}{1-r} .
\end{equation*}
Dengan demikian, barisan jumlah parsial dari deret $\sum_{n=1}^\infty \sabs{x_n}$ terbatas, dan
deret tersebut konvergen.  Oleh karena itu, $\sum_{n=1}^\infty x_n$ konvergen mutlak.
\end{proof}

\subsection{Uji deret berselang-seling}

Uji-uji yang telah kita lihat sejauh ini hanya membahas kekonvergenan mutlak.  Uji
berikut menghasilkan banyak deret yang konvergen bersyarat.

\begin{prop}[Deret berselang-seling]
Misalkan $\{ x_n \}_{n=1}^\infty$ adalah barisan bilangan real positif yang monoton turun sedemikian
sehingga $\lim\limits_{n\to\infty} x_n = 0$.  Maka
\begin{equation*}
\sum_{n=1}^\infty {(-1)}^n x_n \qquad \text{konvergen.}
\end{equation*}
\end{prop}

\begin{proof}
Misalkan $s_m \coloneqq \sum_{n=1}^m {(-1)}^n x_n$ adalah jumlah parsial ke-$m$.  Maka
\begin{equation*}
s_{2k} =
\sum_{n=1}^{2k} {(-1)}^n x_n
=
(-x_1 + x_2) + \cdots + (-x_{2k-1} + x_{2k})
=
\sum_{\ell=1}^{k} (-x_{2\ell-1} + x_{2\ell}) .
\end{equation*}
Barisan $\{ x_n \}_{n=1}^\infty$ menurun, sehingga $(-x_{2\ell-1}+x_{2\ell}) \leq 0$
untuk semua $\ell$.
Dengan demikian, subbarisan $\{ s_{2k} \}_{k=1}^\infty$ dari jumlah-jumlah parsial
merupakan barisan menurun.  Demikian pula, $(x_{2\ell}-x_{2\ell+1}) \geq 0$, sehingga
\begin{equation*}
s_{2k} = - x_1 + ( x_2 - x_3 ) + \cdots + ( x_{2k-2} - x_{2k-1} ) + x_{2k}
\geq -x_1 .
\end{equation*}
Intuisi di balik batas
$0 \geq s_{2k} \geq -x_1$ diilustrasikan dalam \figureref{fig:seralternate}.
\begin{myfigureht}
\myincludegraphics{ser-alternate}{%
Sebuah diagram pada bidang yang menampilkan nilai-nilai barisan plus-minus x sub n
untuk n dari 1 hingga 8.  Dengan garis penuh yang dimulai dari sumbu horizontal,
tampak sebuah "grafik" minus x sub 1, x sub 2, minus x sub 3, x sub 4, dan seterusnya.
Terlihat bahwa setiap garis penuh berikutnya lebih pendek
karena setiap x sub n berikutnya lebih kecil daripada sebelumnya.  Dengan garis putus-putus,
juga ditampilkan minus x sub 2, minus x sub 4, dan seterusnya.
Selisih
minus x sub 1 ditambah x sub 2,
minus x sub 3 ditambah x sub 4,
minus x sub 5 ditambah x sub 6,
dan
minus x sub 7 ditambah x sub 8 ditampilkan dan merupakan selisih
antara minus x sub 1 dan minus x sub 2 yang bersesuaian.
Perhatikan bahwa selisih-selisih ini hanya mengisi sebagian panjang total
garis untuk minus x sub 1.}
\caption{Menunjukkan bahwa $0 \geq s_{2k} \geq -x_1$ dengan $k=4$ untuk suatu deret berselang-seling.\label{fig:seralternate}}
\end{myfigureht}

Karena $\{ s_{2k} \}_{k=1}^\infty$ menurun dan terbatas di bawah, barisan tersebut konvergen.
Misalkan $a \coloneqq \lim_{k\to\infty} s_{2k}$.
Kita hendak menunjukkan bahwa $\lim_{m\to\infty} s_m = a$ (dan bukan hanya untuk subbarisan tersebut).
Untuk $\epsilon > 0$, pilih $M$ sedemikian sehingga $\sabs{s_{2k}-a} <
\nicefrac{\epsilon}{2}$ setiap kali $k \geq M$.
Karena $\lim_{n\to\infty} x_n = 0$, kita juga dapat
memperbesar $M$ jika perlu
untuk memperoleh
$x_{2k+1} < \nicefrac{\epsilon}{2}$ setiap kali $k \geq M$.  
Misalkan $m \geq 2M+1$.  Jika $m=2k$, maka $k \geq M+\nicefrac{1}{2} \geq M$ dan
$\sabs{s_{m}-a}=\sabs{s_{2k}-a} < \nicefrac{\epsilon}{2} < \epsilon$.
Jika $m=2k+1$, maka $k \geq M$ juga.  Perhatikan bahwa
$s_{2k+1} = s_{2k} - x_{2k+1}$.
Jadi
\begin{equation*}
\sabs{s_{m}-a} = 
\sabs{s_{2k+1}-a} = 
\sabs{s_{2k}-a - x_{2k+1}} \leq
\sabs{s_{2k}-a} + x_{2k+1} < 
\nicefrac{\epsilon}{2}+ \nicefrac{\epsilon}{2} = \epsilon .  \qedhere
\end{equation*}
\end{proof}

Patut diperhatikan bahwa terdapat deret yang konvergen bersyarat
dengan nilai mutlak suku-sukunya menuju nol selambat apa pun.
Deret
\begin{equation*}
\sum_{n=1}^\infty \frac{{(-1)}^n}{n^p}
\end{equation*}
konvergen untuk $p > 0$ yang sekecil apa pun, tetapi tidak konvergen
mutlak ketika $p \leq 1$.

\subsection{Penyusunan ulang}

Deret yang konvergen mutlak berperilaku seperti yang kita harapkan.  Sebagai contoh,
deret yang konvergen mutlak dapat dijumlahkan dalam urutan apa pun.  Hal semacam ini
tidak berlaku untuk deret yang konvergen bersyarat
(lihat \exampleref{example:harmonsumanything}
dan \exerciseref{exercise:seriesconvergestoanything}).

Perhatikan suatu deret
\begin{equation*}
\sum_{n=1}^\infty x_n .
\end{equation*}
Diberikan fungsi bijektif $\sigma \colon \N \to \N$, penyusunan
ulang\index{penyusunan ulang deret} yang bersesuaian adalah
deret:
\begin{equation*}
\sum_{k=1}^\infty x_{\sigma(k)} .
\end{equation*}
Kita cukup menjumlahkan deret tersebut dalam urutan yang berbeda.

\begin{prop}
Misalkan $\sum_{n=1}^\infty x_n$ merupakan deret konvergen mutlak yang konvergen ke suatu bilangan
$x$.  Misalkan $\sigma \colon \N \to \N$ merupakan suatu bijeksi.  Maka
$\sum_{n=1}^\infty x_{\sigma(n)}$ konvergen mutlak dan konvergen ke $x$.
\end{prop}

Dengan kata lain,
setiap penyusunan ulang deret yang konvergen mutlak juga konvergen (secara mutlak)
ke bilangan yang sama.

\begin{proof}
Ambil sebarang $\epsilon > 0$.  Karena $\sum_{n=1}^\infty x_n$ konvergen mutlak, ambil $M$ sedemikian sehingga
\begin{equation*}
\abs{\left(\sum_{n=1}^M x_n \right) - x} < \frac{\epsilon}{2}
\qquad \text{dan} \qquad
\sum_{n=M+1}^\infty \sabs{x_n} < \frac{\epsilon}{2} .
\end{equation*}
Karena $\sigma$ merupakan bijeksi,
terdapat bilangan $K$ sedemikian sehingga untuk setiap
$n \leq M$, terdapat $k \leq K$ sedemikian sehingga $\sigma(k) = n$.
Dengan kata lain,
$\{ 1,2,\ldots,M \} \subset \sigma\bigl(\{ 1,2,\ldots,K \} \bigr)$.

Untuk $N \geq K$, misalkan $Q \coloneqq \max \sigma\bigl(\{ 1,2,\ldots,N \}\bigr)$.
Hitung
\begin{equation*}
\begin{split}
\abs{\left( \sum_{n=1}^N x_{\sigma(n)} \right) - x}
& =
\abs{ \left( \sum_{n=1}^M x_n
+
\sum_{\substack{n=1\\\sigma(n) > M}}^N x_{\sigma(n)} \right) - x}
\\
& \leq
\abs{ \left( \sum_{n=1}^M x_n \right) - x}
+
\sum_{\substack{n=1\\\sigma(n) > M}}^N \sabs{x_{\sigma(n)}}
\\
& \leq
\abs{ \left( \sum_{n=1}^M x_n \right) - x}
+
\sum_{n=M+1}^Q \sabs{x_{n}}
\\
& < \nicefrac{\epsilon}{2} + \nicefrac{\epsilon}{2} = \epsilon .
\end{split}
\end{equation*}
Jadi,
$\sum_{n=1}^\infty x_{\sigma(n)}$ konvergen ke $x$.  Untuk melihat bahwa kekonvergenannya
bersifat mutlak, kita terapkan argumen di atas pada $\sum_{n=1}^\infty \sabs{x_n}$ untuk menunjukkan
bahwa $\sum_{n=1}^\infty \sabs{x_{\sigma(n)}}$ konvergen.
\end{proof}

\begin{example} \label{example:harmonsumanything}
Mari kita tunjukkan bahwa deret harmonik berselang-seling
$\sum_{n=1}^\infty \frac{{(-1)}^{n+1}}{n}$, yang tidak konvergen mutlak, dapat
disusun ulang agar konvergen ke sebarang bilangan.
Suku-suku berindeks ganjil dan suku-suku berindeks genap masing-masing divergen ke
tak hingga positif dan tak hingga negatif (buktikanlah!):
\begin{equation*}
\sum_{m=1}^\infty \frac{1}{2m-1} = \infty, \qquad \text{dan} \qquad
\sum_{m=1}^\infty \frac{-1}{2m} = -\infty .
\end{equation*}
Untuk menyederhanakan penulisan, misalkan $a_n \coloneqq \frac{{(-1)}^{n+1}}{n}$,
ambil sebarang bilangan $L \in \R$, dan tetapkan $\sigma(1) \coloneqq 1$.
Andaikan kita telah
mendefinisikan $\sigma(n)$ untuk semua $n \leq N$.  Jika
\begin{equation*}
\sum_{n=1}^N a_{\sigma(n)} \leq L ,
\end{equation*}
maka tetapkan $\sigma(N+1) \coloneqq k$, dengan $k \in \N$ bilangan ganjil terkecil
yang belum kita gunakan,
yaitu $\sigma(n) \neq k$ untuk semua $n \leq N$.
Jika tidak, tetapkan $\sigma(N+1) \coloneqq k$, dengan $k$
bilangan genap terkecil yang belum kita gunakan.

Menurut konstruksi, $\sigma \colon \N \to \N$ bersifat satu-satu.
Fungsi ini juga surjektif, sebab jika kita terus menambahkan suku-suku berindeks ganjil (atau genap),
pada akhirnya jumlah tersebut akan melewati $L$ dan kita beralih
ke suku-suku berindeks genap (atau ganjil).  Jadi, kita beralih tak berhingga kali.

Terakhir, misalkan $N$ adalah indeks ketika kita baru saja melewati $L$ lalu beralih.
Andaikan kita baru saja beralih dari indeks ganjil ke indeks genap, sehingga kita mulai
mengurangkan.
Andaikan kita terus mengurangkan hingga suatu
$N' > N$ dan akan beralih kembali dari indeks genap ke indeks ganjil pada langkah berikutnya.
Maka
\begin{equation*}
L + \frac{1}{\sigma(N)} \geq \sum_{n=1}^{N} a_{\sigma(n)}
> \sum_{n=1}^{N'} a_{\sigma(n)} > L- \frac{1}{\sigma(N')}.
\end{equation*}
Hal yang serupa berlaku ketika kita beralih ke arah sebaliknya.
Oleh karena itu,
jumlah hingga $N'$ berjarak kurang dari $\frac{1}{\min \{ \sigma(N), \sigma(N') \}}$
terhadap $L$.  Karena
kita beralih tak berhingga kali, $\sigma(N) \to \infty$
dan $\sigma(N') \to \infty$. Dengan demikian,
\begin{equation*}
\sum_{n=1}^\infty a_{\sigma(n)} = 
\sum_{n=1}^\infty \frac{{(-1)}^{\sigma(n)+1}}{\sigma(n)} = L .
\end{equation*}

Berikut adalah contoh untuk mengilustrasikan pembuktian tersebut.  Ketika $L=1.2$, urutannya
adalah
\begin{equation*}
1+\nicefrac{1}{3}-\nicefrac{1}{2}+\nicefrac{1}{5}+\nicefrac{1}{7}+\nicefrac{1}{9}-\nicefrac{1}{4}+\nicefrac{1}{11}+\nicefrac{1}{13}-\nicefrac{1}{6}
+\nicefrac{1}{15}+\nicefrac{1}{17}+\nicefrac{1}{19} - \nicefrac{1}{8} + \cdots .
\end{equation*}
Pada tahap ini, jarak jumlah tersebut dari limit tidak lebih dari $\nicefrac{1}{8}$.
Lihat \figureref{fig:serrearrange}.
\begin{myfigureht}
\myincludegraphics{ser-rearrange}{%
Diagram suatu barisan yang digambarkan sebagai titik-titik.  Sebuah garis putus-putus horizontal ditampilkan.
Barisan tersebut bermula di bawah garis putus-putus, lalu berada di atasnya pada
langkah berikutnya.  Kemudian barisan itu kembali ke bawah garis putus-putus, lalu memerlukan beberapa langkah
sebelum berada di atas garis putus-putus, dan langkah berikutnya kembali berada di bawah
garis tersebut.  Demikian seterusnya.  Ukuran langkahnya semakin kecil.}
\caption{Empat belas jumlah parsial pertama dari penyusunan ulang yang konvergen
ke $1.2$.\label{fig:serrearrange}}
\end{myfigureht}
\end{example}

\subsection{Perkalian deret}

Seperti yang telah
kita sebutkan,
perkalian deret agak lebih sulit daripada penjumlahan.
Jika setidaknya salah satu deret konvergen
mutlak, kita dapat menggunakan teorema berikut.  Untuk hasil ini, lebih
mudah jika deret dimulai dari 0, bukan dari 1.

\begin{thm}[\myindex{teorema Mertens}\footnote{Dibuktikan oleh
matematikawan Jerman
\href{https://en.wikipedia.org/wiki/Franz_Mertens}{Franz Mertens}
(1840--1927).}]
Andaikan $\sum_{n=0}^\infty a_n$ dan $\sum_{n=0}^\infty b_n$ merupakan dua deret konvergen, yang masing-masing konvergen
ke $A$ dan $B$.  Andaikan setidaknya salah satu deret tersebut
konvergen mutlak.  Definisikan
\begin{equation*}
c_n \coloneqq a_0 b_n + a_1 b_{n-1} + \cdots + a_n b_0 = \sum_{i=0}^n a_i b_{n-i} .
\end{equation*}
Maka deret $\sum_{n=0}^\infty c_n$
konvergen ke $AB$.
\end{thm}

Deret $\sum_{n=0}^\infty c_n$ disebut \emph{\myindex{hasil kali Cauchy}} dari
$\sum_{n=0}^\infty a_n$ dan $\sum_{n=0}^\infty b_n$.

\begin{proof}
Andaikan $\sum_{n=0}^\infty a_n$ konvergen mutlak, dan ambil sebarang
$\epsilon > 0$.
Dalam pembuktian ini, daripada memilih taksiran yang rumit hanya agar
taksiran akhirnya kurang dari $\epsilon$,
mari kita peroleh suatu taksiran yang bergantung pada $\epsilon$
dan dapat dibuat sekecil apa pun.

Tuliskan
\begin{equation*}
A_m \coloneqq \sum_{n=0}^m a_n , \qquad B_m \coloneqq \sum_{n=0}^m b_n .
\end{equation*}
Kita susun ulang jumlah parsial ke-$m$ dari $\sum_{n=0}^\infty c_n$:
\begin{equation*}
\begin{split}
\abs{\left(\sum_{n=0}^m c_n \right) - AB}
& =
\abs{\left( \sum_{n=0}^m \sum_{i=0}^n a_i b_{n-i} \right) - AB}
\\
& =
\abs{\left( \sum_{n=0}^m
  B_n a_{m-n} \right) - AB}
\\
& =
\abs{\left( \sum_{n=0}^m
  ( B_n -  B ) a_{m-n} \right)
    + B A_m - AB}
\\
& \leq
\left(
\sum_{n=0}^m
  \sabs{ B_n -  B } \sabs{a_{m-n}}
\right)
+
\sabs{B}\sabs{A_m - A}
\end{split}
\end{equation*}
Suku kedua di ruas kanan tentu dapat kita buat kecil.
Kuncinya adalah menangani suku pertama.
Pilih $K$ sedemikian sehingga untuk semua $m \geq K$, berlaku
$\sabs{A_m - A} < \epsilon$ dan
juga
$\sabs{B_m - B} < \epsilon$.
Karena $\sum_{n=0}^\infty a_n$ konvergen mutlak,
terdapat $K$ yang cukup besar sedemikian sehingga
untuk semua $m \geq K$,
\begin{equation*}
\sum_{n=K}^m \sabs{a_n} < \epsilon .
\end{equation*}
Karena $\sum_{n=0}^\infty b_n$ konvergen,
$B_{\text{max}} \coloneqq \sup \bigl\{ \sabs{ B_n - B } : n = 0,1,2,\ldots
\bigr\}$
berhingga.  Ambil $m \geq 2K$.
Khususnya, $m-K+1 > K$.  Jadi,
\begin{equation*}
\begin{split}
\sum_{n=0}^m
  \sabs{ B_n -  B } \sabs{a_{m-n}}
& =
\left(
\sum_{n=0}^{m-K}
  \sabs{ B_n -  B } \sabs{a_{m-n}}
\right)
+
\left(
\sum_{n=m-K+1}^m
  \sabs{ B_n -  B } \sabs{a_{m-n}}
\right)
\\
& \leq
\left(
\sum_{n=K}^m
\sabs{a_{n}}
\right)
B_{\text{max}}
+
\left(
\sum_{n=0}^{K-1}
  \epsilon \sabs{a_{n}}
\right)
\\
& \leq
\epsilon
B_{\text{max}}
+
\epsilon
\left(
\sum_{n=0}^\infty \sabs{a_{n}}
\right) .
\end{split}
\end{equation*}
Oleh karena itu, untuk $m \geq 2K$, kita peroleh
\begin{equation*}
\begin{split}
\abs{\left(\sum_{n=0}^m c_n \right) - AB}
& \leq
\left(
\sum_{n=0}^m
  \sabs{ B_n -  B } \sabs{a_{m-n}}
\right)
+
\sabs{B}\sabs{A_m - A}
\\
& \leq
\epsilon
B_{\text{max}}
+
\epsilon
\left(
\sum_{n=0}^\infty \sabs{a_{n}}
\right)
+
\sabs{B}\epsilon
=
\epsilon 
\left(
B_{\text{max}}
+
\left(
\sum_{n=0}^\infty \sabs{a_{n}}
\right)
+
\sabs{B}
\right) .
\end{split}
\end{equation*}
Ekspresi di dalam tanda kurung pada ruas kanan
merupakan bilangan tetap.
Jadi,
kita dapat membuat ruas kanan sekecil apa pun dengan memilih
$\epsilon> 0$ yang cukup kecil.  Maka $\sum_{n=0}^\infty c_n$ konvergen ke $AB$.
\end{proof}

\begin{example}
Jika kedua deret hanya konvergen bersyarat, deret hasil kali Cauchy
bahkan belum tentu konvergen.
Andaikan kita mengambil $a_n = b_n = {(-1)}^n \frac{1}{\sqrt{n+1}}$.
Deret $\sum_{n=0}^\infty a_n = \sum_{n=0}^\infty b_n$
konvergen
menurut uji deret berselang-seling; namun, deret tersebut tidak konvergen
mutlak, seperti dapat dilihat dari uji-$p$.  Mari kita perhatikan
hasil kali Cauchy-nya.
\begin{equation*}
%mbxlatex \begin{aligned}
c_n
%mbxlatex &
= 
{(-1)}^n
\left(
\frac{1}{\sqrt{n+1}} + 
\frac{1}{\sqrt{2n}} + 
\frac{1}{\sqrt{3(n-1)}} + \cdots +
%\frac{1}{\sqrt{2n}} + 
\frac{1}{\sqrt{n+1}}
\right)
%mbxlatex \\
%mbxlatex &
=
{(-1)}^n
\sum_{i=0}^n \frac{1}{\sqrt{(i+1)(n-i+1)}} .
%mbxlatex \end{aligned}
\end{equation*}
Oleh karena itu,
\begin{equation*}
\sabs{c_n} 
=
\sum_{i=0}^n \frac{1}{\sqrt{(i+1)(n-i+1)}} 
\geq
\sum_{i=0}^n \frac{1}{\sqrt{(n+1)(n+1)}} 
= 1 .
\end{equation*}
Suku-sukunya tidak menuju nol sehingga $\sum_{n=0}^\infty c_n$ tidak mungkin konvergen.
\end{example}

\subsection{Deret pangkat}

Tetapkan $x_0 \in \R$.
Sebuah \emph{\myindex{deret pangkat}} di sekitar $x_0$
adalah deret berbentuk
\begin{equation*}
\sum_{n=0}^\infty a_n {(x-x_0)}^n .
\end{equation*}
Deret pangkat sesungguhnya merupakan fungsi dari $x$, dan
banyak fungsi penting dalam analisis dapat dituliskan
sebagai deret pangkat.  Kita menggunakan konvensi bahwa
$0^0 = 1$ (yakni jika $x=x_0$ dan $n=0$).

Suatu deret pangkat disebut
\emph{konvergen}\index{konvergen!deret pangkat} jika
terdapat setidaknya satu $x \neq x_0$ yang membuat deret tersebut konvergen.
Jika $x=x_0$, deret tersebut selalu
konvergen karena semua suku kecuali suku pertama bernilai nol.
Jika deret tersebut tidak konvergen pada titik mana pun $x \neq x_0$, kita menyebut
deret tersebut \emph{divergen}\index{divergen!deret pangkat}.

\begin{example} \label{ps:expex}
Deret
\begin{equation*}
\sum_{n=0}^\infty \frac{1}{n!} x^n
\end{equation*}
konvergen mutlak untuk semua $x \in \R$.  Untuk $x=0$, hal ini langsung berlaku;
untuk $x \neq 0$, uji rasio memberikan
\begin{equation*}
\lim_{n \to \infty}
\abs{\frac{\bigl(1/(n+1)!\bigr) \, x^{n+1}}{(1/n!) \, x^{n}}}
=
\lim_{n \to \infty}
\frac{\sabs{x}}{n+1}
=
0.
\end{equation*}
Ingat dari kalkulus bahwa deret ini konvergen ke $e^x$.
\end{example}

\begin{example} \label{ps:1kex}
Deret
\begin{equation*}
\sum_{n=1}^\infty \frac{1}{n} x^n
\end{equation*}
konvergen mutlak untuk semua $x \in (-1,1)$.  Untuk $x=0$, hal ini langsung berlaku; untuk $0 < \sabs{x} < 1$, uji rasio memberikan:
\begin{equation*}
\lim_{n \to \infty}
\abs{
\frac{\bigl(1/(n+1) \bigr) \, x^{n+1}}{(1/n) \, x^{n}}
}
=
\lim_{n \to \infty}
\sabs{x} \frac{n}{n+1}
=
\sabs{x} < 1 .
\end{equation*}
Deret tersebut konvergen pada $x=-1$ karena
$\sum_{n=1}^\infty \frac{{(-1)}^n}{n}$ konvergen
berdasarkan uji deret
berselang-seling.
Akan tetapi, deret pangkat tersebut tidak konvergen mutlak pada $x=-1$ karena
$\sum_{n=1}^\infty \frac{1}{n}$ tidak konvergen.
Deret tersebut
divergen pada $x=1$.
Jika $\sabs{x} > 1$, deret tersebut divergen berdasarkan uji rasio.
\end{example}

\begin{example} \label{ps:divergeex}
Deret
\begin{equation*}
\sum_{n=1}^\infty n^n x^n
\end{equation*}
divergen untuk semua $x \neq 0$.  Mari kita terapkan uji akar:
\begin{equation*}
\limsup_{n\to\infty}
\,
\sabs{n^n x^n}^{1/n}
=
\limsup_{n\to\infty}
\,
n \sabs{x}
= \infty .
\end{equation*}
Oleh karena itu, deret tersebut divergen untuk semua $x \neq 0$.
\end{example}

Secara umum, perilaku konvergensi deret pangkat serupa dengan
ketiga contoh di atas.

\begin{prop} \label{prop:powerserrealradius}
Misalkan $\sum_{n=0}^\infty a_n {(x-x_0)}^n$ suatu deret pangkat.
Jika deret tersebut konvergen, maka deret itu konvergen mutlak pada
semua $x \in \R$, atau
terdapat bilangan $\rho$ sedemikian sehingga
deret tersebut konvergen mutlak pada interval
$(x_0-\rho,x_0+\rho)$ dan divergen jika $x < x_0-\rho$ atau $x > x_0+\rho$.
\end{prop}

Bilangan $\rho$ disebut \emph{\myindex{jari-jari konvergensi}} dari
deret pangkat tersebut.  Kita menulis $\rho = \infty$ jika deret tersebut konvergen untuk
semua $x$, dan kita menulis $\rho = 0$ jika deret tersebut divergen.
Pada titik-titik ujung, yakni jika $x = x_0+\rho$ atau $x = x_0-\rho$,
proposisi tersebut tidak memberikan kesimpulan,
dan deretnya mungkin konvergen atau mungkin tidak.
Lihat \figureref{ps:convfig}.
Dalam \exampleref{ps:1kex},
jari-jari konvergensinya adalah $\rho=1$;
dalam \exampleref{ps:expex}, jari-jari konvergensinya adalah $\rho=\infty$;
dan dalam \exampleref{ps:divergeex}, jari-jari konvergensinya adalah $\rho=0$.

\begin{myfigureht}
\myincludepdft{ps-conv}{%
Diagram garis bilangan real dengan tiga titik yang ditandai, dari kiri ke
kanan, x subskrip 0 dikurangi rho, x subskrip 0, dan x subskrip 0 ditambah rho.
Bagian garis di sebelah kiri x subskrip 0 dikurangi rho dan
bagian garis di sebelah kanan x subskrip 0 ditambah rho diberi label
"divergen".  Interval antara x subskrip 0 dikurangi rho dan
x subskrip 0 ditambah rho diarsir dan diberi label "konvergen mutlak". 
Kedua titik x subskrip 0 dikurangi rho dan x subskrip 0 ditambah rho diberi label
tanda tanya.}
\caption{Kekonvergenan suatu deret pangkat.\label{ps:convfig}}
\end{myfigureht}

\begin{proof}
Tuliskan
\begin{equation*}
R \coloneqq \limsup_{n\to\infty} {\sabs{a_n}}^{1/n} .
\end{equation*}
Kita menerapkan uji akar,
\begin{equation*}
L = \limsup_{n\to\infty} {\babs{a_n {(x-x_0)}^n}}^{1/n} 
=
\sabs{x-x_0} \limsup_{n\to\infty} {\sabs{a_n}}^{1/n}
=
\sabs{x-x_0} R .
\end{equation*}
Jika $R = \infty$, maka $L=\infty$ untuk setiap $x \neq x_0$, dan
deret tersebut divergen berdasarkan uji akar.
Di sisi lain,
jika $R = 0$, maka $L=0$ untuk setiap $x$,
dan deret tersebut konvergen mutlak untuk semua $x$.

Misalkan $0 < R < \infty$.
Deret tersebut
konvergen mutlak jika
$1 > L = R \sabs{x-x_0}$,
yakni,
\begin{equation*}
\sabs{x-x_0} < \nicefrac{1}{R} .
\end{equation*}
Deret tersebut divergen jika
$1 < L = R \sabs{x-x_0}$,
atau
\begin{equation*}
\sabs{x-x_0} > \nicefrac{1}{R} .
\end{equation*}
Dengan menetapkan $\rho \coloneqq \nicefrac{1}{R}$, bukti selesai.
\end{proof}

Kiranya berguna untuk menyatakan kembali apa yang telah kita peroleh dalam bukti
sebagai proposisi tersendiri.

\begin{prop}
Misalkan $\sum_{n=0}^\infty a_n {(x-x_0)}^n$ suatu deret pangkat, dan misalkan
\begin{equation*}
R \coloneqq \limsup_{n\to\infty} {\sabs{a_n}}^{1/n} .
\end{equation*}
Jika $R = \infty$, deret pangkat tersebut divergen.  Jika
$R=0$, deret pangkat tersebut konvergen di mana-mana.   Dalam kasus lainnya,
jari-jari konvergensinya adalah $\rho = \nicefrac{1}{R}$.
\end{prop}

Sering kali, jari-jari konvergensi ditulis sebagai $\rho = \nicefrac{1}{R}$ dalam
ketiga kasus tersebut, dengan
kesepakatan mengenai nilai $\rho$ jika $R = 0$ atau $R =
\infty$.

\pagebreak[2]
Deret pangkat yang konvergen dapat dijumlahkan dan dikalikan satu sama lain, serta dikalikan
dengan konstanta.
Proposisi berikut memiliki bukti langsung berdasarkan apa yang kita ketahui tentang deret
secara umum dan deret pangkat secara khusus.  Kita serahkan buktinya kepada pembaca.

\begin{prop}
Misalkan $\sum_{n=0}^\infty a_n {(x-x_0)}^n$ dan
$\sum_{n=0}^\infty b_n {(x-x_0)}^n$ dua deret pangkat konvergen
dengan jari-jari konvergensi sekurang-kurangnya $\rho > 0$, dan misalkan $\alpha \in \R$.  Maka,
untuk semua $x$ sedemikian sehingga $\sabs{x-x_0} < \rho$, berlaku
\begin{equation*}
\left(\sum_{n=0}^\infty a_n {(x-x_0)}^n\right)
+
\left(\sum_{n=0}^\infty b_n {(x-x_0)}^n\right)
=
\sum_{n=0}^\infty (a_n+b_n) {(x-x_0)}^n ,
\end{equation*}
\begin{equation*}
\alpha
\left(\sum_{n=0}^\infty a_n {(x-x_0)}^n\right)
=
\sum_{n=0}^\infty \alpha a_n {(x-x_0)}^n ,
\end{equation*}
dan
\begin{equation*}
\left(\sum_{n=0}^\infty a_n {(x-x_0)}^n\right)
\,
\left(\sum_{n=0}^\infty b_n {(x-x_0)}^n\right)
=
\sum_{n=0}^\infty c_n {(x-x_0)}^n ,
\end{equation*}
dengan
$c_n = a_0b_n + a_1 b_{n-1} + \cdots + a_n b_0$.
\end{prop}

Untuk semua $x$ dengan $\sabs{x-x_0} < \rho$, kita mempunyai dua deret konvergen, sehingga
penjumlahan suku demi suku dan perkalian dengan konstanta
mengikuti dari bagian sebelumnya.
Karena untuk $x$ seperti itu deret-deret tersebut konvergen mutlak,
kita dapat menerapkan teorema Mertens untuk menentukan hasil kali kedua deret.
Akibatnya, setelah melakukan operasi aljabar, jari-jari
konvergensi deret yang dihasilkan sekurang-kurangnya $\rho$.
Jari-jari konvergensi hasilnya dapat lebih besar secara ketat daripada
jari-jari konvergensi salah satu deret semula.
Lihat latihan.

Mari kita lihat beberapa contoh deret pangkat.
Polinom hanyalah deret pangkat berhingga:  Suatu polinom
merupakan deret pangkat dengan
$a_n$ bernilai nol untuk semua $n$ yang cukup besar.  Kita mengembangkan
suatu polinom menjadi deret pangkat di sekitar sembarang titik $x_0$ dengan menuliskan
polinom tersebut sebagai polinom dalam $(x-x_0)$.  Sebagai contoh,
$2x^2-3x+4$ sebagai deret pangkat di sekitar $x_0 = 1$ adalah
\begin{equation*}
2x^2-3x+4 = 3 + (x-1) + 2{(x-1)}^2 .
\end{equation*}

Kita juga dapat mengembangkan
\emph{\myindex{fungsi rasional}} (yakni hasil bagi polinom)
menjadi deret pangkat, meskipun kita tidak akan membuktikan fakta ini sepenuhnya di sini.
Perhatikan bahwa deret yang merepresentasikan suatu fungsi rasional hanya mendefinisikan fungsi tersebut
pada suatu interval, meskipun fungsi itu terdefinisi di tempat lain.  Sebagai contoh, untuk
\emph{\myindex{deret geometri}} dan
$x \in (-1,1)$,
\begin{equation*}
\frac{1}{1-x} =
\sum_{n=0}^\infty x^n .
\end{equation*}
deret tersebut divergen jika $\sabs{x} > 1$, meskipun $\frac{1}{1-x}$
terdefinisi untuk semua $x \neq 1$.

Kita dapat menggunakan deret geometri bersama aturan penjumlahan dan
perkalian deret pangkat untuk mengembangkan fungsi rasional menjadi deret
pangkat di sekitar $x_0$,
asalkan penyebutnya tidak nol pada $x_0$.  Kita nyatakan tanpa
bukti bahwa hal ini selalu dapat dilakukan, lalu memberikan contoh
perhitungan semacam itu menggunakan deret geometri.

\begin{example}
Mari kita kembangkan $\frac{x}{1+2x+x^2}$ menjadi deret pangkat di sekitar titik asal ($x_0 = 0$) dan
tentukan jari-jari konvergensinya.

Tuliskan $1+2x+x^2 = {(1+x)}^2 = {\bigl(1-(-x)\bigr)}^2$, dan misalkan
$\sabs{x} < 1$.  Hitung
\begin{equation*}
\begin{split}
\frac{x}{1+2x+x^2}
&=
x \,
{\left(
\frac{1}{1-(-x)}
\right)}^2
\\
&=
x \,
{\left( 
\sum_{n=0}^\infty {(-1)}^n x^n 
\right)}^2
\\
&=
x \,
\left(
\sum_{n=0}^\infty c_n x^n 
\right)
\\
&=
\sum_{n=0}^\infty c_n x^{n+1} .
\end{split}
\end{equation*}
Dengan menggunakan rumus hasil kali deret,
kita memperoleh $c_0 = 1$, $c_1 = -1 -1 = -2$, $c_2 = 1+1+1 = 3$, dan seterusnya.
Dengan demikian, untuk $\sabs{x} < 1$,
\begin{equation*}
\frac{x}{1+2x+x^2}
=
\sum_{n=1}^\infty {(-1)}^{n+1} n x^n .
\end{equation*}
Jari-jari konvergensinya sekurang-kurangnya 1.  Kita serahkan kepada pembaca untuk
memverifikasi bahwa jari-jari konvergensinya tepat sama dengan 1.
\end{example}

Anda dapat menggunakan metode pecahan parsial yang Anda kenal dari kalkulus.
Sebagai contoh, untuk menentukan deret pangkat bagi $\frac{x^3+x}{x^2-1}$ di sekitar 0, tuliskan
\begin{equation*}
\frac{x^3+x}{x^2-1}
=
x + \frac{1}{1+x} - \frac{1}{1-x}
=
x + \sum_{n=0}^\infty {(-1)}^n x^n - \sum_{n=0}^\infty x^n .
\end{equation*}

\subsection{Latihan}

\begin{exercise}
Tentukan kekonvergenan atau kedivergenan deret-deret berikut.

\medskip

\noindent
a)
$\displaystyle \sum_{n=1}^\infty \frac{1}{2^{2n+1}}$
\qquad
b)
$\displaystyle \sum_{n=1}^\infty \frac{{(-1)}^{n}(n-1)}{n}$
\qquad
c)
$\displaystyle \sum_{n=1}^\infty \frac{{(-1)}^n}{n^{1/10}}$
\qquad
d)
$\displaystyle \sum_{n=1}^\infty \frac{n^n}{{(n+1)}^{2n}}$
\end{exercise}

\begin{exercise}
Misalkan $\sum_{n=0}^\infty a_n$ dan $\sum_{n=0}^\infty b_n$ keduanya
konvergen mutlak.
Tunjukkan bahwa deret hasil kali $\sum_{n=0}^\infty c_n$, dengan
$c_n = a_0 b_n + a_1 b_{n-1} + \cdots + a_n b_0$, juga konvergen mutlak.
\end{exercise}

\begin{exercise}[Menantang] \label{exercise:seriesconvergestoanything}
Misalkan $\sum_{n=1}^\infty a_n$ konvergen bersyarat.
Tunjukkan bahwa untuk sebarang $x \in \R$
terdapat suatu penyusunan ulang dari $\sum_{n=1}^\infty a_n$
sedemikian sehingga deret yang disusun ulang tersebut konvergen ke $x$.
Petunjuk: Lihat \exampleref{example:harmonsumanything}.
\end{exercise}

\begin{exercise}
\pagebreak[2]
\leavevmode
\begin{enumerate}[a)]
\item
Tunjukkan bahwa deret harmonik berselang-seling $\sum_{n=1}^\infty \frac{{(-1)}^{n+1}}{n}$
memiliki suatu penyusunan ulang
sedemikian sehingga untuk setiap interval $(x,y)$, terdapat jumlah parsial $s_n$
dari deret yang disusun ulang tersebut sedemikian sehingga $s_n \in (x,y)$.
\item
Tunjukkan bahwa penyusunan ulang yang Anda temukan tidak konvergen.
Lihat \exampleref{example:harmonsumanything}.
\item
Tunjukkan bahwa untuk setiap $x \in \R$, terdapat subbarisan
jumlah-jumlah parsial $\{ s_{n_k} \}_{k=1}^\infty$ dari penyusunan ulang Anda sedemikian sehingga
$\lim\limits_{k\to\infty} s_{n_k} = x$.
\end{enumerate}
\end{exercise}

\begin{exercise}
Untuk setiap deret pangkat berikut, tentukan apakah deret tersebut konvergen;
jika ya, tentukan jari-jari konvergensinya.
Pada bagian f), perhatikan faktorial pada eksponen.

\medskip

\noindent
a)
$\displaystyle \sum_{n=0}^\infty 2^n x^n$
\qquad
b) $\displaystyle \sum_{n=0}^\infty n x^n$
\qquad 
c) 
$\displaystyle \sum_{n=0}^\infty n! \, x^n$
%mbxSTARTIGNORE
\qquad
%mbxENDIGNORE
%mbx <rabr/>
d) $\displaystyle \sum_{n=0}^\infty \frac{1}{(2n)!} {(x-10)}^n$
\qquad
e) $\displaystyle \sum_{n=0}^\infty x^{2n}$
\qquad
f) $\displaystyle \sum_{n=0}^\infty n! \, x^{n!}$
\end{exercise}

\begin{exercise}
Misalkan $\sum_{n=0}^\infty a_n x^n$ konvergen pada $x=1$.
\begin{enumerate}[a)]
\item
Apa yang dapat Anda simpulkan tentang jari-jari konvergensinya?
\item
Jika juga diketahui bahwa pada $x=1$ kekonvergenannya tidak mutlak,
apa yang dapat Anda simpulkan?
\end{enumerate}
\end{exercise}

\begin{exercise}
Kembangkan
$\dfrac{x}{4-x^2}$ menjadi deret pangkat di sekitar $x_0 = 0$,
dan hitung jari-jari konvergensinya.
\end{exercise}

\begin{exercise}
\leavevmode
\begin{enumerate}[a)]
\item
Berikan contoh ketika jari-jari konvergensi $\sum_{n=0}^\infty a_n x^n$ dan
$\sum_{n=0}^\infty b_n x^n$ keduanya 1, tetapi jari-jari konvergensi
jumlah kedua deret tersebut tak hingga.
\item
(Lebih sulit)
Berikan contoh ketika jari-jari konvergensi $\sum_{n=0}^\infty a_n x^n$ dan
$\sum_{n=0}^\infty b_n x^n$ keduanya 1, tetapi jari-jari konvergensi
hasil kali kedua deret tersebut tak hingga.
\end{enumerate}
\end{exercise}

\begin{exercise}
Tentukan cara menghitung jari-jari konvergensi dengan uji rasio.
Artinya, misalkan $\sum_{n=0}^\infty a_n x^n$ adalah deret pangkat dan
$R \coloneqq \lim_{n\to\infty} \frac{\sabs{a_{n+1}}}{\sabs{a_n}}$ ada atau sama dengan $\infty$.
Tentukan jari-jari konvergensinya dan buktikan pernyataan Anda.
\end{exercise}

\begin{samepage}
\begin{exercise}
\leavevmode
\begin{enumerate}[a)]
\item
Buktikan bahwa $\lim_{n\to\infty} n^{1/n} = 1$ dengan prosedur berikut: Tuliskan $n^{1/n} = 1+b_n$ dan
perhatikan bahwa $b_n > 0$ ketika $n > 1$.  Kemudian tunjukkan bahwa ${(1+b_n)}^n \geq
\frac{n(n-1)}{2}b_n^2$ dan gunakan ini untuk menunjukkan bahwa $\lim\limits_{n\to\infty} b_n = 0$.
\item
Gunakan hasil bagian a) untuk menunjukkan bahwa
jika $\sum_{n=0}^\infty a_n x^n$ adalah deret pangkat konvergen dengan jari-jari konvergensi $R$,
maka $\sum_{n=0}^\infty n a_n x^n$ juga konvergen dengan jari-jari konvergensi yang sama.
\end{enumerate}
\end{exercise}
\end{samepage}

\begin{exnote}
Ada berbagai pengertian tentang keterjumlahan (kekonvergenan)
suatu deret
selain pengertian yang telah kita lihat.
Salah satu yang umum adalah \emph{\myindex{keterjumlahan Ces\`aro}}%
\footnote{Dinamai menurut matematikawan Italia
\href{https://en.wikipedia.org/wiki/Ernesto_Ces\%C3\%A0ro}{Ernesto Ces\`aro}
(1859--1906).}.  Misalkan $\sum_{n=1}^\infty a_n$ adalah deret
dan $s_n$ adalah jumlah parsial ke-$n$.  Deret tersebut dikatakan
terjumlah secara Ces\`aro ke $a$ jika
\begin{equation*}
a = \lim_{n\to \infty} \frac{s_1 + s_2 + \cdots + s_n}{n} .
\end{equation*}
\end{exnote}

\begin{exercise}[Menantang]
\pagebreak[2]
\leavevmode
\begin{enumerate}[a)]
\item
Jika $\sum_{n=1}^\infty a_n$ konvergen ke $a$ (dalam pengertian biasa), tunjukkan bahwa
$\sum_{n=1}^\infty a_n$ terjumlah secara Ces\`aro (lihat di atas) ke $a$.
\item
Tunjukkan bahwa dalam pengertian Ces\`aro $\sum_{n=1}^\infty {(-1)}^n$ terjumlah ke
$\nicefrac{-1}{2}$.
\item
Misalkan $a_n \coloneqq k$ ketika $n = k^3$ untuk suatu $k \in \N$,
$a_n \coloneqq -k$ ketika $n = k^3+1$ untuk suatu $k \in \N$,
dan selain itu
misalkan $a_n \coloneqq 0$.  Tunjukkan bahwa $\sum_{n=1}^\infty a_n$ divergen dalam pengertian biasa
(sebenarnya, baik barisan suku-sukunya maupun barisan jumlah parsialnya tidak terbatas), tetapi deret tersebut
terjumlah secara Ces\`aro ke 0 (sekilas tampak agak paradoksal).
\end{enumerate}
\end{exercise}

\begin{exercise}[Menantang]
Tunjukkan bahwa kemonotonan dalam uji deret berselang-seling
diperlukan.  Artinya, berikan barisan bilangan real positif
$\{ x_n \}_{n=1}^\infty$ dengan $\lim_{n\to\infty} x_n = 0$, tetapi
$\sum_{n=1}^\infty {(-1)}^n x_n$ divergen.
\end{exercise}

\begin{exercise}
Berikan suatu deret
$\sum_{n=1}^\infty x_n$ yang konvergen,
tetapi $\sum_{n=1}^\infty x_n^2$ divergen.
Petunjuk: Bandingkan dengan \exerciseref{exercise:squareseriesconv}.
\end{exercise}

\begin{exercise}
Misalkan $\{ c_n \}_{n=1}^\infty$ adalah suatu barisan.  Buktikan bahwa untuk setiap $r \in (0,1)$,
terdapat barisan bilangan asli yang naik secara ketat $\{ n_k \}_{k=1}^\infty$
($n_{k+1} > n_k$) sedemikian sehingga
\begin{equation*}
\sum_{k=1}^\infty c_k x^{n_k}
\end{equation*}
konvergen mutlak untuk semua $x \in [-r,r]$.
\end{exercise}

\begin{exercise}[Tonelli/Fubini untuk penjumlahan, menantang]
\label{exercise:tonellifubiniforsums}%
\index{Fubini untuk penjumlahan}\index{Tonelli untuk penjumlahan}%
Misalkan $\{ x_{k,\ell} \}_{k=1,\ell=1}^\infty$ menyatakan barisan
berindeks ganda dan misalkan $\sigma \colon \N \to \N^2$ adalah suatu bijeksi.
Tinjau deret-deret
\begin{equation*}
\text{i)}~\sum_{i=1}^\infty x_{\sigma(i)},
\qquad
\text{ii)}~\sum_{k=1}^\infty \left( \sum_{\ell=1}^\infty x_{k,\ell} \right),
\qquad
\text{iii)}~\sum_{\ell=1}^\infty \left( \sum_{k=1}^\infty x_{k,\ell} \right) .
\end{equation*}
Ekspresi ii) dan iii) merupakan deret yang suku-sukunya sendiri berupa deret; karena itu, kita katakan ekspresi tersebut konvergen
jika setiap deret di bagian dalam selalu konvergen dan deret di bagian luar kemudian juga
konvergen.
\begin{enumerate}[a)]
\item
(Tonelli) Misalkan $x_{k,\ell} \geq 0$ untuk semua $k,\ell$.
Tunjukkan bahwa ketiga deret i), ii), iii) semuanya divergen (ke $\infty$)
atau semuanya konvergen ke bilangan yang sama.
Dalam kasus divergensi, beberapa deret \myquote{dalam} mungkin bernilai tak hingga
dan dalam hal itu kita menganggap seluruh jumlah tersebut divergen.
\item
(Fubini) Misalkan i) konvergen mutlak.  Tunjukkan bahwa ii) dan iii)
konvergen dan keduanya konvergen ke bilangan yang sama dengan i).
\end{enumerate}
\end{exercise}
